% Descriptive Statistics — Pure Mathematics with Statistics Upper Sixth — Cours signé OnBuch+
% Official MINESEC syllabus. Compiler avec : tectonic PMS08-descriptive-statistics.tex
\newcommand{\DOCMATIERE}{Pure Mathematics with Statistics}
\newcommand{\DOCNIVEAU}{Upper Sixth}
\newcommand{\DOCMODULE}{Upper Sixth — Pure mathematics with statistics}
\newcommand{\DOCLECON}{8}
\newcommand{\DOCTITRE}{Descriptive Statistics}
% OnBuch+ — préambule « cours signés » ENRICHI (compilé avec tectonic / XeLaTeX)
% Système visuel « Pop » d'OnBuch+ : fond crème (jamais blanc), bordures encre
% épaisses, ombres « dures » décalées, titres Archivo Black en capitales.
% Version MATHÉMATIQUES : graphes (pgfplots/tikz), boîtes pédagogiques
% (définition, propriété, méthode, attention, le savais-tu, exercice résolu,
% exercice, TP, bilan), page de garde et signature OnBuch+ ancrée en bas.
%
% Macros attendues dans main.tex AVANT \input{preamble} :
%   \DOCMATIERE  \DOCNIVEAU  \DOCMODULE  \DOCLECON (numéro)  \DOCTITRE
\documentclass[11pt]{article}

\usepackage{fontspec}
\usepackage[english]{babel}
\usepackage[a4paper,margin=2cm,top=3.4cm,bottom=2.8cm,headheight=1.4cm,footskip=1.3cm]{geometry}
\usepackage{amsmath,amssymb,amsfonts}
\usepackage{mathtools} % \xmapsto, \coloneqq... souvent utilisés par les modèles
\usepackage{cancel} % simplifications barrées (bilans de forces)
% \abs{z} (module, valeur absolue) et \norm{u} (norme), utilisés par les modèles
\providecommand{\abs}{}\let\abs\relax\DeclarePairedDelimiter{\abs}{\lvert}{\rvert}
\providecommand{\norm}{}\let\norm\relax\DeclarePairedDelimiter{\norm}{\lVert}{\rVert}
\usepackage[table]{xcolor} % [table] : fournit aussi \rowcolors (alternance de lignes)
\usepackage{pagecolor}
\usepackage{tikz}
\usetikzlibrary{arrows.meta,positioning,calc,shapes.geometric,shapes.misc,shapes.arrows,decorations.pathmorphing,decorations.pathreplacing,decorations.markings,patterns,shadows,mindmap,trees,fit,backgrounds,matrix,intersections,angles,quotes}
\usepackage{pgfplots}
\usepackage[version=4]{mhchem} % équations nucléaires \ce{^{235}_{92}U}
\usepackage[european,siunitx]{circuitikz} % schémas électriques
\pgfplotsset{compat=1.18}
\usepgfplotslibrary{fillbetween,groupplots} % aires entre courbes (calcul intégral)
\usepackage{enumitem}
\usepackage{fancyhdr}
% [most] charge toutes les bibliothèques tcolorbox (listings, minted, raster,
% documentation...) et gonfle inutilement la mémoire TeX sur un document long
% (~300 boîtes) : on ne charge que ce qui est réellement utilisé.
\usepackage[skins,breakable]{tcolorbox}
\usepackage{graphicx}
\usepackage{colortbl}
\usepackage{booktabs}
\usepackage{tabularx}
\usepackage{diagbox} % case d'en-tête diagonale des tableaux à double entrée
% Colonnes extensibles centrée / gauche / droite (C, L, R), souvent utilisées par les modèles
\newcolumntype{C}{>{\centering\arraybackslash}X}
\newcolumntype{L}{>{\raggedright\arraybackslash}X}
\newcolumntype{R}{>{\raggedleft\arraybackslash}X}
\usepackage{array}
\usepackage{multicol}
\usepackage{siunitx}
% parse-numbers=false : accepte aussi \SI{2,5 \times 10^{-3}}{...}
\sisetup{locale=UK,output-decimal-marker={.},group-minimum-digits=4,parse-numbers=false}
\DeclareSIUnit\ohm{\ensuremath{\symup{\Omega}}} % U+2126 absent de la police
\DeclareSIUnit\division{div} % réglages d'oscilloscope (ms/div, V/div)
\DeclareSIUnit\tr{tr} % tours (vitesse de rotation, ex. \SI{1500}{\tr\per\minute})
\usepackage{qrcode}
% \degree : commande fréquemment utilisée par les modèles pour un angle ou une
% température en dehors de \SI{}{} (non fournie par les paquets chargés).
\providecommand{\degree}{^{\circ}}
% \qed (fin de démonstration), utilisé par les modèles sans amsthm
\providecommand{\qed}{\hfill$\square$}
% \barre : double barre des valeurs interdites dans un tableau de variations (tkz-tab)
\providecommand{\barre}{\ensuremath{\|}}
% environnement proof (amsthm non chargé) utilisé par les modèles
\ifdefined\proof\else\newenvironment{proof}[1][Proof]{\par\noindent\textit{#1.}\ }{\qed\par}\fi
\usepackage{lastpage}
\usepackage{hyperref}
\hypersetup{hidelinks}

% ============================================================
% Palette « Pop » OnBuch+ — mêmes hex que la classe P (plus_ui.dart)
% ============================================================
\definecolor{poporange}{HTML}{F59321}
\definecolor{popdark}{HTML}{E07A0C}
\definecolor{popcream}{HTML}{FFF6E8}
\definecolor{popink}{HTML}{1C1714}
\definecolor{poppurple}{HTML}{7A5AE0}
\definecolor{popgreen}{HTML}{2E9E6A}
\definecolor{poppink}{HTML}{E0457B}
\definecolor{popblue}{HTML}{2F7DE1}
\definecolor{popgold}{HTML}{C98A12}
\definecolor{popmuted}{HTML}{8A7F76}
\definecolor{popline}{HTML}{EADFD2}
\definecolor{popgray}{HTML}{8A7F76}
\colorlet{popgrey}{popgray}
% Alias tolérants (le modèle nomme parfois les couleurs différemment)
\colorlet{popwhite}{white}
\colorlet{popblack}{popink}
\colorlet{popred}{poppink}
\colorlet{popyellow}{popgold}
% Teintes claires pour fonds de tableaux / zones
\colorlet{popblueL}{popblue!10!white}
\colorlet{poporangeL}{poporange!14!white}
\colorlet{popgreenL}{popgreen!12!white}
\colorlet{poppurpleL}{poppurple!12!white}
\colorlet{poppinkL}{poppink!12!white}
\colorlet{popgoldL}{popgold!14!white}

\pagecolor{popcream}

% ============================================================
% Typographie
% ============================================================
\defaultfontfeatures{Ligatures=TeX}
\setmainfont{PlusJakartaSans-Medium.ttf}[Path=../../../fonts/,BoldFont=PlusJakartaSans-Bold.ttf]
% Maths en sans-serif (Fira Math) avec chiffres et lettres droites de Plus
% Jakarta, pour que les formules se fondent dans le texte.
\usepackage[math-style=ISO,bold-style=ISO]{unicode-math}
\setmathfont{FiraMath-Regular.otf}[Path=../../../fonts/]
\setmathfont{PlusJakartaSans-Medium.ttf}[Path=../../../fonts/,range={up/num,up/{Latin,latin},\mathup}]
% (icomma retiré : décimales avec point en anglais)
% Caractères mathématiques Unicode absents des polices de texte (Plus Jakarta,
% Archivo Black) : rendus par leur équivalent mathématique, en texte comme en
% formule — sinon ils s'affichent en carré vide (ex. « Arithmétique dans ℤ »).
\usepackage{newunicodechar}
\newunicodechar{ℝ}{\ensuremath{\mathbb{R}}}
\newunicodechar{ℤ}{\ensuremath{\mathbb{Z}}}
\newunicodechar{ℂ}{\ensuremath{\mathbb{C}}}
\newunicodechar{ℕ}{\ensuremath{\mathbb{N}}}
\newunicodechar{ℚ}{\ensuremath{\mathbb{Q}}}
\newunicodechar{𝔻}{\ensuremath{\mathbb{D}}}
\newunicodechar{α}{\ensuremath{\alpha}}
\newunicodechar{β}{\ensuremath{\beta}}
\newunicodechar{γ}{\ensuremath{\gamma}}
\newunicodechar{δ}{\ensuremath{\delta}}
\newunicodechar{ε}{\ensuremath{\varepsilon}}
\newunicodechar{θ}{\ensuremath{\theta}}
\newunicodechar{λ}{\ensuremath{\lambda}}
\newunicodechar{μ}{\ensuremath{\mu}}
\newunicodechar{σ}{\ensuremath{\sigma}}
\newunicodechar{τ}{\ensuremath{\tau}}
\newunicodechar{φ}{\ensuremath{\varphi}}
\newunicodechar{ψ}{\ensuremath{\psi}}
\newunicodechar{ω}{\ensuremath{\omega}}
\newunicodechar{Σ}{\ensuremath{\Sigma}}
% majuscules produites par \MakeUppercase dans les titres (α-aminés -> Α)
\newunicodechar{Α}{\ensuremath{\alpha}}
\newunicodechar{Β}{\ensuremath{\beta}}
\newunicodechar{Γ}{\ensuremath{\Gamma}}
\newunicodechar{Δ}{\ensuremath{\Delta}}
\newunicodechar{Ε}{\ensuremath{\varepsilon}}
\newunicodechar{Θ}{\ensuremath{\Theta}}
\newunicodechar{Λ}{\ensuremath{\Lambda}}
\newunicodechar{Μ}{\ensuremath{\mu}}
\newunicodechar{Τ}{\ensuremath{\tau}}
\newunicodechar{Φ}{\ensuremath{\Phi}}
\newunicodechar{Ψ}{\ensuremath{\Psi}}
\newunicodechar{Ω}{\ensuremath{\Omega}}
\newunicodechar{↦}{\ensuremath{\mapsto}}
\newunicodechar{∈}{\ensuremath{\in}}
\newunicodechar{∉}{\ensuremath{\notin}}
\newunicodechar{∧}{\ensuremath{\wedge}}
\newunicodechar{⇔}{\ensuremath{\Leftrightarrow}}
\newunicodechar{⟺}{\ensuremath{\Longleftrightarrow}}
\newunicodechar{⇒}{\ensuremath{\Rightarrow}}
\newunicodechar{⟹}{\ensuremath{\Longrightarrow}}
\newunicodechar{∘}{\ensuremath{\circ}}
\newunicodechar{∪}{\ensuremath{\cup}}
\newunicodechar{∩}{\ensuremath{\cap}}
\newunicodechar{≡}{\ensuremath{\equiv}}
\newunicodechar{⊂}{\ensuremath{\subset}}
\newunicodechar{⊥}{\ensuremath{\perp}}
\newunicodechar{∃}{\ensuremath{\exists}}
\newunicodechar{∀}{\ensuremath{\forall}}
\newunicodechar{∛}{\ensuremath{\sqrt[3]{\,}}}
\newunicodechar{∜}{\ensuremath{\sqrt[4]{\,}}}
\newunicodechar{⃗}{\ensuremath{{}^{\to}}}
\newunicodechar{ⁿ}{\ifmmode^{n}\else\textsuperscript{\ensuremath{n}}\fi}
\newunicodechar{ˣ}{\ifmmode^{x}\else\textsuperscript{\ensuremath{x}}\fi}
\newunicodechar{ᵇ}{\ifmmode^{b}\else\textsuperscript{\ensuremath{b}}\fi}
\newunicodechar{ʳ}{\ifmmode^{r}\else\textsuperscript{\ensuremath{r}}\fi}
\newunicodechar{ᵘ}{\ifmmode^{u}\else\textsuperscript{\ensuremath{u}}\fi}
\newunicodechar{ʸ}{\ifmmode^{y}\else\textsuperscript{\ensuremath{y}}\fi}
\newunicodechar{ᵃ}{\ifmmode^{a}\else\textsuperscript{\ensuremath{a}}\fi}
\newunicodechar{ᵉ}{\ifmmode^{e}\else\textsuperscript{\ensuremath{e}}\fi}
\newunicodechar{ᵏ}{\ifmmode^{k}\else\textsuperscript{\ensuremath{k}}\fi}
\newunicodechar{ᵖ}{\ifmmode^{p}\else\textsuperscript{\ensuremath{p}}\fi}
\newunicodechar{ᴺ}{\ifmmode^{N}\else\textsuperscript{\ensuremath{N}}\fi}
\newunicodechar{ᶜ}{\ifmmode^{c}\else\textsuperscript{\ensuremath{c}}\fi}
\newunicodechar{ʰ}{\ifmmode^{h}\else\textsuperscript{\ensuremath{h}}\fi}
\newunicodechar{ᵠ}{\ifmmode^{q}\else\textsuperscript{\ensuremath{q}}\fi}
\newunicodechar{ₙ}{\ifmmode_{n}\else\textsubscript{\ensuremath{n}}\fi}
\newunicodechar{ₐ}{\ifmmode_{a}\else\textsubscript{\ensuremath{a}}\fi}
\newunicodechar{ᵢ}{\ifmmode_{i}\else\textsubscript{\ensuremath{i}}\fi}
\newunicodechar{ᵣ}{\ifmmode_{r}\else\textsubscript{\ensuremath{r}}\fi}
\newunicodechar{ₖ}{\ifmmode_{k}\else\textsubscript{\ensuremath{k}}\fi}
\newunicodechar{ᵦ}{\ifmmode_{\beta}\else\textsubscript{\ensuremath{\beta}}\fi}
\newunicodechar{ₓ}{\ifmmode_{x}\else\textsubscript{\ensuremath{x}}\fi}
\newfontfamily\popdisplay{ArchivoBlack-Regular.ttf}[Path=../../../fonts/]
\newfontfamily\poplogo{SpaceGrotesk-Bold.ttf}[Path=../../../fonts/]
\newfontfamily\popsemi{PlusJakartaSans-SemiBold.ttf}[Path=../../../fonts/]
\newfontfamily\popxbold{PlusJakartaSans-ExtraBold.ttf}[Path=../../../fonts/]
\newfontfamily\popxboldls{PlusJakartaSans-ExtraBold.ttf}[Path=../../../fonts/,LetterSpace=3.0]

\setlist[enumerate]{leftmargin=1.7em,itemsep=5pt,topsep=4pt}
\setlist[itemize]{leftmargin=1.7em,itemsep=4pt,topsep=4pt,label={\color{poporange}\textbullet}}
\setlength{\parindent}{0pt}
\setlength{\parskip}{5pt}
\renewcommand{\arraystretch}{1.25}

% pgfplots : style Pop par défaut
\pgfplotsset{
  every axis/.append style={axis line style={popink,line width=1pt},tick style={popink},
    grid style={popline,line width=0.5pt},label style={font=\small},tick label style={font=\footnotesize}},
  popcurve/.style={poporange,line width=1.8pt,no markers,smooth},
}
% Même style disponible dans un \draw TikZ simple (hors pgfplots)
\tikzset{popcurve/.style={poporange,line width=1.8pt}}
\tikzset{popgrid/.style={popline,very thin}}
\colorlet{popcurve}{poporange} % le nom du style est parfois utilisé comme couleur

% ============================================================
% Pill / squiggle
% ============================================================
\newcommand{\poppill}[3]{%
  \tikz[baseline=-0.55ex]\node[draw=popink,line width=1.2pt,rounded corners=2.6pt,%
    fill=#2,inner xsep=6.5pt,inner ysep=3pt]{%
    \popxboldls\fontsize{7}{7}\selectfont\color{#3}\MakeUppercase{#1}};%
}
\newcommand{\popsquiggle}[1]{%
  \tikz[baseline=-0.4ex]{
    \draw[#1,line width=2.6pt,line cap=round]
      plot [smooth] coordinates {(0,0.09) (0.34,0.01) (0.68,0.21) (1.02,0.01) (1.36,0.10)};
  }%
}
% Mot-clé surligné (marqueur orange sous le texte)
\newcommand{\cle}[1]{{\popxbold\color{popdark}#1}}

% ============================================================
% En-tête / pied
% ============================================================
\pagestyle{fancy}
\fancyhf{}
\renewcommand{\headrulewidth}{0pt}
\renewcommand{\footrulewidth}{0pt}
\lhead{\raisebox{-0.15ex}{\poplogo\fontsize{13}{13}\selectfont\color{popink}OnBuch\color{poporange}+}}
\rhead{\poppill{\DOCMATIERE\ \textbullet\ \DOCNIVEAU\ \textbullet\ Chapter \DOCLECON}{white}{popink}}
\lfoot{\popxbold\fontsize{7.6}{7.6}\selectfont\color{popmuted}\MakeUppercase{Signed course OnBuch+}\ \textbullet\ {\popsemi\DOCTITRE}}
\rfoot{\tikz[baseline=-0.6ex]\node[rounded corners=8.5pt,fill=popink,minimum height=17pt,minimum width=17pt,inner xsep=5pt,inner sep=0]{\popxbold\fontsize{8}{8}\selectfont\color{popcream}\thepage\,/\,\pageref*{LastPage}};}

% ============================================================
% Titres
% ============================================================
\newcommand{\chaptitre}[2]{%
  \begin{tcolorbox}[enhanced,colback=poporange,colframe=popink,boxrule=2.6pt,arc=11pt,
      drop shadow=popink,left=15pt,right=15pt,top=12pt,bottom=11pt,before skip=3pt,after skip=18pt]
    \poppill{#2}{popink}{popcream}\par\vspace{9pt}
    {\raggedright\hyphenpenalty=10000\exhyphenpenalty=10000\popdisplay\fontsize{22}{24}\selectfont\color{popcream}\MakeUppercase{#1}\par}
  \end{tcolorbox}
}
% Section numérotée : pastille orange avec le numéro + titre Archivo
\newcounter{popsec}
\newcommand{\coursec}[1]{%
  \stepcounter{popsec}\setcounter{popsub}{0}%
  \par\vspace{16pt}\noindent
  \begin{minipage}{\linewidth}
  \tikz[baseline=-0.7ex]\node[draw=popink,line width=1.6pt,fill=poporange,rounded corners=4pt,minimum size=22pt,inner sep=0]{\popdisplay\fontsize{12}{12}\selectfont\color{popcream}\thepopsec};%
  \hspace{8pt}{\popdisplay\fontsize{15}{17}\selectfont\color{popink}\MakeUppercase{#1}}\par
  \vspace{4pt}\noindent{\color{poporange}\rule{36pt}{3.6pt}}
  \end{minipage}\par\nopagebreak\vspace{8pt}
}
\newcounter{popsub}
\newcommand{\courssub}[1]{%
  \stepcounter{popsub}%
  \par\vspace{9pt}\noindent{\popxbold\fontsize{11.5}{13}\selectfont\color{popink}{\color{poporange}\thepopsec.\thepopsub}\hspace{6pt}#1}\par\nopagebreak\vspace{4pt}
}

% ============================================================
% Boîtes pédagogiques — même recette : fond blanc, bordure épaisse colorée,
% ombre dure encre, badge plein. Titre optionnel [#1].
% ============================================================
\tcbset{popbase/.style={breakable,enhanced,colback=white,boxrule=2.3pt,arc=9pt,
  shadow={3pt}{-3pt}{0pt}{popink},left=14pt,right=14pt,top=11pt,bottom=11pt,before skip=11pt,after skip=15pt,
  fontupper=\popsemi\fontsize{10.6}{14.5}\selectfont\color{popink},
  fontlower=\popsemi\fontsize{10.6}{14.5}\selectfont\color{popink}}}
% Coche dessinée (le glyphe ✓ n'existe pas dans les polices du document)
\newcommand{\popcheck}[1]{\tikz[baseline=-0.5ex]\draw[#1,line width=1.6pt,line cap=round,line join=round](-0.13,0.0)--(-0.04,-0.09)--(0.13,0.1);}
\providecommand{\coche}{\popcheck{popgreen}}
\providecommand{\resultat}[1]{\par\smallskip\noindent\fcolorbox{popgreen}{popgreenL}{\parbox{\dimexpr\linewidth-2\fboxsep-2\fboxrule\relax}{\textbf{Result.} #1}}\par\smallskip}
\newcommand{\popboxhead}[3]{\poppill{#1}{#2}{white}\ifx\relax#3\relax\else\hspace{6pt}{\popxbold\fontsize{10.6}{12}\selectfont\color{popink}#3}\fi\par\vspace{7pt}}

\newtcolorbox{aretenir}[1][]{popbase,colframe=popgreen,before upper={\popboxhead{Key points}{popgreen}{#1}}}
\newtcolorbox{exemplebox}[1][]{popbase,colframe=poporange,before upper={\popboxhead{Exemple}{poporange}{#1}}}
\newtcolorbox{definition}[1][]{popbase,colframe=popblue,before upper={\popboxhead{Definition}{popblue}{#1}}}
\newtcolorbox{propriete}[1][]{popbase,colframe=poppurple,before upper={\popboxhead{Property}{poppurple}{#1}}}
\newtcolorbox{methode}[1][]{popbase,colframe=popgold,before upper={\popboxhead{Method}{popgold}{#1}}}
\newtcolorbox{attention}[1][]{popbase,colframe=poppink,before upper={\popboxhead{Warning}{poppink}{#1}}}
\newtcolorbox{experience}[1][]{popbase,colframe=popblue,colback=popblueL,before upper={\popboxhead{Experiment}{popblue}{#1}}}
% « Le savais-tu ? » : carte violette pleine (contexte, culture, Cameroun)
\newtcolorbox{savaistu}[1][]{popbase,colback=poppurple,colframe=popink,
  fontupper=\popsemi\fontsize{10.4}{14.2}\selectfont\color{white},
  before upper={\poppill{Did you know?}{white}{poppurple}\ifx\relax#1\relax\else\hspace{6pt}{\popxbold\color{white}#1}\fi\par\vspace{7pt}}}
% Exercice résolu : énoncé en haut, \tcblower puis la solution sur fond crème
\newcounter{popexo}\newcounter{popexr}
\newtcolorbox{exoresolu}[1][]{popbase,colframe=poporange,colbacklower=popcream,
  segmentation style={popink,line width=1.2pt,dashed},
  before upper={\stepcounter{popexr}\popboxhead{Worked example \thepopexr}{poporange}{#1}},
  before lower={\poppill{Solution}{popink}{popcream}\par\vspace{6pt}}}
% Exercice (non résolu en place) — la correction est dans la partie Corrigés
\newtcolorbox{exercice}[1][]{popbase,colframe=popink,
  before upper={\stepcounter{popexo}\popboxhead{Exercise \thepopexo}{popink}{#1}}}
% Corrigé
\newtcolorbox{corrige}[1][]{popbase,colframe=popgreen,colback=popgreenL,
  before upper={\popboxhead{Solution}{popgreen}{#1}}}
% Objectifs / prérequis / bilan
\newtcolorbox{objectifs}{popbase,colframe=popink,colback=poporangeL,
  before upper={\popboxhead{Chapter objectives}{popink}{}}}
\newtcolorbox{prerequis}{popbase,colframe=popink,colback=popblueL,
  before upper={\popboxhead{Prerequisites}{popblue}{}}}
\newtcolorbox{bilan}{popbase,colframe=popink,colback=white,boxrule=2.8pt,
  before upper={\popboxhead{Fiche bilan}{poporange}{L'essentiel à connaître pour l'examen}}}
% Figure encadrée avec légende
\newtcolorbox{popfigure}[1][]{enhanced,breakable=false,colback=white,colframe=popline,boxrule=1.4pt,arc=8pt,
  left=10pt,right=10pt,top=10pt,bottom=8pt,before skip=10pt,after skip=12pt,halign=center,
  lower separated=false,halign lower=center,
  fontlower=\popsemi\fontsize{9}{11.5}\selectfont\color{popmuted},
  #1}
\newcommand{\legende}[1]{\par\vspace{6pt}{\popsemi\fontsize{9}{11.5}\selectfont\color{popmuted}#1}}

% ============================================================
% Signature OnBuch+
% ============================================================
\newcommand{\poplogoinline}[1]{{\poplogo\fontsize{#1}{#1}\selectfont\color{popink}OnBuch\color{poporange}+}}
\newcommand{\signatureonbuch}{%
  \begin{tcolorbox}[enhanced,colback=popink,colframe=popink,boxrule=2.2pt,arc=10pt,
      drop shadow=poporange,left=14pt,right=14pt,top=10pt,bottom=10pt,before skip=0pt,after skip=0pt]
    \tikz[baseline=-0.7ex]\node[circle,fill=popgreen,draw=popcream,line width=1.3pt,minimum size=22pt,inner sep=0]{\popcheck{white}};%
    \hspace{9pt}\begin{minipage}[c]{0.62\linewidth}
      {\popxbold\fontsize{12}{13}\selectfont\color{popcream}Signed\ \textbullet\ {\poplogo OnBuch\color{poporange}+}}\par\vspace{2pt}
      {\raggedright\popsemi\fontsize{8}{10.5}\selectfont\color{popline}\MakeUppercase{OnBuch+ teaching team}\\\MakeUppercase{Follows the official MINESEC syllabus}\par}
    \end{minipage}\hfill
    {\poplogo\fontsize{22}{22}\selectfont\color{popcream}OnBuch\color{poporange}+}
  \end{tcolorbox}%
}

% ============================================================
% Page de garde
% #1 titre · #2 matière · #3 niveau · #4 module · #5 n° leçon · #6 durée ·
% #7 description / objectifs (texte ou itemize)
% ============================================================
\newcommand{\pagegarde}[7]{%
  \thispagestyle{empty}
  \newgeometry{a4paper,margin=1.7cm,top=1.6cm,bottom=1.4cm}%
  \begin{tikzpicture}[remember picture,overlay]
    \fill[poporange!18] ([xshift=-3.2cm,yshift=-3.6cm]current page.north east) circle (4.6cm);
    \fill[poppurple!14] ([xshift=2.4cm,yshift=7.6cm]current page.south west) circle (3.8cm);
  \end{tikzpicture}%
  {\poplogo\fontsize{30}{30}\selectfont\color{popink}OnBuch\color{poporange}+}\hfill
  \raisebox{0.6ex}{\poppill{Signed course \textbullet\ Premium}{popink}{popcream}}\par
  \vspace{1pt}\popsquiggle{poppurple}\par
  \vspace{8pt}
  \begin{tcolorbox}[enhanced,colback=poporange,colframe=popink,boxrule=2.8pt,arc=13pt,
      drop shadow=popink,left=18pt,right=18pt,top=12pt,bottom=13pt,after skip=10pt]
    \poppill{#2 \textbullet\ #3}{popink}{popcream}\hspace{5pt}\poppill{#4}{white}{popink}\par\vspace{8pt}
    {\popdisplay\fontsize{14}{14}\selectfont\color{popink}CHAPTER #5}\par\vspace{4pt}
    {\raggedright\hyphenpenalty=10000\exhyphenpenalty=10000\popdisplay\fontsize{25}{27}\selectfont\color{popcream}\MakeUppercase{#1}\par}
    \vspace{8pt}
    {\popxbold\fontsize{9.5}{9.5}\selectfont\color{popink}\MakeUppercase{Suggested time: #6}}
  \end{tcolorbox}
  \begin{tcolorbox}[enhanced,colback=white,colframe=popink,boxrule=2.2pt,arc=10pt,
      drop shadow=popink,left=16pt,right=16pt,top=10pt,bottom=10pt,after skip=8pt]
    \poppill{In this chapter}{poppurple}{white}\par\vspace{5pt}
    \popsemi\fontsize{9.3}{12.6}\selectfont\color{popink}#7
  \end{tcolorbox}
  \noindent\begin{minipage}[t]{0.487\linewidth}
    \begin{tcolorbox}[enhanced,colback=popgreen,colframe=popink,boxrule=2.2pt,arc=10pt,
        drop shadow=popink,left=11pt,right=11pt,top=10pt,bottom=10pt,height=3.1cm,valign=center]
      \tcbox[colback=white,colframe=popink,boxrule=1.3pt,arc=5pt,boxsep=0pt,nobeforeafter,
        left=3pt,right=3pt,top=3pt,bottom=3pt,valign=center]{\qrcode[height=1.9cm]{https://play.google.com/store/apps/details?id=com.luvvix.onbuch&hl=en}}%
      \hspace{8pt}\begin{minipage}[c]{0.5\linewidth}
        {\popdisplay\fontsize{11}{12.5}\selectfont\color{white}\MakeUppercase{Download OnBuch}}\par\vspace{4pt}
        {\raggedright\popsemi\fontsize{8.4}{11}\selectfont\color{white}Free \textbullet\ Google Play\\AI tutor, past papers, results\par}
      \end{minipage}
    \end{tcolorbox}
  \end{minipage}\hfill
  \begin{minipage}[t]{0.487\linewidth}
    \begin{tcolorbox}[enhanced,colback=poppurple,colframe=popink,boxrule=2.2pt,arc=10pt,
        drop shadow=popink,left=11pt,right=11pt,top=10pt,bottom=10pt,height=3.1cm,valign=center]
      \tikz[baseline=-0.7ex]\node[circle,fill=white,draw=popink,line width=1.3pt,minimum size=1.6cm,inner sep=0]{\popdisplay\fontsize{24}{24}\selectfont\color{poppurple}+};%
      \hspace{8pt}\begin{minipage}[c]{0.6\linewidth}
        {\popdisplay\fontsize{11}{12.5}\selectfont\color{white}\MakeUppercase{Go OnBuch+}}\par\vspace{4pt}
        {\raggedright\popsemi\fontsize{8.4}{11}\selectfont\color{white}All signed courses, unlimited AI tutor, PDF sheets — OnBuch+ tab in the app\par}
      \end{minipage}
    \end{tcolorbox}
  \end{minipage}
  \vspace{8pt}
  \popcontenu
  \vfill
  \signatureonbuch
  \restoregeometry
  \newpage
}

% Rangée « Ce cours contient » (page de garde)
\newcommand{\poptile}[3]{% #1 couleur #2 gros texte #3 légende
  \begin{tcolorbox}[enhanced,colback=#1,colframe=popink,boxrule=2pt,arc=9pt,shadow={2.5pt}{-2.5pt}{0pt}{popink},
      left=6pt,right=6pt,top=9pt,bottom=9pt,halign=center,valign=center,height=1.75cm,width=\linewidth,nobeforeafter]
    {\popdisplay\fontsize{12.5}{13}\selectfont\color{white}\MakeUppercase{#2}}\par\vspace{3pt}
    {\centering\popsemi\fontsize{7.8}{9.5}\selectfont\color{white}#3\par}
  \end{tcolorbox}}
\newcommand{\popcontenu}{%
  \par\noindent{\popxboldls\fontsize{7.5}{7.5}\selectfont\color{popmuted}THIS SIGNED COURSE CONTAINS}\par\vspace{7pt}
  \noindent\begin{minipage}[t]{0.235\linewidth}\poptile{popblue}{Course}{complete, illustrated, step by step}\end{minipage}\hfill
  \begin{minipage}[t]{0.235\linewidth}\poptile{poporange}{Methods}{and worked examples}\end{minipage}\hfill
  \begin{minipage}[t]{0.235\linewidth}\poptile{poppink}{Exercises}{exam style + solutions}\end{minipage}\hfill
  \begin{minipage}[t]{0.235\linewidth}\poptile{popgreen}{Summary sheet}{the essentials to remember}\end{minipage}\par}

% Sommaire
\newtcolorbox{sommaire}{enhanced,colback=white,colframe=popink,boxrule=2.2pt,arc=10pt,
  drop shadow=popink,left=18pt,right=18pt,top=13pt,bottom=13pt,before skip=6pt,after skip=20pt,
  before upper={\poppill{Contents}{poppurple}{white}\par\vspace{9pt}\popsemi\fontsize{10.6}{14.5}\selectfont\color{popink}}}

% Signature de fin de cours — poussée tout en bas de la dernière page
\newcommand{\findecours}{%
  \par\vfill\nopagebreak
  \begin{center}\popsquiggle{poporange}\end{center}\vspace{4pt}
  \signatureonbuch
}
\AtBeginDocument{\def\llbracket{\mathopen{[\mkern-3.5mu[}}\def\rrbracket{\mathclose{]\mkern-3.5mu]}}} % intervalles d'entiers (glyphe ⟦ absent de la police)
% Glyphes absents de Fira Math : redéfinis à partir de glyphes disponibles
\AtBeginDocument{%
  \def\setminus{\mathbin{\backslash}}\let\smallsetminus\setminus
  \def\vdots{\mathord{\vbox{\baselineskip3.2pt\lineskiplimit0pt\kern4pt\hbox{.}\hbox{.}\hbox{.}}}}%
  \def\ddots{\mathinner{\mkern1mu\raise6.5pt\hbox{.}\mkern2mu\raise3.6pt\hbox{.}\mkern2mu\raise0.7pt\hbox{.}\mkern1mu}}%
  \def\triangle{\mathord{\scalebox{0.9}{$\symup{\Delta}$}}}%
  \def\nabla{\mathord{\raisebox{\depth}{\scalebox{1}[-1]{$\symup{\Delta}$}}}}%
  \def\checkmark{\popcheck{popdark}}%
  \def\star{\mathbin{\ast}}\def\bigstar{\mathord{\ast}}%
  \def\diamond{\mathbin{\scalebox{0.6}{$\Diamond$}}}%
  \def\bigcup{\mathop{\mathchoice{\raisebox{-0.3em}{\scalebox{1.7}{$\cup$}}}{\scalebox{1.25}{$\cup$}}{\cup}{\cup}}\displaylimits}%
  \def\bigcap{\mathop{\mathchoice{\raisebox{-0.3em}{\scalebox{1.7}{$\cap$}}}{\scalebox{1.25}{$\cap$}}{\cap}{\cap}}\displaylimits}%
  \def\lesseqqgtr{\mathrel{\lessgtr}}\def\bigoplus{\mathop{\oplus}\displaylimits}%
}
\providecommand{\ssi}{if and only if} % abréviation fréquente
\AtBeginDocument{\providecommand{\wideparen}{\overparen}} % arc de cercle

% Fonctions usuelles que les modèles utilisent sans que amsmath les fournisse
\providecommand{\arsinh}{\operatorname{arsinh}}\providecommand{\arcosh}{\operatorname{arcosh}}
\providecommand{\artanh}{\operatorname{artanh}}\providecommand{\arsech}{\operatorname{arsech}}
\providecommand{\arcsch}{\operatorname{arcsch}}\providecommand{\arcoth}{\operatorname{arcoth}}
\providecommand{\arcsinh}{\operatorname{arsinh}}\providecommand{\arccosh}{\operatorname{arcosh}}
\providecommand{\arctanh}{\operatorname{artanh}}\providecommand{\sech}{\operatorname{sech}}
\providecommand{\csch}{\operatorname{csch}}\providecommand{\arcsec}{\operatorname{arcsec}}
\providecommand{\arccsc}{\operatorname{arccsc}}\providecommand{\arccot}{\operatorname{arccot}}
\providecommand{\sgn}{\operatorname{sgn}}\providecommand{\lcm}{\operatorname{lcm}}
\providecommand{\Var}{\operatorname{Var}}\providecommand{\Cov}{\operatorname{Cov}}

%% FIGFIX-BEGIN v5 — correctifs globaux des figures (docs/onbuch-generation/tools/figfix)
\usepackage{adjustbox}\usepackage{etoolbox}\tcbuselibrary{raster}
% toute figure TikZ/pgfplots plus large que la ligne est réduite à la largeur disponible
\BeforeBeginEnvironment{tikzpicture}{\begin{adjustbox}{max width=\linewidth}}
\AfterEndEnvironment{tikzpicture}{\end{adjustbox}}
% légendes pgfplots lisibles par-dessus les courbes
\pgfplotsset{every axis/.append style={legend style={fill=white,fill opacity=0.92,text opacity=1,draw=black!15,inner sep=3pt}}}
% étiquettes d'arêtes (syntaxe edge["..."]) avec fond blanc
\tikzset{every edge quotes/.append style={fill=white,inner sep=1.2pt}}
%% FIGFIX-END
\begin{document}

\pagegarde{Descriptive Statistics}{Pure Mathematics with Statistics}{Upper Sixth}{Upper Sixth — Pure mathematics with statistics}{8}{23 periods}{This chapter equips Upper Sixth students with the complete machinery of descriptive statistics: data collection and presentation, measures of central tendency and dispersion, index numbers and skewness. It prepares learners to handle any GCE Advanced Level statistics question with rigour and confidence.
\vspace{4pt}
\begin{multicols}{2}\small
\begin{itemize}
\item By the end of this chapter I can define and correctly use the basic vocabulary of statistics (population, sample, variable, parameter, statistic).
\item By the end of this chapter I can classify data and choose an appropriate method of collection.
\item By the end of this chapter I can construct frequency tables, including grouped frequency distributions with correct class boundaries.
\item By the end of this chapter I can represent data graphically using bar charts, pie charts, line graphs, histograms, frequency polygons and ogives.
\item By the end of this chapter I can compute and interpret the mode, median, mean, range, quartiles, deciles and percentiles of grouped and ungrouped data.
\item By the end of this chapter I can compute and interpret the variance and standard deviation of grouped and ungrouped data, and combine means of several groups.
\end{itemize}\end{multicols}}

\begin{sommaire}
\begin{multicols}{2}
\begin{enumerate}
\item Basic Terminologies and Types of Data
\item Tabular Presentation and Grouped Frequency Distributions
\item Graphical Representation I: Bar, Pie and Line Charts
\item Graphical Representation II: Histogram, Frequency Polygon and Ogive
\item Measures of Central Tendency I: Mode, Median, Range
\item Measures of Central Tendency II: Mean and Combined Mean
\item Measures of Dispersion: Variance and Standard Deviation
\item Quartiles, Deciles, Percentiles and Interquartile Range
\item Index Numbers
\item Skewness and Shape of Distributions
\item Integration activity
\item Methods and worked examples
\item Exercises
\item Solutions to the exercises
\item Summary sheet
\end{enumerate}
\end{multicols}
\end{sommaire}

\begin{objectifs}
\begin{itemize}
\item By the end of this chapter I can define and correctly use the basic vocabulary of statistics (population, sample, variable, parameter, statistic).
\item By the end of this chapter I can classify data and choose an appropriate method of collection.
\item By the end of this chapter I can construct frequency tables, including grouped frequency distributions with correct class boundaries.
\item By the end of this chapter I can represent data graphically using bar charts, pie charts, line graphs, histograms, frequency polygons and ogives.
\item By the end of this chapter I can compute and interpret the mode, median, mean, range, quartiles, deciles and percentiles of grouped and ungrouped data.
\item By the end of this chapter I can compute and interpret the variance and standard deviation of grouped and ungrouped data, and combine means of several groups.
\item By the end of this chapter I can construct and interpret simple, weighted, Laspeyres and Paasche index numbers.
\item By the end of this chapter I can identify the shape of a distribution and measure skewness using Pearson's and the quartile coefficients.
\end{itemize}
\end{objectifs}

\begin{prerequis}
\begin{itemize}
\item Arithmetic operations with fractions, decimals and percentages (Lower Sixth).
\item Reading and plotting points in the Cartesian plane (Lower Sixth).
\item Basic algebraic manipulation, including the sigma (Σ) notation for sums (Lower Sixth, Sequences and Series).
\item Simple proportional reasoning and ratios (Form 5 / Lower Sixth).
\item Square roots and use of a scientific calculator.
\end{itemize}
\end{prerequis}

\newpage

\coursec{Basic Terminologies and Types of Data}

Every morning, the Douala City Council statistician must decide where to send the waste-collection trucks: which neighbourhood produces the most refuse, and how much does it vary from day to day? At the same time, the Ministry of Public Health compares the average weight of babies born in Yaoundé and Bamenda, while the National Institute of Statistics publishes the consumer price index that determines whether your pocket money still buys the same quantity of puff-puff. All these decisions rest on one toolkit: \cle{descriptive statistics}. In this chapter, you will learn to collect, organise, summarise and interpret data exactly as Cameroonian institutions do — and exactly as the GCE Advanced Level examiners expect.

But before we can summarise data, we must answer more basic questions: \emph{What exactly is data? Where does it come from? How do we collect it without being misled?} A wrong answer at this first stage poisons everything that follows. This section therefore builds the vocabulary and the critical judgement on which the whole chapter stands.

\courssub{Statistics as a Science}

\textbf{Intuition.} Suppose you record the heights of all 40 students in your Upper Sixth class and compute the average. You have \emph{described} your class. Now suppose you measure only 40 students chosen at random from a whole school of 1200 students, and you use their average to draw a conclusion about \emph{all} the students. You have gone beyond description: you have made an \emph{inference}. This distinction is the first great divide in statistics.

\begin{definition}[Statistics, Descriptive and Inferential]
\cle{Statistics} is the science of collecting, organising, presenting, analysing and interpreting numerical data in order to make decisions under uncertainty.
\begin{itemize}
\item \cle{Descriptive statistics} deals with the collection, organisation, presentation and summary of data (tables, charts, averages, measures of spread) \emph{without} drawing conclusions beyond the data at hand.
\item \cle{Inferential statistics} uses information from a \emph{sample} to draw conclusions, estimates or predictions about a larger \emph{population}, together with a measure of the reliability of those conclusions.
\end{itemize}
\end{definition}

This chapter is devoted to descriptive statistics; inference (estimation, hypothesis testing) comes later in your course.

\courssub{Population, Census and Sample}

\textbf{Intuition.} A cook tasting a pot of \emph{ndolé} does not eat the whole pot to know whether there is enough salt: she stirs well and tastes one spoonful. The pot is the \emph{population}; the spoonful is the \emph{sample}. The tasting works only if the spoonful fairly represents the pot — that is why she stirs first.

\begin{definition}[Population, Census, Sample]
\begin{itemize}
\item A \cle{population} is the complete set of all individuals or items under study. It need not be people: the population may be all bags of cement produced by a factory in Douala in June, or all candidates sitting the GCE Advanced Level in a given year.
\item A \cle{census} is a study that collects information from \emph{every} member of the population.
\item A \cle{sample} is a subset of the population, chosen to be studied in place of the whole population.
\item The \cle{sampling frame} is the actual list (or device) from which the sample is drawn, e.g.\ the school register, the electoral roll, a list of households.
\item A \cle{sampling unit} is a single member of the sampling frame that can be selected.
\end{itemize}
\end{definition}

\begin{exemplebox}[Population or Sample?]
A market inspector in Bamenda wants to know the average mass of 50\,kg bags of rice sold at the main market. There are 800 bags. He weighs 30 bags chosen at random.
\begin{itemize}
\item Population: all 800 bags of rice at the market.
\item Sampling frame: the list (or physical arrangement) of the 800 bags.
\item Sample: the 30 bags actually weighed.
\item Sampling unit: one bag of rice.
\end{itemize}
\end{exemplebox}

\begin{popfigure}
\begin{tikzpicture}[font=\small]
\draw[fill=popblueL, draw=popblue, thick] (0,0) ellipse (3.4 and 1.9);
\node[popdark] at (0,1.35) {\textbf{Population}};
\node[popmuted] at (0,-1.35) {all items under study};
\draw[fill=popgreenL, draw=popgreen, thick] (0.3,0.1) ellipse (2.1 and 1.1);
\node[popdark] at (0.3,0.75) {\textbf{Sampling frame}};
\draw[fill=popgold!40, draw=poporange, thick] (0.6,-0.05) ellipse (0.95 and 0.5);
\node[popdark] at (0.6,-0.05) {\textbf{Sample}};
\draw[-{latex}, thick, popink] (-1.2,-1.15) -- (-1.6,-0.6);
\draw[-{latex}, thick, popink] (1.7,-0.75) -- (1.35,-0.35);
\end{tikzpicture}
\legende{The sample is drawn from the sampling frame, which should cover the whole population.}
\end{popfigure}

\begin{attention}[A frame that does not match the population]
If the sampling frame omits part of the population, the sample cannot represent it. Surveying only students who own smartphones about school fees excludes poorer families and biases the results. Always check: \emph{does my frame cover the whole population?}
\end{attention}

\textbf{Why sample at all?} A census gives exact information but is expensive, slow, and sometimes destructive (testing the lifetime of a bulb destroys it; checking every bag of cement for strength breaks every bag). A well-chosen sample gives reliable information quickly and cheaply. The simplest way to choose fairly is \cle{simple random sampling}: every member of the population has the same chance of being selected, for example by numbering the sampling frame and drawing numbers from a hat or from a table of random numbers.

\courssub{Statistical Variables and Their Types}

When we collect data, we record some \emph{characteristic} of each individual: height, religion, number of siblings, favourite subject. Such a characteristic is called a statistical variable.

\begin{definition}[Statistical Variable]
A \cle{statistical variable} (or simply \emph{variable}) is a characteristic or attribute that can take different values from one member of the population to another. The set of all recorded values is the \cle{raw data}.
\end{definition}

\begin{definition}[Qualitative and Quantitative Variables]
\begin{itemize}
\item A variable is \cle{qualitative} (categorical) if its values are \emph{qualities or categories}, not numbers on which arithmetic makes sense.
\begin{itemize}
\item \cle{Nominal}: categories with no natural order (e.g.\ region of origin, blood group, religion).
\item \cle{Ordinal}: categories with a natural order but no measurable spacing (e.g.\ grade A, B, C; level of satisfaction: poor / fair / good).
\end{itemize}
\item A variable is \cle{quantitative} (numerical) if its values are numbers on which arithmetic is meaningful.
\begin{itemize}
\item \cle{Discrete}: takes only isolated values, usually obtained by \emph{counting} (e.g.\ number of children per family: $0,1,2,\dots$).
\item \cle{Continuous}: can take any value in an interval, usually obtained by \emph{measuring} (e.g.\ height, mass, time, temperature).
\end{itemize}
\end{itemize}
\end{definition}

\begin{popfigure}
\begin{tikzpicture}[font=\small,
 grow=right, level distance=3.6cm,
 level 1/.style={sibling distance=3.2cm},
 level 2/.style={sibling distance=1.6cm},
 every node/.style={draw=popline, rounded corners=2pt, fill=white, align=center, inner sep=3pt}]
\node[fill=popblueL, draw=popblue, thick, align=center]{\textbf{Statistical}\\ \textbf{variable}}
 child { node[fill=popgreenL, draw=popgreen, thick, align=center]{\textbf{Quantitative}\\ (numerical)}
   child { node[align=center]{Continuous\\ e.g.\ height} }
   child { node[align=center]{Discrete\\ e.g.\ goals scored} } }
 child { node[fill=popgold!40, draw=poporange, thick, align=center]{\textbf{Qualitative}\\ (categorical)}
   child { node[align=center]{Ordinal\\ e.g.\ exam grade} }
   child { node[align=center]{Nominal\\ e.g.\ blood group} } };
\end{tikzpicture}
\legende{Classification of statistical variables.}
\end{popfigure}

\begin{exoresolu}[Classifying Variables]
Classify each variable as qualitative (nominal or ordinal) or quantitative (discrete or continuous):
\begin{enumerate}
\item the number of goals scored by each team in the MTN Elite One championship;
\item the time taken by students to run 100\,m;
\item the marital status of teachers in a school;
\item the rank (1st, 2nd, 3rd, \dots) of candidates in an examination.
\end{enumerate}
\tcblower
\begin{enumerate}
\item Goals are \emph{counted} and take whole-number values: quantitative \textbf{discrete}.
\item Time is \emph{measured} and can take any value in an interval (e.g.\ $12.37$\,s): quantitative \textbf{continuous}.
\item Marital status (single, married, widowed, divorced) is a set of unordered categories: qualitative \textbf{nominal}.
\item Ranks have a natural order but the ``distance'' between 1st and 2nd is not numerically defined: qualitative \textbf{ordinal}.
\end{enumerate}
\end{exoresolu}

\begin{attention}[Numbers are not always quantitative]
A jersey number, a telephone number or a candidate's index number is made of digits, yet adding two jersey numbers is meaningless: these are \textbf{nominal} labels, not quantitative variables. The test is always: \emph{does arithmetic on the values make sense?}
\end{attention}

\courssub{Parameter versus Statistic; Raw versus Processed Data}

\begin{definition}[Parameter and Statistic]
A \cle{parameter} is a numerical summary (mean, proportion, standard deviation, \dots) computed from the \emph{whole population}. A \cle{statistic} is the corresponding summary computed from a \emph{sample}. We usually use the sample statistic to \emph{estimate} the unknown population parameter.
\end{definition}

\begin{attention}[Common mistake: confusing parameter and statistic]
The mean mass of \emph{all} 800 bags of rice is a \textbf{parameter}; the mean mass of the 30 weighed bags is a \textbf{statistic}. In the GCE exam, read the question twice and ask: \emph{was this number computed from the whole population or from a sample?} Writing ``the population mean is the sample mean'' is a classic error: the sample mean only \emph{estimates} the population mean.
\end{attention}

\cle{Raw data} are the values exactly as collected, before any organisation (e.g.\ a list of 40 heights in centimetres, in the order the students were measured). \cle{Processed data} are raw data that have been organised or summarised: sorted, tallied, grouped into frequency tables, presented in charts, or reduced to averages. The rest of this chapter is precisely the art of turning raw data into processed data without losing the essential information.

\courssub{Methods of Data Collection}

\begin{definition}[Primary and Secondary Data]
\cle{Primary data} are collected \emph{first-hand} by the investigator for the current study. \cle{Secondary data} already exist, having been collected earlier by someone else for another purpose (published reports, hospital records, NIS publications, past examination results).
\end{definition}

The main methods of collecting primary data are:

\begin{itemize}
\item \textbf{Direct observation}: the investigator watches and records events (e.g.\ counting vehicles passing the Wouri bridge between 7\,a.m.\ and 8\,a.m.). \emph{Advantage}: accurate, no dependence on what people say. \emph{Limitation}: slow, costly, and unsuitable for opinions or past events.
\item \textbf{Interview}: a trained enumerator asks questions face-to-face or by telephone. \emph{Advantage}: high response rate, questions can be explained, suitable for respondents who cannot read. \emph{Limitation}: expensive, time-consuming, and the interviewer's tone may influence answers.
\item \textbf{Questionnaire}: a written list of questions distributed to respondents who fill it in themselves. \emph{Advantage}: cheap, fast, reaches many people, anonymous (encourages honesty). \emph{Limitation}: low response rate, no chance to clarify misunderstood questions.
\item \textbf{Experiment}: the investigator controls conditions and records the outcome (e.g.\ testing two fertilisers on identical plots of maize in Buea). \emph{Advantage}: can establish cause and effect. \emph{Limitation}: may be costly, impractical or unethical.
\item \textbf{Use of existing records} (secondary data): consulting registers, files, publications. \emph{Advantage}: immediate and cheap. \emph{Limitation}: definitions may differ, data may be outdated or collected with a different aim.
\end{itemize}

\begin{methode}[Choosing an Appropriate Data-Collection Method]
\begin{enumerate}
\item \textbf{Define the variable} precisely and the population concerned.
\item \textbf{Check whether secondary data already exist} that answer the question reliably; if yes, use them.
\item \textbf{Match the method to the variable}: opinions and attitudes $\to$ questionnaire or interview; observable behaviour $\to$ observation; cause-and-effect questions $\to$ experiment.
\item \textbf{Weigh cost, time and accuracy}: a census is ideal but often impossible; choose a sample method you can afford that still represents the population.
\item \textbf{Pilot the instrument}: test your questionnaire on a few respondents first and correct ambiguous or leading questions before the full survey.
\end{enumerate}
\end{methode}

\courssub{Sources of Error and Bias}

Even a carefully planned survey can go wrong. The main dangers are:

\begin{itemize}
\item \cle{Non-response bias}: people who refuse to answer (or cannot be reached) often differ systematically from those who answer. If only 40\,\% of questionnaires are returned, the missing 60\,\% may hold a different opinion.
\item \cle{Leading questions}: wording that pushes the respondent towards an answer, e.g.\ ``Do you agree that the excellent new school fees policy is fair?''
\item \textbf{Ambiguous or double questions}: ``Do you like maths and physics?'' — a student may like only one.
\item \textbf{Sensitive questions}: respondents may lie about income, age or health unless anonymity is guaranteed.
\item \textbf{Faulty sampling frame}: as seen above, a frame that omits part of the population.
\item \textbf{Measurement and recording errors}: a badly calibrated scale, a mistyped figure.
\end{itemize}

\begin{exemplebox}[A biased versus a well-designed questionnaire]
\emph{Biased version} (circulated in a school): ``Don't you think the boring canteen food should be improved? Tick: Yes / No.'' Problems: the question is \textbf{leading} (``boring'' pushes towards ``Yes''), and it assumes the respondent eats at the canteen.

\emph{Improved version}: ``How often do you eat at the school canteen? Tick one: Daily / Sometimes / Never. If you eat there, how would you rate the food? Tick one: Very good / Good / Fair / Poor.'' Each question is neutral, single, and offers ordered response categories.
\end{exemplebox}

\begin{aretenir}[Key Points of this Section]
\begin{itemize}
\item \textbf{Descriptive} statistics summarises data; \textbf{inferential} statistics generalises from a sample to a population.
\item \textbf{Population} = all items under study; \textbf{census} = study of all of them; \textbf{sample} = the subset actually studied, drawn from a \textbf{sampling frame}.
\item Variables are \textbf{qualitative} (nominal or ordinal) or \textbf{quantitative} (discrete or continuous); digits used as labels are still qualitative.
\item A \textbf{parameter} describes the population; a \textbf{statistic} describes a sample and estimates the parameter.
\item \textbf{Primary} data are collected first-hand (observation, interview, questionnaire, experiment); \textbf{secondary} data come from existing records.
\item Beware of \textbf{non-response} and \textbf{leading questions}: a biased instrument makes all later calculations worthless.
\end{itemize}
\end{aretenir}

\begin{exercice}[Terminology Check]
A GCE Board official wishes to study the performance of all Advanced Level candidates in the Centre Region. She selects 200 candidates at random from the official results database and computes their average mark in Pure Mathematics, obtaining $52.4$.
\begin{enumerate}
\item Identify the population, the sampling frame, the sample and the sampling unit.
\item Is $52.4$ a parameter or a statistic? Justify.
\item Is the average mark of \emph{all} candidates in the region known? What is it called?
\item State one advantage and one limitation of using the results database (secondary data) instead of interviewing candidates.
\end{enumerate}
\end{exercice}

\begin{corrige}[Exercise 1]
\begin{enumerate}
\item Population: all Advanced Level candidates in the Centre Region. Sampling frame: the official results database. Sample: the 200 selected candidates. Sampling unit: one candidate.
\item $52.4$ is a \textbf{statistic}, because it was computed from a sample of 200 candidates, not from the whole population.
\item No; it is unknown. It is the population mean, a \textbf{parameter}, which the statistic $52.4$ estimates.
\item Advantage: the data already exist, so the study is fast and cheap. Limitation: the official was not in control of how the data were produced and cannot obtain information not recorded in the database (e.g.\ candidates' study habits).
\end{enumerate}
\end{corrige}

\begin{exercice}[Variables and Methods]
For each study below, state the variable, its type, and the most appropriate method of data collection:
\begin{enumerate}
\item a farmer in Ngaoundéré wants to compare the yields of two varieties of maize;
\item a school principal wants to know how students travel to school (on foot, by bike, by taxi, by bus);
\item a nurse records the birth weights of babies delivered in a district hospital during one month.
\end{enumerate}
\end{exercice}

\begin{corrige}[Exercise 2]
\begin{enumerate}
\item Variable: yield of maize (e.g.\ in kg per plot); quantitative \textbf{continuous}. Method: \textbf{experiment} — grow both varieties on identical plots under the same conditions and measure the yields.
\item Variable: means of transport to school; qualitative \textbf{nominal}. Method: \textbf{questionnaire} (or direct observation at the school gate); the categories have no natural order.
\item Variable: birth weight; quantitative \textbf{continuous}. Method: \textbf{direct measurement/observation} recorded in the hospital register (the register itself is a source of secondary data for later studies).
\end{enumerate}
\end{corrige}

\begin{savaistu}[Statistics in Cameroon]
Cameroon's National Institute of Statistics (NIS) coordinates official statistics, conducts the national population and housing census, and publishes the consumer price index that tracks the cost of living. A general census is one of the largest operations in the country: tens of thousands of enumerators visit households across all ten regions, because a census — unlike a survey — aims to reach \emph{every} member of the population. The results guide the drawing of electoral districts, the siting of schools and hospitals, and the allocation of state resources to the regions.
\end{savaistu}

\coursec{Tabular Presentation and Grouped Frequency Distributions}

In the previous section we learnt to recognise the different \cle{types of data} (qualitative, quantitative, discrete, continuous) and the \cle{methods of collection}. But raw data, as collected, are usually a long, disordered list of numbers: they tell us almost nothing at first sight. Imagine the marks of $40$ Upper Sixth students written on a single line. To answer simple questions such as \emph{``How many students scored between $50$ and $60$?''} or \emph{``What is the most frequent mark?''}, we must first \textbf{organise} the data. The first and most fundamental tool of organisation is the \cle{frequency table}. This section teaches you how to tally raw data, how to build simple and grouped frequency tables, and how to prepare the cumulative columns that will be used later for the ogive.

\courssub{Tallying and the Simple Frequency Table}

\begin{definition}[Frequency, relative frequency, percentage frequency]
Let a statistical variable take the values $x_1, x_2, \dots, x_k$ in a population of size $N$.
\begin{itemize}
\item The \cle{frequency} (or absolute frequency) $f_i$ of the value $x_i$ is the number of times the value $x_i$ occurs in the data.
\item The \cle{relative frequency} of $x_i$ is the ratio $r_i = \dfrac{f_i}{N}$.
\item The \cle{percentage frequency} of $x_i$ is $r_i \times 100\ \% = \dfrac{f_i}{N}\times 100\ \%$.
\end{itemize}
\end{definition}

\begin{propriete}[Sum of frequencies]
The frequencies always satisfy
\[ \sum_{i=1}^{k} f_i = N \qquad \text{and} \qquad \sum_{i=1}^{k} r_i = \sum_{i=1}^{k}\frac{f_i}{N} = \frac{1}{N}\sum_{i=1}^{k} f_i = 1. \]
Thus \textbf{the sum of the relative frequencies equals $1$} (and the sum of the percentage frequencies equals $100\ \%$). This is an excellent check of any table you construct in the examination.
\end{propriete}

The practical technique for counting is \cle{tallying}: for each data value read, we draw one stroke in the row of that value. Strokes are grouped in bundles of five, the fifth stroke crossing the first four (\text{||||} crossed), so that counting at the end is fast and reliable.

\begin{exemplebox}[Tallying the marks of 40 students]
The marks (out of $20$) obtained by $40$ Upper Sixth students of GBHS Bamenda in a Mathematics test are:
\begin{center}
$12, 8, 15, 9, 12, 7, 14, 12, 10, 8, 16, 12, 9, 11, 13, 8, 12, 10, 15, 7,$\\
$11, 12, 9, 14, 8, 12, 10, 13, 16, 9, 12, 11, 8, 14, 12, 10, 13, 9, 12, 11$
\end{center}
Reading the data one by one and tallying gives:
\begin{center}
\begin{tabularx}{\linewidth}{|c|l|c|X|c|}
\hline
\rowcolor{popblueL} Mark $x_i$ & Tally & $f_i$ & Relative frequency $f_i/N$ & \% freq. \\
\hline
$7$ & || & $2$ & $0.050$ & $5.0$ \\
\hline
$8$ & \text{||||} crossed, | & $6$ & $0.150$ & $15.0$ \\
\hline
$9$ & \text{||||} crossed & $5$ & $0.125$ & $12.5$ \\
\hline
$10$ & |||| & $4$ & $0.100$ & $10.0$ \\
\hline
$11$ & |||| & $4$ & $0.100$ & $10.0$ \\
\hline
$12$ & \text{||||} crossed, \text{||||} crossed & $10$ & $0.250$ & $25.0$ \\
\hline
$13$ & ||| & $3$ & $0.075$ & $7.5$ \\
\hline
$14$ & ||| & $3$ & $0.075$ & $7.5$ \\
\hline
$15$ & || & $2$ & $0.050$ & $5.0$ \\
\hline
$16$ & | & $1$ & $0.025$ & $2.5$ \\
\hline
\rowcolor{popblueL} Total & & $40$ & $1.000$ & $100.0$ \\
\hline
\end{tabularx}
\end{center}
Check: $2+6+5+4+4+10+3+3+2+1 = 40 = N$ and the relative frequencies sum to $1$. Immediately we see that the mark $12$ is the most frequent (it will later be called the \emph{mode}).
\end{exemplebox}

\courssub{Why and When to Group Data}

The table above works well because the marks take only $10$ distinct values. But suppose the variable is the \emph{height} of $200$ students, measured in centimetres: almost every value would be different, and a simple frequency table would have nearly $200$ rows, each with frequency $1$. The table would be useless. The remedy is to \textbf{group} the values into intervals called \cle{classes}.

\begin{aretenir}[When to group data]
Group data into classes when:
\begin{itemize}
\item the variable is \textbf{continuous} (heights, masses, times, temperatures), or
\item the variable is discrete but takes \textbf{many distinct values} (marks out of $100$, ages of a large population, incomes in FCFA).
\end{itemize}
Grouping loses some detail (individual values disappear) but reveals the \textbf{shape} of the distribution.
\end{aretenir}

\textbf{How many classes?} Too few classes hide the pattern; too many classes recreate the disorder. A common guide is \cle{Sturges' rule}:
\[ k \approx 1 + 3.322 \log_{10} N, \]
where $N$ is the number of observations and $k$ the suggested number of classes. For $N = 40$: $k \approx 1 + 3.322\log_{10}40 \approx 1 + 3.322(1.602) \approx 6.3$, so about $6$ classes. Sturges' rule is only a \emph{guide}: in practice we choose a convenient number (often between $5$ and $15$) giving a \textbf{convenient class width} (a whole number, a multiple of $5$ or $10$, etc.).

\begin{savaistu}[Sturges' rule and Cameroon's census]
Sturges' rule was published in 1926 by Herbert Sturges. It works well for moderately sized samples ($N < 200$). For very large data sets (like Cameroon's population census conducted by the National Institute of Statistics), other rules such as the \emph{Rice rule} ($k = 2\sqrt[3]{N}$) or the \emph{Freedman–Diaconis rule} are often preferred. In GCE examinations, Sturges' rule is the expected reference unless the question specifies otherwise.
\end{savaistu}

\courssub{Class Limits, Boundaries, Width and Class Mark}

\begin{definition}[Vocabulary of a class]
A class is an interval of values, written for instance $50$--$59$ or $[50; 60[$.
\begin{itemize}
\item The \cle{class limits} are the stated end values: lower class limit $50$, upper class limit $59$.
\item The \cle{class boundaries} are the true separation points between consecutive classes when the variable is continuous or rounded: lower boundary $49.5$, upper boundary $59.5$.
\item The \cle{class width} is the difference between the upper and lower \emph{boundaries}: $59.5 - 49.5 = 10$.
\item The \cle{class mark} (or mid-point) is the centre of the class:
\[ \text{class mark} = \frac{\text{lower limit} + \text{upper limit}}{2} = \frac{50+59}{2} = 54.5. \]
\end{itemize}
\end{definition}

Why boundaries? Marks recorded ``to the nearest unit'' mean that a recorded $59$ stands for any true value in $[58.5; 59.5[$. The boundary $59.5$ is therefore the genuine frontier with the next class $60$--$69$, whose lower boundary is also $59.5$. Boundaries \textbf{remove the gaps} between classes and guarantee that every possible value falls in exactly one class.

\begin{popfigure}
\begin{tikzpicture}[scale=0.9]
\draw[->, popline, thick] (0,0) -- (13.5,0);
\foreach \x/\lab in {2/49.5, 5/59.5, 8/69.5, 11/79.5}
  \draw[popdark, thick] (\x,0.12) -- (\x,-0.12) node[below, font=\small] {\lab};
\draw[popblue, very thick] (2,0.35) -- (5,0.35);
\draw[poporange, very thick] (5,0.35) -- (8,0.35);
\draw[popgreen, very thick] (8,0.35) -- (11,0.35);
\node[popblue, font=\small, align=center] at (3.5,0.75) {class $50$--$59$};
\node[poporange, font=\small, align=center] at (6.5,0.75) {class $60$--$69$};
\node[popgreen, font=\small, align=center] at (9.5,0.75) {class $70$--$79$};
\fill[poppurple] (3.5,0.35) circle (2.2pt) node[below=6pt, font=\small, poppurple] {mark $54.5$};
\fill[poppurple] (6.5,0.35) circle (2.2pt) node[below=6pt, font=\small, poppurple] {mark $64.5$};
\fill[poppurple] (9.5,0.35) circle (2.2pt) node[below=6pt, font=\small, poppurple] {mark $74.5$};
\node[font=\small, popmuted, align=center] at (6.5,-1.3) {No gap: boundary $59.5$ belongs to both class edges};
\end{tikzpicture}
\legende{Class limits (stated) versus class boundaries (true frontiers) for a continuous variable, with class marks at the centres.}
\end{popfigure}

\begin{attention}[Common mistake: overlapping classes or gaps]
Two classic errors destroy a grouped table:
\begin{itemize}
\item \textbf{Overlapping classes}, e.g.\ $50$--$60$, $60$--$70$: where does a value of $60$ go? Use $50$--$59$, $60$--$69$ for discrete data, or the convention $[50; 60[$, $[60; 70[$ (lower limit included, upper excluded) for continuous data.
\item \textbf{Gaps}, e.g.\ $50$--$59$, $61$--$70$: the value $60$ has no class. Consecutive classes must touch at a boundary.
\end{itemize}
Also, never use classes of different widths in the same table unless the syllabus question explicitly does so: unequal widths distort the histogram.
\end{attention}

\courssub{Constructing a Grouped Frequency Distribution Step by Step}

\begin{methode}[Constructing a grouped frequency distribution]
\begin{enumerate}
\item Find the \textbf{range}: $R = x_{\max} - x_{\min}$.
\item Choose the \textbf{number of classes} $k$ (Sturges' rule as a guide, then round sensibly).
\item Compute the \textbf{class width}: $w \approx \dfrac{R}{k}$, rounded \emph{up} to a convenient value.
\item Fix the \textbf{lower limit of the first class} (at or just below $x_{\min}$), then write all classes.
\item Write the \textbf{class boundaries} and the \textbf{class marks}.
\item \textbf{Tally} each observation into its class and record the \textbf{frequencies} $f_i$.
\item Check that $\sum f_i = N$, then add relative and percentage frequency columns if required.
\end{enumerate}
\end{methode}

\begin{exoresolu}[Grouping the marks of 40 students]
The marks (out of $100$) of $40$ candidates in a mock GCE paper range from $23$ to $88$. Construct a grouped frequency distribution with boundaries, class marks, relative frequencies and a ``less than'' cumulative column.
\tcblower
\textbf{Solution.} \textbf{Step 1:} range $R = 88 - 23 = 65$. \textbf{Step 2:} Sturges gives $k \approx 6.3$; we take $k = 7$. \textbf{Step 3:} $w \approx 65/7 \approx 9.3$; we round up to the convenient width $w = 10$. \textbf{Step 4:} first class $20$--$29$ (starts just below $23$). \textbf{Steps 5--7:} tallying the $40$ marks gives:

\begin{center}
\begin{tabularx}{\linewidth}{|c|c|c|c|c|X|c|}
\hline
\rowcolor{popblueL} Class & Boundaries & Mark $m_i$ & Tally & $f_i$ & Rel.\ freq.\ $f_i/40$ & Cum.\ $<$ \\
\hline
$20$--$29$ & $19.5$--$29.5$ & $24.5$ & ||| & $3$ & $0.075$ & $3$ \\
\hline
$30$--$39$ & $29.5$--$39.5$ & $34.5$ & |||| & $4$ & $0.100$ & $7$ \\
\hline
$40$--$49$ & $39.5$--$49.5$ & $44.5$ & \text{||||} crossed, | & $6$ & $0.150$ & $13$ \\
\hline
$50$--$59$ & $49.5$--$59.5$ & $54.5$ & \text{||||} crossed, ||| & $8$ & $0.200$ & $21$ \\
\hline
$60$--$69$ & $59.5$--$69.5$ & $64.5$ & \text{||||} crossed, || & $7$ & $0.175$ & $28$ \\
\hline
$70$--$79$ & $69.5$--$79.5$ & $74.5$ & \text{||||} crossed, | & $6$ & $0.150$ & $34$ \\
\hline
$80$--$89$ & $79.5$--$89.5$ & $84.5$ & \text{||||} crossed, | & $6$ & $0.150$ & $40$ \\
\hline
\rowcolor{popblueL} Total & & & & $40$ & $1.000$ & \\
\hline
\end{tabularx}
\end{center}

Checks: $\sum f_i = 3+4+6+8+7+6+6 = 40$; relative frequencies sum to $1$. Each class mark is the mean of the limits, e.g.\ $\frac{50+59}{2} = 54.5$; each width is $59.5-49.5 = 10$. The last cumulative entry equals $N = 40$, as it must.
\end{exoresolu}

\courssub{Discrete versus Continuous Grouping}

\begin{aretenir}[Two styles of grouping]
\begin{itemize}
\item \textbf{Discrete data} (counts: number of children per family, marks): classes are written with whole-number limits, $0$--$4$, $5$--$9$, $10$--$14$, and boundaries fall at half-integers: $-0.5$--$4.5$, $4.5$--$9.5$, \dots
\item \textbf{Continuous data} (measurements: height, mass, time): classes are written with the convention $[a; b[$, e.g.\ $[160; 170[$ cm, $[170; 180[$ cm, where the lower bound is included and the upper bound excluded; limits and boundaries then coincide.
\end{itemize}
In both cases the class mark is the centre of the class and the width is the difference of the boundaries.
\end{aretenir}

\courssub{Cumulative Frequency Columns: Preparing the Ogive}

\begin{definition}[Cumulative frequencies]
\begin{itemize}
\item The \cle{less than cumulative frequency} of a class is the total number of observations falling \textbf{below the upper boundary} of that class. It is obtained by adding frequencies from the top of the table downwards.
\item The \cle{more than cumulative frequency} of a class is the total number of observations falling \textbf{at or above the lower boundary} of that class. It is obtained by adding frequencies from the bottom upwards.
\end{itemize}
\end{definition}

For the table of the worked example, the ``more than'' column reads: $40, 37, 33, 27, 19, 12, 6$ (starting from $40$ and subtracting $f_i$ successively). Note the duality: for every class, $(\text{less than}) + (\text{more than}) = N$. These columns, plotted against the \emph{upper} (respectively \emph{lower}) boundaries, give the \textbf{ogive} studied in Section~4, from which the median and quartiles will be read graphically.

\begin{exercice}[Practice]
The masses (in kg, to the nearest kg) of $30$ sacks of cocoa at a buying centre in Kumba range from $41$ kg to $80$ kg.
\begin{enumerate}
\item Using Sturges' rule, suggest a number of classes and a class width.
\item Construct a grouped frequency table with classes $40$--$44$, $45$--$49$, \dots, giving boundaries, class marks and frequencies (invent plausible frequencies summing to $30$).
\item Add the ``less than'' and ``more than'' cumulative columns and verify that their sum in each row equals $30$.
\end{enumerate}
\end{exercice}

\begin{corrige}[Exercise 1]
\textbf{1.} $k \approx 1 + 3.322\log_{10}30 \approx 1 + 3.322(1.477) \approx 5.9$, so $k = 6$ classes; range $= 80-41 = 39$, width $\approx 39/6 = 6.5$; the given width $5$ with $9$ classes ($40$--$44$ to $80$--$84$) is also acceptable and more convenient.
\textbf{2.--3.} Example of a valid table (frequencies may vary, only the checks matter): $f = 2, 3, 5, 6, 5, 4, 3, 1, 1$ for the classes $40$--$44$ to $80$--$84$; $\sum f = 30$. Boundaries: $39.5$--$44.5$, $44.5$--$49.5$, \dots; marks: $42, 47, 52, \dots$ Less-than column: $2, 5, 10, 16, 21, 25, 28, 29, 30$; more-than column: $30, 28, 25, 20, 14, 9, 5, 2, 1$; each pair sums to $30$. \checkmark
\end{corrige}

\begin{aretenir}[Key points of the section]
\begin{itemize}
\item Tallying in bundles of five gives reliable frequencies; always check $\sum f_i = N$ and $\sum r_i = 1$.
\item Group when the variable is continuous or takes many values; Sturges' rule $k \approx 1 + 3.322\log_{10}N$ guides the number of classes.
\item Limits are stated; boundaries are the true frontiers (half a unit beyond the limits for data rounded to the nearest unit); width $=$ difference of boundaries; class mark $=$ centre.
\item Classes must neither overlap nor leave gaps.
\item ``Less than'' and ``more than'' cumulative columns prepare the ogive and satisfy: their sum in each row equals $N$.
\end{itemize}
\end{aretenir}

\coursec{Graphical Representation I: Bar, Pie and Line Charts}

In the previous section we learnt how to organise raw data into frequency tables and grouped frequency distributions. A table is precise, but it is not always \emph{speaking}: a parent reading a school report, a trader in Mokolo market, or a minister presenting a budget rarely has time to scan columns of figures. A well-chosen \cle{diagram} conveys the same information at a glance. In this section we study the three diagrams used for \emph{qualitative} or \emph{discrete} data: the \cle{bar chart}, the \cle{pie chart} and the \cle{line graph}. Diagrams for continuous grouped data (histogram, frequency polygon, ogive) will follow in the next section.

\courssub{Bar Charts}

\textbf{Intuition.}\par
Suppose a teacher asks the 40 students of an Upper Sixth class their favourite subject. The natural picture is one vertical bar per subject, the height of the bar showing how many students chose it. Longer bar, more students: the comparison is immediate.

\begin{definition}[Bar chart]
A \cle{bar chart} is a diagram in which each category of a qualitative (or discrete) variable is represented by a rectangular bar of \textbf{equal width}, separated by \textbf{gaps}, whose height (or length) is proportional to the frequency of that category.
\end{definition}

The two golden rules are: \textbf{equal widths} and \textbf{gaps between bars}. The gaps are not decorative: they signal that the categories are separate and that there is no continuous scale between them.

\begin{exemplebox}[Simple bar chart: favourite subjects]
The favourite subjects of 40 Upper Sixth students are: Mathematics 12, Physics 8, Biology 9, Economics 6, Literature 5. The simple bar chart below displays these frequencies directly.
\end{exemplebox}

\begin{popfigure}
\begin{tikzpicture}
\begin{axis}[
    ybar, bar width=0.7cm,
    width=11cm, height=6.5cm,
    ymin=0, ymax=14,
    ylabel={Number of students},
    symbolic x coords={Maths, Physics, Biology, Economics, Literature},
    xtick=data,
    x tick label style={font=\small},
    ytick={0,2,4,6,8,10,12,14},
    axis y line*=left, axis x line*=bottom,
    nodes near coords,
    every node near coord/.append style={font=\small},
]
\addplot[fill=popblue, draw=popdark] coordinates {(Maths,12) (Physics,8) (Biology,9) (Economics,6) (Literature,5)};
\end{axis}
\end{tikzpicture}
\legende{Simple bar chart: favourite subjects of 40 Upper Sixth students.}
\end{popfigure}

\textbf{Multiple (side-by-side) bar charts.}\par
When we want to compare two or more groups across the same categories, we place the bars \emph{side by side}, one colour per group, with a legend. For instance, comparing boys and girls in the same survey of favourite subjects: for each subject we draw two adjacent bars.

\begin{exemplebox}[Multiple bar chart: favourite subjects by gender]
The table below shows the favourite subjects of 20 boys and 20 girls.
\begin{center}
\begin{tabularx}{\linewidth}{|l|X|X|X|}
\hline
\rowcolor{popblueL} Subject & Boys & Girls & Total \\
\hline
Mathematics & 8 & 4 & 12 \\
Physics & 5 & 3 & 8 \\
Biology & 3 & 6 & 9 \\
Economics & 2 & 4 & 6 \\
Literature & 2 & 3 & 5 \\
\hline
\end{tabularx}
\end{center}
The multiple bar chart allows immediate visual comparison: boys prefer Mathematics and Physics, while girls prefer Biology and Economics.
\end{exemplebox}

\begin{popfigure}
\begin{tikzpicture}
\begin{axis}[
    ybar, bar width=0.5cm,
    width=11cm, height=6.5cm,
    ymin=0, ymax=10,
    ylabel={Number of students},
    symbolic x coords={Maths, Physics, Biology, Economics, Literature},
    xtick=data,
    x tick label style={font=\small},
    ytick={0,2,4,6,8,10},
    axis y line*=left, axis x line*=bottom,
    legend style={at={(0.5,1.02)}, anchor=south, legend columns=2, font=\small},
    nodes near coords,
    every node near coord/.append style={font=\tiny},
]
\addplot[fill=popblue, draw=popdark] coordinates {(Maths,8) (Physics,5) (Biology,3) (Economics,2) (Literature,2)};
\addplot[fill=poppink, draw=popdark] coordinates {(Maths,4) (Physics,3) (Biology,6) (Economics,4) (Literature,3)};
\legend{Boys, Girls}
\end{axis}
\end{tikzpicture}
\legende{Multiple bar chart: favourite subjects by gender.}
\end{popfigure}

\textbf{Component (stacked) bar charts.}\par
When we want to show both the \emph{total} of each category and its \emph{composition}, we stack the components on top of one another inside a single bar per category. For example, one bar per class (Upper Sixth Arts, Upper Sixth Science), each bar split into ``pass'' and ``fail'' segments. The total height gives the class size; the segments give the breakdown.

\begin{exemplebox}[Reading a stacked bar chart]
A stacked bar for Upper Sixth Science has total height 50, with a ``pass'' segment of height 38. We read off immediately: 50 candidates, 38 passes, hence $50-38=12$ failures. Extracting frequencies from a diagram is a routine GCE task: always read heights against the axis, and use totals to deduce missing segments.
\end{exemplebox}

\begin{attention}[Percentage component bar chart]
A variant is the \cle{percentage component bar chart} where each bar has the same height (representing 100\%) and the segments show the percentage composition. This is useful when totals differ widely but we want to compare proportions.
\end{attention}

\courssub{Pie Charts}

\textbf{Intuition.}\par
A pie chart represents the \emph{whole} (total frequency) by a full circle and each category by a slice (sector) whose angle is proportional to its frequency. It answers the question: ``what \emph{share} of the whole does each category represent?''

\begin{propriete}[Angle of a sector]
If a category has frequency $f$ out of a total frequency $N$, the angle of its sector is
\[ \theta = \frac{f}{N}\times 360^{\circ}. \]
This follows because the full circle ($360^{\circ}$) represents the whole frequency $N$, and angles are shared proportionally. The sum of all sector angles is exactly $360^{\circ}$ — a built-in check.
\end{propriete}

\begin{methode}[Computing pie-chart angles]
\begin{enumerate}
\item Compute the total frequency $N$.
\item For each category, compute $\theta=\dfrac{f}{N}\times 360^{\circ}$.
\item Round each angle (usually to the nearest degree), then \textbf{check} that the rounded angles sum to $360^{\circ}$; if not, adjust the largest sector by $1^{\circ}$.
\item Draw a circle, mark a starting radius, and construct the sectors one after another with a protractor.
\item Label each sector (or use a legend) with the category and, if useful, its percentage.
\end{enumerate}
\end{methode}

\begin{exoresolu}[Household expenditure of a Cameroonian family]
Statement. A family in Yaound\'e spends its monthly budget of $200\,000$ FCFA as follows: Food $80\,000$, Rent $50\,000$, Transport $20\,000$, School fees $30\,000$, Savings $20\,000$. Construct the angles of a pie chart representing this budget.
\tcblower
Solution. Total $N=200\,000$ FCFA. Applying $\theta=\dfrac{f}{N}\times 360^{\circ}$:
\begin{center}
\begin{tabularx}{\linewidth}{|l|X|X|}
\hline
\rowcolor{popblueL} Category & Amount (FCFA) & Angle \\
\hline
Food & $80\,000$ & $\frac{80000}{200000}\times360^{\circ}=144^{\circ}$ \\
\hline
Rent & $50\,000$ & $90^{\circ}$ \\
\hline
School fees & $30\,000$ & $54^{\circ}$ \\
\hline
Transport & $20\,000$ & $36^{\circ}$ \\
\hline
Savings & $20\,000$ & $36^{\circ}$ \\
\hline
\textbf{Total} & $200\,000$ & $360^{\circ}$ $\checkmark$ \\
\hline
\end{tabularx}
\end{center}
The angles sum to $360^{\circ}$, so the chart can be drawn as below.
\end{exoresolu}

\begin{popfigure}
\begin{tikzpicture}[scale=1.1]
% Food 144, Rent 90, Fees 54, Transport 36, Savings 36
\fill[popblue]   (0,0) -- (90:2.4) arc (90:234:2.4) -- cycle;
\fill[poporange] (0,0) -- (234:2.4) arc (234:324:2.4) -- cycle;
\fill[popgreen]  (0,0) -- (324:2.4) arc (324:378:2.4) -- cycle;
\fill[popgold]   (0,0) -- (378:2.4) arc (378:414:2.4) -- cycle;
\fill[poppurple] (0,0) -- (414:2.4) arc (414:450:2.4) -- cycle;
\draw[popdark] (0,0) circle (2.4);
\node[font=\small, text=white] at (162:1.5) {Food 40\%};
\node[font=\small, text=white] at (279:1.5) {Rent 25\%};
\node[font=\small, text=white] at (351:1.6) {Fees 15\%};
\node[font=\small, text=white] at (33:1.7) {Tr. 10\%};
\node[font=\small, text=white] at (69:1.7) {Sav. 10\%};
\end{tikzpicture}
\legende{Pie chart of a monthly family budget of $200\,000$ FCFA.}
\end{popfigure}

\begin{attention}[Angles and percentages]
Do not confuse the \emph{angle} with the \emph{percentage}. The percentage is $\frac{f}{N}\times100$; the angle is $\frac{f}{N}\times360^{\circ}$. Also, a pie chart is meaningless if the categories do not form a complete whole (e.g. favourite subjects where students could pick two): the slices would not add up to one pie.
\end{attention}

\courssub{Line Graphs and Time Series}

\begin{definition}[Time series and line graph]
A \cle{time series} is a set of observations of a variable recorded at successive points in time (days, months, years). Its \cle{line graph} is obtained by plotting the values against time and joining consecutive points by straight line segments. Time is always on the horizontal axis.
\end{definition}

The line graph is the natural tool for spotting \emph{trends} (is the quantity rising or falling overall?), \emph{seasonal patterns} (regular peaks, e.g. rainfall in the rainy season in Buea) and \emph{turning points}.

\begin{exemplebox}[Cocoa prices over 12 months]
The average farm-gate price of cocoa (in FCFA per kg) recorded monthly in the South-West Region is plotted below. The rising trend from January to May, the dip around July, and the recovery towards December are visible at a glance — information that a bare table would hide.
\end{exemplebox}

\begin{popfigure}
\begin{tikzpicture}
\begin{axis}[
    width=11.5cm, height=6.5cm,
    xlabel={Month}, ylabel={Price (FCFA/kg)},
    xmin=1, xmax=12, ymin=700, ymax=1100,
    xtick={1,2,3,4,5,6,7,8,9,10,11,12},
    xticklabels={Jan,Feb,Mar,Apr,May,Jun,Jul,Aug,Sep,Oct,Nov,Dec},
    x tick label style={font=\small},
    ytick={700,800,900,1000,1100},
    grid=both, grid style={popline, very thin},
    axis lines=left,
]
\addplot[color=poporange, very thick, mark=*, mark size=1.8pt] coordinates {
(1,780) (2,810) (3,860) (4,920) (5,980) (6,950) (7,880) (8,900) (9,940) (10,1000) (11,1040) (12,1060)};
\end{axis}
\end{tikzpicture}
\legende{Line graph of average cocoa prices over 12 months.}
\end{popfigure}

\begin{exemplebox}[Reading a line graph: interpolation]
From the cocoa price graph, estimate the price in mid-April. April (month 4) is 920, May (month 5) is 980. Mid-April corresponds to $x=4.5$. Linear interpolation gives $\frac{920+980}{2}=950$ FCFA/kg. Always state that this is an \emph{estimate} assuming linear change between recorded points.
\end{exemplebox}

\courssub{Choosing the Right Diagram}

\begin{methode}[Choosing an appropriate diagram]
\begin{enumerate}
\item Identify the \textbf{type of data}: qualitative categories, discrete counts, or a variable evolving in time.
\item To \textbf{compare frequencies} between categories: use a bar chart (simple, multiple or stacked).
\item To show \textbf{shares of a total}: use a pie chart.
\item To show \textbf{evolution over time}: use a line graph.
\item For \textbf{continuous grouped data}: use a histogram or frequency polygon (next section) — never a gapped bar chart.
\end{enumerate}
\end{methode}

\begin{attention}[Bar chart or histogram?]
A very common GCE mistake is to draw bars \emph{touching each other} for categorical data, or conversely to leave gaps when the data are continuous. Rule: \textbf{gaps} for separate categories (bar chart); \textbf{touching bars} for continuous data grouped into classes (histogram). If your classes are $10$--$19$, $20$--$29$, \dots, the bars must touch (using class boundaries $9.5$--$19.5$, \dots), and the diagram is a histogram, not a bar chart.
\end{attention}

\courssub{Reading, Interpreting and Detecting Misleading Graphs}

Reading a chart means extracting exact frequencies: read bar heights against the vertical axis, sector angles with the formula $f=\frac{\theta}{360^{\circ}}\times N$, and points of a line graph against both axes. Interpreting means going further: stating the mode category, comparing shares, describing a trend.

\begin{attention}[Misleading graphs — GCE exam awareness]
Examiners love to show a ``bad'' graph and ask why it misleads. Watch for:
\begin{itemize}
\item \textbf{Truncated vertical axis}: if the frequency axis starts at, say, 90 instead of 0, a bar of height 100 looks twice as tall as a bar of height 95, although the real difference is only 5\%. Always check where the axis starts.
\item \textbf{Unequal bar widths}: a bar drawn twice as wide \emph{and} slightly taller suggests (by its area) a difference far greater than the data justify. All bars must have equal width.
\item \textbf{Unequal time intervals} on a line graph, which distort the apparent speed of change.
\item \textbf{Missing labels, scales or titles}: a graph without a scale cannot be interpreted at all.
\end{itemize}
The honest remedy: start frequency axes at 0, keep widths and intervals uniform, and label everything.
\end{attention}

\begin{aretenir}[Key points]
\begin{itemize}
\item Bar chart: equal widths, gaps between bars, height $=$ frequency; simple, multiple or stacked versions.
\item Pie chart: sector angle $\theta=\dfrac{f}{N}\times360^{\circ}$; the angles must sum to $360^{\circ}$.
\item Line graph: the diagram of a time series; time on the horizontal axis; reveals trends and seasonal patterns.
\item Choose the diagram according to the type of data and the message to convey.
\item Beware of truncated axes and unequal widths: they make small differences look large.
\end{itemize}
\end{aretenir}

\begin{exercice}[Pie chart and bar chart]
The 60 students of a school canteen were asked their preferred meal: \emph{ndol\'e} 18, \emph{eru} 15, rice and beans 12, \emph{achu} 9, \emph{koki} 6.
\begin{enumerate}
\item Compute the pie-chart angle for each meal and check your answer.
\item Draw a simple bar chart of the same data.
\item Which diagram better shows each meal's share of all preferences? Justify.
\end{enumerate}
\end{exercice}

\begin{corrige}[Exercise 1]
\begin{enumerate}
\item $N=60$, so $\theta=6f$ degrees: \emph{ndol\'e} $108^{\circ}$, \emph{eru} $90^{\circ}$, rice and beans $72^{\circ}$, \emph{achu} $54^{\circ}$, \emph{koki} $36^{\circ}$. Sum: $108+90+72+54+36=360^{\circ}$. $\checkmark$
\item Five bars of equal width, separated by gaps, of heights 18, 15, 12, 9, 6 on a frequency axis starting at 0.
\item The pie chart, because it displays shares of a total: each slice is directly the fraction $\frac{f}{60}$ of the whole, e.g. \emph{ndol\'e} represents $\frac{18}{60}=30\%$ of preferences.
\end{enumerate}
\end{corrige}

\coursec{Graphical Representation II: Histogram, Frequency Polygon and Ogive}

In the previous section we represented data with bar charts, pie charts and line charts. Those tools are excellent for \emph{discrete} or \emph{categorical} data: each bar stands apart, each sector of a pie shows a share of a whole. But what happens when the data are \cle{continuous} — heights, masses, examination marks measured on a continuous scale — and have been organised into a grouped frequency distribution? Separated bars would give the false impression that there are gaps between the classes. In this section we meet the three graphical tools designed for grouped continuous data: the \cle{histogram}, the \cle{frequency polygon} and the \cle{ogive} (cumulative frequency curve). The ogive is especially powerful: it allows us to \emph{read off} the median and the quartiles graphically, anticipating the calculations of Sections 6 and 8.

\courssub{The Histogram for Grouped Continuous Data}

\textbf{From observation to idea.}\par
A Biology teacher in Buea measures the heights of 50 Upper Sixth students and groups them into classes $150$--$155$, $155$--$160$, \dots, $175$--$180$ cm. If we draw a bar chart with gaps, we suggest that no student can have a height between $155.0$ and $155.1$ cm — which is absurd, since height is continuous. The natural fix is to let the rectangles \emph{touch} each other. The resulting diagram is called a \cle{histogram}.

\begin{definition}[Histogram]
A \cle{histogram} is a graphical representation of a grouped frequency distribution of a continuous variable, consisting of \textbf{adjacent rectangles} (no gaps) whose bases are the class intervals on the horizontal axis and whose \textbf{areas are proportional to the class frequencies}.
\end{definition}

\begin{attention}[Histogram or bar chart?]
Do not confuse the two. In a \textbf{bar chart} the bars are separated, the horizontal axis is categorical or discrete, and it is the \emph{height} of each bar that shows the frequency. In a \textbf{histogram} the bars touch, the horizontal axis is a continuous numerical scale, and it is the \emph{area} of each rectangle that represents the frequency. When all classes have the same width, height and area are proportional, so the two ideas coincide — but only then.
\end{attention}

\textbf{Equal class widths.}\par
When every class has the same width, we simply take
\[
\text{height of rectangle} = \text{frequency},
\]
because area $=$ width $\times$ height and the width is constant. The horizontal axis must be a true continuous scale, marked with the \emph{class boundaries} ($150, 155, 160, \dots$), not the class mid-points.

\begin{exemplebox}[Heights of 50 students]
The table below gives the grouped heights (in cm) of 50 students.
\begin{center}
\begin{tabularx}{\linewidth}{|l|X|X|X|X|X|X|}
\hline
\rowcolor{popblueL} Height (cm) & $150$--$155$ & $155$--$160$ & $160$--$165$ & $165$--$170$ & $170$--$175$ & $175$--$180$ \\
\hline
Frequency & $4$ & $10$ & $16$ & $12$ & $6$ & $2$ \\
\hline
\end{tabularx}
\end{center}
All classes have width $5$, so we plot rectangles of height equal to the frequency over each interval $[150,155]$, $[155,160]$, \dots
\end{exemplebox}

\begin{popfigure}
\begin{tikzpicture}
\begin{axis}[
  width=12cm, height=7cm,
  ybar, bar width=1.6cm,
  ymin=0, ymax=18,
  xmin=148, xmax=182,
  xlabel={Height (cm)}, ylabel={Frequency},
  xtick={150,155,160,165,170,175,180},
  ytick={0,2,4,6,8,10,12,14,16},
  grid=major, grid style={popline!40},
  axis lines=left,
  every axis plot/.append style={fill=popblueL, draw=popblue},
]
\addplot coordinates {(152.5,4) (157.5,10) (162.5,16) (167.5,12) (172.5,6) (177.5,2)};
\end{axis}
\end{tikzpicture}
\legende{Histogram of the heights of 50 students (equal class widths of $5$ cm).}
\end{popfigure}

\courssub{Unequal Class Widths: Frequency Density}

\textbf{Why height alone can mislead.}\par
Suppose one class is twice as wide as the others. If we draw its rectangle with height equal to its frequency, its \emph{area} will look twice as big as it should, and the eye judges areas, not heights. The honest solution is to make the \textbf{area} of each rectangle equal (or proportional) to the frequency. Since area $=$ class width $\times$ height, the height must be adjusted.

\begin{definition}[Frequency density]
For a class in a grouped distribution, the \cle{frequency density} is
\[
\text{frequency density} = \frac{\text{class frequency}}{\text{class width}}.
\]
In a histogram, the height of each rectangle is the frequency density, so that
\[
\text{area of rectangle} = \text{class width} \times \text{frequency density} = \text{frequency}.
\]
\end{definition}

\begin{methode}[Drawing a histogram with unequal class widths]
\begin{enumerate}
\item Check that the classes are continuous; if necessary, convert class limits to \textbf{class boundaries} (e.g.\ $10$--$19$, $20$--$29$ become $9.5$--$19.5$, $19.5$--$29.5$).
\item Compute the \textbf{width} of each class.
\item Compute the \textbf{frequency density} $= \dfrac{\text{frequency}}{\text{class width}}$ for each class.
\item Draw the horizontal axis as a continuous scale covering all boundaries; label the vertical axis ``frequency density''.
\item Over each class interval, draw a rectangle whose height is the frequency density. The rectangles must touch.
\end{enumerate}
\end{methode}

\begin{exoresolu}[Histogram with unequal classes]
The ages of members of a community club in Bamenda are summarised below. Compute the frequency densities and describe the histogram.
\begin{center}
\begin{tabularx}{\linewidth}{|l|X|X|X|X|}
\hline
\rowcolor{popblueL} Age (years) & $10$--$19$ & $20$--$29$ & $30$--$39$ & $40$--$59$ \\
\hline
Frequency & $20$ & $30$ & $25$ & $15$ \\
\hline
\end{tabularx}
\end{center}
\tcblower
\textbf{Solution.} Using boundaries, the classes are $9.5$--$19.5$, $19.5$--$29.5$, $29.5$--$39.5$ and $39.5$--$59.5$.
\begin{center}
\begin{tabularx}{\linewidth}{|l|X|X|X|}
\hline
\rowcolor{popblueL} Class boundaries & Width & Frequency & Frequency density \\
\hline
$9.5$--$19.5$ & $10$ & $20$ & $20/10 = 2.0$ \\
\hline
$19.5$--$29.5$ & $10$ & $30$ & $30/10 = 3.0$ \\
\hline
$29.5$--$39.5$ & $10$ & $25$ & $25/10 = 2.5$ \\
\hline
$39.5$--$59.5$ & $20$ & $15$ & $15/20 = 0.75$ \\
\hline
\end{tabularx}
\end{center}
The histogram has rectangles of heights $2.0$, $3.0$, $2.5$ and $0.75$ over the four intervals. Notice the last rectangle: although $15$ members is a substantial frequency, the class is twice as wide, so the rectangle is drawn \emph{low and wide}; its area, $20 \times 0.75 = 15$, correctly represents the frequency.
\end{exoresolu}

\courssub{The Frequency Polygon}

\begin{definition}[Frequency polygon]
A \cle{frequency polygon} is obtained from a histogram by plotting the points
\[
(\text{class mid-point},\ \text{frequency})
\]
and joining consecutive points by straight line segments. The polygon is usually closed by joining the end points to the mid-points of the empty classes immediately below the first class and above the last class (where the frequency is zero).
\end{definition}

The frequency polygon is ideal for \textbf{comparing two distributions on the same axes}: two histograms drawn together would overlap confusingly, but two polygons can be drawn in different colours and compared at a glance.

\begin{exemplebox}[Comparing two schools]
For the heights of the 50 students above, the mid-points are $152.5, 157.5, 162.5, 167.5, 172.5, 177.5$ with frequencies $4, 10, 16, 12, 6, 2$. Suppose a second school of 50 students gives frequencies $2, 6, 14, 16, 8, 4$ over the same classes. Plotting both polygons on the same axes shows immediately that the second school's heights are shifted slightly to the right: its students are on average taller.
\end{exemplebox}

\begin{popfigure}
\begin{tikzpicture}
\begin{axis}[
  width=12cm, height=7cm,
  ymin=0, ymax=18,
  xmin=145, xmax=185,
  xlabel={Height (cm)}, ylabel={Frequency},
  xtick={150,155,160,165,170,175,180},
  grid=major, grid style={popline!40},
  axis lines=left,
  legend style={at={(0.02,0.97)}, anchor=north west, font=\small},
]
\addplot[ybar, bar width=1.5cm, fill=popblueL, draw=popblue!60] coordinates {(152.5,4) (157.5,10) (162.5,16) (167.5,12) (172.5,6) (177.5,2)};
\addlegendentry{Histogram (School A)}
\addplot[mark=*, color=popdark, thick] coordinates {(147.5,0) (152.5,4) (157.5,10) (162.5,16) (167.5,12) (172.5,6) (177.5,2) (182.5,0)};
\addlegendentry{Polygon (School A)}
\addplot[mark=square*, color=poporange, thick, dashed] coordinates {(147.5,0) (152.5,2) (157.5,6) (162.5,14) (167.5,16) (172.5,8) (177.5,4) (182.5,0)};
\addlegendentry{Polygon (School B)}
\end{axis}
\end{tikzpicture}
\legende{Frequency polygons of two schools superimposed on the histogram of School A.}
\end{popfigure}

\begin{aretenir}[Frequency polygon — key points]
\begin{itemize}
\item Points are plotted at \textbf{class mid-points}, joined by straight lines.
\item The polygon is closed at both ends on the horizontal axis.
\item The area under the frequency polygon equals the area under the corresponding histogram.
\item It is the best tool for comparing two (or more) distributions with the same classes.
\end{itemize}
\end{aretenir}

\courssub{Cumulative Frequency and the Ogive}

\textbf{From observation to idea.}\par
A school principal asks: ``How many candidates are \emph{shorter than} $165$ cm? \emph{Shorter than} $170$ cm?'' To answer such ``less than'' questions quickly, we accumulate the frequencies as we move up the classes. This gives the \cle{cumulative frequency distribution}.

\begin{definition}[Cumulative frequency]
The \cle{cumulative frequency} of a class is the total frequency of all values \textbf{less than the upper class boundary} of that class. It is obtained by adding successively the class frequencies.
\end{definition}

\begin{exemplebox}[Cumulative frequency table for the heights]
For the 50 students (boundaries $150, 155, \dots, 180$):
\begin{center}
\begin{tabularx}{\linewidth}{|l|X|X|}
\hline
\rowcolor{popblueL} Height (cm) & Frequency & Cumulative frequency \\
\hline
$150$--$155$ & $4$ & $4$ \\
\hline
$155$--$160$ & $10$ & $14$ \\
\hline
$160$--$165$ & $16$ & $30$ \\
\hline
$165$--$170$ & $12$ & $42$ \\
\hline
$170$--$175$ & $6$ & $48$ \\
\hline
$175$--$180$ & $2$ & $50$ \\
\hline
\end{tabularx}
\end{center}
Each cumulative frequency is plotted against the \textbf{upper class boundary}: $(155,4)$, $(160,14)$, $(165,30)$, $(170,42)$, $(175,48)$, $(180,50)$, together with the starting point $(150,0)$.
\end{exemplebox}

\begin{definition}[Ogive]
The \cle{ogive} (or cumulative frequency curve) is the curve obtained by plotting cumulative frequency against the \textbf{upper class boundaries} and joining the points (by straight segments for a polygonal ogive, or by a smooth S-shaped curve).
\end{definition}

\begin{methode}[Constructing an ogive from a cumulative frequency table]
\begin{enumerate}
\item Write down the \textbf{upper class boundary} of every class.
\item Build the cumulative frequency column by successive addition; check that the last value equals $N$, the total frequency.
\item Plot the points $(\text{upper boundary},\ \text{cumulative frequency})$, and add the point $(\text{lower boundary of first class},\ 0)$.
\item Join the points with straight segments (or a smooth curve). The ogive always rises from $0$ to $N$ and never decreases.
\item Label both axes clearly: the vertical axis is ``cumulative frequency''.
\end{enumerate}
\end{methode}

\begin{attention}[The most common mistake]
Never plot cumulative frequency against the \textbf{class mid-point}. A cumulative frequency answers the question ``how many observations are \emph{less than} this value?'', and that value must be the \textbf{upper class boundary}. Plotting against mid-points shifts the whole ogive half a class to the left and gives wrong medians and quartiles. Examiners penalise this heavily.
\end{attention}

\begin{popfigure}
\begin{tikzpicture}
\begin{axis}[
  width=12cm, height=8cm,
  ymin=0, ymax=55,
  xmin=148, xmax=182,
  xlabel={Height (cm)}, ylabel={Cumulative frequency},
  xtick={150,155,160,165,170,175,180},
  ytick={0,12.5,25,37.5,50},
  grid=major, grid style={popline!40},
  axis lines=left,
]
\addplot[mark=*, color=popblue, very thick] coordinates {(150,0) (155,4) (160,14) (165,30) (170,42) (175,48) (180,50)};
\addplot[dashed, color=poporange, thick] coordinates {(148,25) (163.4,25)};
\addplot[dashed, color=poporange, thick] coordinates {(163.4,25) (163.4,0)};
\addplot[dashed, color=popgreen, thick] coordinates {(148,12.5) (159.3,12.5)};
\addplot[dashed, color=popgreen, thick] coordinates {(159.3,12.5) (159.3,0)};
\addplot[dashed, color=poppurple, thick] coordinates {(148,37.5) (168.1,37.5)};
\addplot[dashed, color=poppurple, thick] coordinates {(168.1,37.5) (168.1,0)};
\node[color=popgreen, font=\small, anchor=west] at (axis cs:150.5,14.5) {$N/4=12.5$};
\node[color=poporange, font=\small, anchor=west] at (axis cs:150.5,27) {$N/2=25$};
\node[color=poppurple, font=\small, anchor=west] at (axis cs:150.5,39.5) {$3N/4=37.5$};
\node[color=popgreen, font=\small, anchor=north] at (axis cs:159.3,1) {$Q_1$};
\node[color=poporange, font=\small, anchor=north] at (axis cs:163.4,1) {$Q_2$};
\node[color=poppurple, font=\small, anchor=north] at (axis cs:168.1,1) {$Q_3$};
\end{axis}
\end{tikzpicture}
\legende{Ogive of the heights, with readings at $N/4$, $N/2$ and $3N/4$ giving $Q_1$, the median $Q_2$ and $Q_3$.}
\end{popfigure}

\courssub{Reading the Median and Quartiles from the Ogive}

The ogive turns ``less than'' questions around: instead of asking ``how many students are below $x$?'', we ask ``below which height do we find half the students?''. That height is the \cle{median}.

\begin{propriete}[Graphical determination of median and quartiles]
For a distribution of $N$ observations, draw the ogive and then:
\begin{itemize}
\item locate $\dfrac{N}{2}$ on the cumulative frequency axis, draw a horizontal line to the ogive, then a vertical line down: the value read is the \textbf{median} $Q_2$;
\item locate $\dfrac{N}{4}$: the value read is the \textbf{lower quartile} $Q_1$;
\item locate $\dfrac{3N}{4}$: the value read is the \textbf{upper quartile} $Q_3$.
\end{itemize}
The \textbf{interquartile range} is then $\mathrm{IQR} = Q_3 - Q_1$. (Proof not examinable; the method follows directly from the definitions of Section 8.)
\end{propriete}

\begin{exoresolu}[Reading the median and quartiles]
From the ogive of the 50 heights, estimate the median, the quartiles and the interquartile range.
\tcblower
\textbf{Solution.} Here $N = 50$, so
\[
\frac{N}{4} = 12.5, \qquad \frac{N}{2} = 25, \qquad \frac{3N}{4} = 37.5.
\]
Drawing horizontal lines at $12.5$, $25$ and $37.5$ up to the curve and projecting down to the height axis (see the figure), we read approximately
\[
Q_1 \approx 159.3\ \text{cm}, \qquad Q_2 \approx 163.4\ \text{cm}, \qquad Q_3 \approx 168.1\ \text{cm}.
\]
These readings are consistent with linear interpolation on the table: for instance $Q_1$ lies in the class $155$--$160$ (cumulative frequencies $4$ to $14$), giving $155 + \dfrac{12.5-4}{10}\times 5 = 159.25 \approx 159.3$ cm, and $Q_3$ lies in the class $165$--$170$ (cumulative frequencies $30$ to $42$), giving $165 + \dfrac{37.5-30}{12}\times 5 = 168.125 \approx 168.1$ cm. Hence
\[
\mathrm{IQR} = Q_3 - Q_1 \approx 168.1 - 159.3 = 8.8\ \text{cm}.
\]
Interpretation: about half of the students have heights between roughly $159.3$ cm and $168.1$ cm, and the middle $50\ \%$ of heights spread over about $8.8$ cm. These graphical estimates will be confirmed by calculation in Sections 6 and 8.
\end{exoresolu}

\begin{attention}[Reading errors to avoid]
\begin{itemize}
\item Use $N/2$, not $(N+1)/2$, when reading from the ogive: the ogive is a continuous-curve estimate.
\item Always read the horizontal line \emph{from the cumulative axis to the curve}, then drop \emph{vertically} to the variable axis — never the reverse.
\item State your answers as \emph{estimates} (use $\approx$): graphical readings are never exact.
\end{itemize}
\end{attention}

\begin{savaistu}[Why ``ogive''?]
The word \emph{ogive} comes from architecture: it names the pointed, S-shaped arch used in Gothic cathedrals. Statisticians borrowed the word because cumulative frequency curves have exactly that graceful S-shape — flat at the extremes where few data accumulate, and steep in the dense middle classes. The steeper the ogive over a class, the more concentrated the data are there.
\end{savaistu}

\begin{aretenir}[Chapter summary — what you must master]
\begin{itemize}
\item A \textbf{histogram} has touching rectangles whose \emph{areas} represent frequencies; with unequal class widths, plot \textbf{frequency density} $=$ frequency $\div$ class width.
\item A \textbf{frequency polygon} joins the points (mid-point, frequency) and is ideal for comparing distributions.
\item An \textbf{ogive} plots cumulative frequency against \textbf{upper class boundaries}; it rises from $0$ to $N$.
\item From the ogive, read $Q_1$ at $N/4$, the median at $N/2$ and $Q_3$ at $3N/4$; then $\mathrm{IQR} = Q_3 - Q_1$.
\end{itemize}
\end{aretenir}

\begin{exercice}[Practice]
The masses (in kg) of 60 bags of beans at a market in Mbouda are grouped as follows:
\begin{center}
\begin{tabularx}{\linewidth}{|l|X|X|X|X|X|}
\hline
\rowcolor{popblueL} Mass (kg) & $40$--$45$ & $45$--$50$ & $50$--$55$ & $55$--$60$ & $60$--$70$ \\
\hline
Frequency & $6$ & $15$ & $20$ & $12$ & $7$ \\
\hline
\end{tabularx}
\end{center}
\begin{enumerate}
\item Compute the frequency density of each class and describe the histogram.
\item Construct the cumulative frequency table.
\item State the coordinates of the points to plot for the ogive.
\item Given that the ogive passes through $(50, 21)$ and $(55, 41)$, estimate the median by linear interpolation between the appropriate points.
\end{enumerate}
\end{exercice}

\begin{corrige}[Practice]
\textbf{1.} Widths: $5, 5, 5, 5, 10$. Frequency densities: $6/5 = 1.2$; $15/5 = 3.0$; $20/5 = 4.0$; $12/5 = 2.4$; $7/10 = 0.7$. The histogram has touching rectangles of these heights; the last rectangle is twice as wide but low, its area $10 \times 0.7 = 7$ correctly showing the frequency.

\textbf{2.} Cumulative frequencies: $6$, $21$, $41$, $53$, $60$.

\textbf{3.} Plot $(40,0)$, $(45,6)$, $(50,21)$, $(55,41)$, $(60,53)$, $(70,60)$ — cumulative frequency against \emph{upper} class boundaries.

\textbf{4.} $N/2 = 30$, which lies between $21$ (at $50$ kg) and $41$ (at $55$ kg). By linear interpolation,
\[
\text{median} \approx 50 + \frac{30 - 21}{41 - 21}\times 5 = 50 + \frac{9}{20}\times 5 = 52.25\ \text{kg}.
\]
So the median mass of a bag of beans is approximately $52.25$ kg.
\end{corrige}

\coursec{Measures of Central Tendency I: Mode, Median, Range}

In the previous section we learned how to \emph{display} a distribution: histograms, frequency polygons and ogives give us a picture of the data. But a picture alone is not enough for the GCE examiner, nor for the economist comparing salaries in Douala and Yaoundé. We need \emph{numbers} that summarise a whole distribution. The first family of such numbers are the \cle{measures of central tendency} (or \emph{averages}): they tell us where the \emph{centre} of the data lies. In this section we study three of them: the \cle{mode}, the \cle{median}, and the simplest measure of spread, the \cle{range}.

\courssub{The Mode of Ungrouped Data}

Imagine a shoe seller at Mokolo market who records the sizes of the last $20$ pairs of shoes sold. To restock wisely, she does not need the ``average'' size in some complicated sense: she needs the size that \emph{sells most often}. That value is the mode.

\begin{definition}[Mode]
The \cle{mode} of a set of data is the value that occurs with the \textbf{greatest frequency}.
\begin{itemize}
\item A distribution with exactly one mode is called \cle{unimodal}.
\item A distribution with two values sharing the greatest frequency is called \cle{bimodal}.
\item If every value occurs the same number of times (or no value repeats), the data set has \textbf{no mode}.
\end{itemize}
\end{definition}

\begin{exemplebox}[Finding the mode]
The marks of $15$ candidates in a quiz out of $10$ are:
\[ 4,\ 7,\ 5,\ 7,\ 6,\ 7,\ 5,\ 8,\ 7,\ 4,\ 6,\ 7,\ 5,\ 9,\ 7. \]
Counting frequencies: $4$ appears $2$ times, $5$ appears $3$ times, $6$ appears $2$ times, $7$ appears $6$ times, $8$ once, $9$ once. The greatest frequency is $6$, so
\[ \text{Mode} = 7. \]
Had both $5$ and $7$ appeared $6$ times each, the data would be \emph{bimodal} with modes $5$ and $7$.
\end{exemplebox}

\begin{attention}[Mode is a value, not a frequency]
A very common error is to give the \emph{frequency} as the answer. In the example above the mode is $7$ (the mark), \textbf{not} $6$ (its frequency). The mode is always a \emph{data value}.
\end{attention}

\courssub{The Modal Class of Grouped Data}

When data are grouped into classes, individual values are lost, so we can no longer read off a mode exactly. We proceed in two stages.

\begin{definition}[Modal class]
For grouped data, the \cle{modal class} is the class with the \textbf{highest frequency} (strictly, the highest frequency \emph{density} when classes have unequal widths; for equal-width classes, the highest frequency suffices).
\end{definition}

If a single estimate of the mode is required from grouped data, we may locate it inside the modal class by \emph{linear interpolation}. The idea: the mode is pulled towards the ``stronger'' neighbour. If the class before the modal class has frequency $f_{-1}$ and the class after has frequency $f_{+1}$, the mode divides the modal class in the ratio $(f_m - f_{-1}) : (f_m - f_{+1})$ measured from the lower boundary.

\begin{propriete}[Mode of grouped data by interpolation — GCE extra]
For a modal class with lower class boundary $L$, class width $c$, modal frequency $f_m$, and neighbouring frequencies $f_{-1}$ (class before) and $f_{+1}$ (class after):
\[ \text{Mode} \approx L + \frac{f_m - f_{-1}}{(f_m - f_{-1}) + (f_m - f_{+1})} \times c. \]
This formula is an \textbf{estimate}; it is not always required at GCE, but it earns marks when ``estimate the mode'' is asked for grouped data.
\end{propriete}

\begin{exemplebox}[Estimating the mode from grouped data]
The heights of $60$ students give: modal class $160$--$165$ cm with $f_m = 20$, previous class frequency $f_{-1} = 14$, next class frequency $f_{+1} = 12$, width $c = 5$, lower boundary $L = 159.5$.
\[ \text{Mode} \approx 159.5 + \frac{20-14}{(20-14)+(20-12)} \times 5 = 159.5 + \frac{6}{14}\times 5 \approx 161.6\ \mathrm{cm}. \]
\end{exemplebox}

\courssub{The Median of Ungrouped Data}

Suppose nine workers in a small company in Bamenda earn monthly salaries (in thousands of FCFA):
\[ 50,\ 55,\ 60,\ 65,\ 70,\ 75,\ 80,\ 85,\ 500. \]
The last value is the manager's salary. The ``typical'' salary is badly distorted by that extreme value if we average naively. The \emph{median} solves this: it is simply the \textbf{middle value} once the data are arranged in order.

\begin{definition}[Median]
The \cle{median} of a set of $n$ observations arranged in \textbf{ascending order} is the middle value:
\begin{itemize}
\item if $n$ is \textbf{odd}, the median is the value in position $\dfrac{n+1}{2}$;
\item if $n$ is \textbf{even}, the median is the \textbf{mean of the two values} in positions $\dfrac{n}{2}$ and $\dfrac{n}{2}+1$.
\end{itemize}
\end{definition}

\begin{exemplebox}[Odd and even $n$]
\textbf{(a) Odd.} For the nine salaries above, already ordered, position $\frac{9+1}{2} = 5$ gives median $= 70$ thousand FCFA. Notice the manager's $500$ has no influence.

\textbf{(b) Even.} For the marks $3, 5, 6, 8, 9, 12$ ($n = 6$), the middle positions are $3$rd and $4$th:
\[ \text{Median} = \frac{6+8}{2} = 7. \]
Note that $7$ is not itself a data value — this is perfectly acceptable for a median.
\end{exemplebox}

\begin{attention}[Order first!]
The single most frequent mistake is to pick the middle value of the data \emph{as given}, without arranging them in ascending order. Always \textbf{sort first}, then locate the middle.
\end{attention}

\courssub{The Median of Grouped Data by Interpolation}

With grouped data we no longer know individual values, but we know how many observations fall below each class boundary — exactly the information contained in the \emph{cumulative frequency} table and the ogive of the previous section. The median is the value below which $50\%$ of the observations lie, i.e.\ the value of the $\frac{N}{2}$-th observation, where $N$ is the total frequency.

\begin{methode}[Locating the median and the median class]
\begin{enumerate}
\item Compute the total frequency $N$ and form the \textbf{cumulative frequency} column.
\item Compute $\dfrac{N}{2}$ (do \emph{not} add $1$, even for odd $N$, when data are grouped).
\item Read down the cumulative frequencies until you find the \textbf{first} cumulative frequency $\geq \frac{N}{2}$: the corresponding class is the \cle{median class}.
\item Apply the interpolation formula below.
\end{enumerate}
\end{methode}

\begin{propriete}[Median of grouped data by linear interpolation]
Let $L$ be the \textbf{lower class boundary} of the median class, $F$ the cumulative frequency \textbf{up to the end of the class before} the median class, $f$ the frequency of the median class, $c$ its width, and $N$ the total frequency. Then
\[ \text{Median} = L + \left( \frac{\frac{N}{2} - F}{f} \right) \times c. \]
\emph{Justification (examinable as understanding):} inside the median class we assume the $f$ observations are spread \emph{uniformly}. We need to climb from $F$ observations (at the lower boundary $L$) up to $\frac{N}{2}$ observations, i.e.\ a fraction $\frac{N/2 - F}{f}$ of the way through the class. Multiplying this fraction by the class width $c$ gives the distance to add to $L$.
\end{propriete}

\begin{exoresolu}[Median of grouped data]
The table shows the masses (kg) of $50$ bags of cement sampled at a depot.
\begin{center}
\begin{tabular}{|c|c|c|}
\hline
\rowcolor{popblueL} Mass (kg) & Frequency $f$ & Cumulative frequency \\
\hline
$45$--$47$ & $4$ & $4$ \\
\hline
$48$--$50$ & $10$ & $14$ \\
\hline
$51$--$53$ & $17$ & $31$ \\
\hline
$54$--$56$ & $12$ & $43$ \\
\hline
$57$--$59$ & $7$ & $50$ \\
\hline
\end{tabular}
\end{center}
Estimate the median mass.
\tcblower
\textbf{Step 1.} $N = 50$, so $\dfrac{N}{2} = 25$.

\textbf{Step 2.} The first cumulative frequency $\geq 25$ is $31$, so the median class is $51$--$53$.

\textbf{Step 3.} Class boundaries: $L = 50.5$, width $c = 3$; cumulative frequency before the median class $F = 14$; median class frequency $f = 17$.

\textbf{Step 4.} Apply the formula:
\begin{align*}
\text{Median} &= 50.5 + \left( \frac{25 - 14}{17} \right) \times 3 \\
&= 50.5 + \frac{11}{17} \times 3 \\
&= 50.5 + 1.941\ldots \\
&\approx 52.44\ \mathrm{kg}.
\end{align*}
Half of the bags weigh less than about $52.44$ kg, and half weigh more.
\end{exoresolu}

\begin{attention}[Using $N/2$ of the number of classes]
A classic trap: some candidates compute $\frac{N}{2}$ using the \emph{number of classes} (here $\frac{5}{2} = 2.5$) or halve the \emph{class frequencies}. The $\frac{N}{2}$ in the formula is always \textbf{half the total frequency} (here $\frac{50}{2} = 25$). Also, $F$ is the cumulative frequency \emph{before} the median class, not that of the median class itself, and $L$ is a class \emph{boundary} ($50.5$), not a class limit ($51$).
\end{attention}

\begin{popfigure}
\begin{tikzpicture}[scale=1.05]
\draw[popline, dashed] (0,0) rectangle (1,3);
\node[below, font=\small] at (0.5,0) {before};
\draw[popblue, thick] (1,0) rectangle (4,3);
\node[below, font=\small, text=popblue] at (2.5,0) {median class};
\draw[popline, dashed] (4,0) rectangle (5,3);
\node[below, font=\small] at (4.5,0) {after};
\draw[<->, thick] (1,-0.7) -- (4,-0.7) node[midway, below, font=\small] {$c$};
\draw[<->, thick] (5.3,0) -- (5.3,3) node[midway, right, font=\small] {$f$};
\fill[poppink] (1,0) circle (2pt) node[below left, font=\small, text=black] {$L$};
\fill[popgreen] (2.94,0) circle (2pt) node[below, font=\small, text=black, xshift=6pt] {median};
\draw[<->, popgreen, thick] (1,3.5) -- (2.94,3.5) node[midway, above, font=\small] {$\dfrac{N/2 - F}{f}\times c$};
\end{tikzpicture}
\legende{Interpolation inside the median class: the median sits a fraction $\frac{N/2 - F}{f}$ of the class width above the lower boundary $L$.}
\end{popfigure}

\courssub{Graphical Estimation of the Median from the Ogive}

Because the ogive plots cumulative frequency against the \textbf{upper class boundaries}, the median is read directly at cumulative frequency $\frac{N}{2}$: draw a horizontal line from $\frac{N}{2}$ on the vertical axis to the curve, then drop vertically to the horizontal axis. This is exactly the graphical counterpart of the interpolation formula, and GCE examiners accept either approach — but the graph gives an \emph{estimate} whose accuracy depends on your drawing.

\begin{popfigure}
\begin{tikzpicture}
\begin{axis}[
  width=11cm, height=7.5cm,
  xlabel={Mass (kg)}, ylabel={Cumulative frequency},
  xmin=44, xmax=60, ymin=0, ymax=55,
  xtick={47.5,50.5,53.5,56.5,59.5},
  ytick={0,10,20,25,30,40,50},
  grid=both, grid style={popline, very thin},
  axis line style={popdark},
]
\addplot[popblue, very thick, mark=*, mark size=1.5pt] coordinates {(47.5,4) (50.5,14) (53.5,31) (56.5,43) (59.5,50)};
\addplot[poppink, thick, dashed] coordinates {(44,25) (52.44,25)};
\addplot[poppink, thick, dashed] coordinates {(52.44,25) (52.44,0)};
\node[font=\small, text=poppink, anchor=south east] at (axis cs:45,25) {$N/2 = 25$};
\node[font=\small, text=poppink, anchor=north] at (axis cs:52.44,0) {$\approx 52.4$};
\fill[poppink] (axis cs:52.44,25) circle (2pt);
\end{axis}
\end{tikzpicture}
\legende{Reading the median from the ogive: locate $\frac{N}{2}=25$ on the cumulative axis, meet the curve, and read down to obtain the median $\approx 52.4$ kg. Points are plotted at upper class boundaries.}
\end{popfigure}

\courssub{The Range}

\begin{definition}[Range]
The \cle{range} of a set of data is the difference between the greatest and the least values:
\[ \text{Range} = \text{maximum value} - \text{minimum value}. \]
\end{definition}

The range is a measure of \emph{dispersion} (spread), not of centre, but it is traditionally studied alongside the mode and median because of its simplicity. For the nine Bamenda salaries, the range is $500 - 50 = 450$ thousand FCFA.

\textbf{Uses.} The range is quick to compute and easy to understand; it is used in quality control (e.g.\ checking that the masses of bags of cement stay within acceptable limits) and in meteorology (daily temperature range).

\textbf{Limitations.} The range depends only on the \emph{two extreme values} and ignores everything in between. Worse, it is extremely sensitive to \cle{outliers}: without the manager, the salary range collapses from $450$ to $35$ thousand FCFA. Two distributions can have the same range and completely different shapes. This is why more robust measures of spread (interquartile range, variance, standard deviation) are developed later in this chapter.

\courssub{Comparing the Mode and the Median: Which Average?}

\begin{center}
\begin{tabularx}{\linewidth}{|l|X|X|}
\hline
\rowcolor{popblueL} \textbf{Average} & \textbf{Strengths} & \textbf{Best used when} \\
\hline
Mode & Easy; unaffected by extreme values; can be used for \emph{non-numerical} data (most popular shoe size, most common blood group) & The ``most typical'' or most frequent category is wanted; data are categorical \\
\hline
Median & Unaffected by outliers; always splits data into two equal halves & Data contain extreme values (salaries, house prices) or are skewed \\
\hline
\end{tabularx}
\end{center}

The mode may not exist, or may be far from the centre (a bimodal distribution has no single ``typical'' value). The median always exists and resists outliers, but it ignores the actual sizes of the values — only their order matters. When every value should count equally, we need the \emph{mean}, studied in the next section.

\begin{aretenir}[Key points]
\begin{itemize}
\item \textbf{Mode} $=$ most frequent value (a data value, not a frequency); data may be unimodal, bimodal, or have no mode.
\item \textbf{Modal class} $=$ class of highest frequency; mode of grouped data by interpolation: $L + \frac{f_m - f_{-1}}{(f_m - f_{-1})+(f_m - f_{+1})}\,c$.
\item \textbf{Median} of ungrouped data: middle value after ordering; position $\frac{n+1}{2}$ (odd $n$), mean of positions $\frac{n}{2}$ and $\frac{n}{2}+1$ (even $n$).
\item \textbf{Median of grouped data}: $\text{Median} = L + \left(\frac{N/2 - F}{f}\right)c$, where $N$ is the \emph{total frequency}.
\item From the ogive, read the median at cumulative frequency $\frac{N}{2}$ (plot against upper class boundaries).
\item \textbf{Range} $=$ max $-$ min; simple but very sensitive to outliers.
\end{itemize}
\end{aretenir}

\begin{exercice}[Practice]
The table gives the times (minutes) taken by $40$ students to solve a problem.
\begin{center}
\begin{tabular}{|c|c|c|c|c|c|}
\hline
\rowcolor{popblueL} Time (min) & $5$--$9$ & $10$--$14$ & $15$--$19$ & $20$--$24$ & $25$--$29$ \\
\hline
Frequency & $5$ & $9$ & $14$ & $8$ & $4$ \\
\hline
\end{tabular}
\end{center}
\begin{enumerate}
\item State the modal class and estimate the mode by interpolation.
\item Construct the cumulative frequency table and estimate the median by interpolation.
\item Estimate the median graphically from an ogive and compare with (b).
\item Explain why the range cannot be computed exactly from this table, and give the largest possible value it could take.
\end{enumerate}
\end{exercice}

\begin{corrige}[Practice]
\textbf{1.} Modal class: $15$--$19$ ($f_m = 14$, $f_{-1} = 9$, $f_{+1} = 8$, $L = 14.5$, $c = 5$).
\[ \text{Mode} \approx 14.5 + \frac{14-9}{(14-9)+(14-8)} \times 5 = 14.5 + \frac{5}{11}\times 5 \approx 16.8\ \mathrm{min}. \]
\textbf{2.} Cumulative frequencies: $5, 14, 28, 36, 40$. $N = 40$, $\frac{N}{2} = 20$. First cumulative frequency $\geq 20$ is $28$: median class $15$--$19$, with $L = 14.5$, $F = 14$, $f = 14$, $c = 5$.
\[ \text{Median} = 14.5 + \frac{20-14}{14}\times 5 = 14.5 + \frac{30}{14} \approx 16.6\ \mathrm{min}. \]
\textbf{3.} Plotting $(9.5,5), (14.5,14), (19.5,28), (24.5,36), (29.5,40)$ and reading across from $20$ on the cumulative axis gives a median of about $16.6$ min, agreeing with (b) up to drawing accuracy.
\textbf{4.} Individual values are unknown after grouping, so the exact maximum and minimum are lost. The largest possible range would be (greatest possible value) $-$ (least possible value) $= 29 - 5 = 24$ minutes (using class limits; $29.5 - 4.5 = 25$ using boundaries).
\end{corrige}

\coursec{Measures of Central Tendency II: Mean and Combined Mean}

In the previous section we met the \cle{mode}, the \cle{median} and the \cle{range}: the mode tells us the most frequent value, the median tells us the middle value, and the range tells us how spread out the data are. Yet none of these measures uses \emph{every} observation. If a trader in Bamenda records his daily sales for a month, the mode ignores all days except the most frequent one, and the median ignores the actual sizes of the sales. We now introduce the most famous average of all, the \cle{arithmetic mean}, which takes every single value into account --- and we study how to combine means of several groups correctly.

\courssub{The Arithmetic Mean of Ungrouped Data}

\textbf{Intuition: sharing equally.}\par
Suppose five students score $8$, $10$, $12$, $14$ and $16$ marks out of $20$ in a class test. If the marks were pooled and shared equally among the five students, each would receive
\[
\frac{8+10+12+14+16}{5}=\frac{60}{5}=12.
\]
This ``equal share'' is exactly the arithmetic mean. It is the value each observation \emph{would} have if the total were redistributed fairly.

\begin{definition}[Arithmetic mean of ungrouped data]
The \cle{arithmetic mean} of $n$ observations $x_1, x_2, \dots, x_n$ is
\[
\bar{x}=\frac{x_1+x_2+\cdots+x_n}{n}=\frac{\sum x}{n}.
\]
If the value $x_i$ occurs with frequency $f_i$ (discrete frequency table), then
\[
\bar{x}=\frac{\sum f_i x_i}{\sum f_i}.
\]
\end{definition}

\textbf{The mean as a balance point.}\par
A beautiful way to picture the mean is to imagine the number line as a light bar with equal unit weights placed at each data point. The bar balances exactly at the mean, because the sum of the signed deviations from the mean is zero:
\[
\sum (x_i-\bar{x})=\sum x_i - n\bar{x}=\sum x_i - \sum x_i = 0.
\]
This property is fundamental and is the reason the mean appears everywhere in statistics.

\begin{popfigure}
\begin{tikzpicture}[scale=1.0]
  \draw[thick, popline] (0,0) -- (10,0);
  \foreach \x/\lab in {1/8, 3/10, 5/12, 7/14, 9/16}{
    \draw[popline] (\x,-0.12) -- (\x,0.12);
    \node[below, popmuted] at (\x,-0.15) {\small \lab};
    \fill[popblue] (\x,0.35) circle (0.14);
  }
  \fill[poporange] (5,0) -- (4.75,-0.55) -- (5.25,-0.55) -- cycle;
  \node[below, poporange] at (5,-0.6) {\small $\bar{x}=12$};
  \draw[<->, poppurple, thick] (1,1.0) -- (5,1.0) node[midway, above]{\small $-4$};
  \draw[<->, poppurple, thick] (5,1.0) -- (9,1.0) node[midway, above]{\small $+4$};
\end{tikzpicture}
\legende{The mean as a balance point: deviations to the left and to the right cancel out, since $\sum(x_i-\bar{x})=0$.}
\end{popfigure}

\begin{exemplebox}[Mean of a discrete frequency table]
The number of goals scored by a school team in $10$ matches is summarised below. Find the mean number of goals per match.
\begin{center}
\begin{tabular}{|c|ccccc|}
\hline
\rowcolor{popblueL} Goals $x$ & $0$ & $1$ & $2$ & $3$ & $4$ \\
\hline
Frequency $f$ & $2$ & $3$ & $3$ & $1$ & $1$ \\
\hline
$fx$ & $0$ & $3$ & $6$ & $3$ & $4$ \\
\hline
\end{tabular}
\end{center}
Here $\sum f = 10$ and $\sum fx = 16$, so
\[
\bar{x}=\frac{\sum fx}{\sum f}=\frac{16}{10}=1.6 \text{ goals per match.}
\]
Note that the mean need not be one of the observed values, nor even a whole number.
\end{exemplebox}

\begin{methode}[Finding the mean from raw data or a discrete frequency table]
\begin{enumerate}
\item List the values $x_i$ and their frequencies $f_i$ (if raw data, each value has frequency $1$).
\item Compute the product $f_i x_i$ for each row.
\item Find $\sum f_i$ (total number of observations) and $\sum f_i x_i$ (sum of all values).
\item Apply $\bar{x} = \dfrac{\sum f_i x_i}{\sum f_i}$.
\end{enumerate}
\end{methode}

\courssub{Mean of Grouped Data Using Class Mid-Points}

When data are grouped into classes, the individual values are lost. We therefore \emph{estimate} the mean by pretending that every observation in a class sits at the \cle{class mid-point} (also called \cle{class mark}). For a class with lower limit $L$ and upper limit $U$, the mid-point is $x=\dfrac{L+U}{2}$.

\begin{attention}[Class limits vs. class boundaries]
If the classes are given as $45$--$47$, $48$--$50$, etc., there is a gap of $1$ unit between them. For continuous data (like mass), the true class boundaries are $44.5$--$47.5$, $47.5$--$50.5$, etc. The mid-point calculated from the \emph{stated limits} $(45+47)/2 = 46$ is exactly the same as from the boundaries $(44.5+47.5)/2 = 46$. In GCE examinations, always use the stated limits to compute the mid-point unless class boundaries are explicitly given.
\end{attention}

\begin{propriete}[Mean of grouped data]
For a grouped frequency distribution with class mid-points $x_i$ and frequencies $f_i$,
\[
\bar{x}=\frac{\sum f_i x_i}{\sum f_i}.
\]
This is an \emph{estimate} of the true mean; the formula itself must be quoted in the examination.
\end{propriete}

\begin{methode}[Estimating the mean from a grouped frequency table]
\begin{enumerate}
\item For each class, compute the mid-point $x = \dfrac{\text{lower limit} + \text{upper limit}}{2}$.
\item Multiply each mid-point by its class frequency to get $fx$.
\item Sum the frequencies ($\sum f$) and the $fx$ column ($\sum fx$).
\item Calculate $\bar{x} = \dfrac{\sum fx}{\sum f}$.
\item State the units and remember this is an \emph{estimate}.
\end{enumerate}
\end{methode}

\begin{exoresolu}[Mean from a grouped table]
The masses (in kg) of $40$ bags of cement checked at a depot are:
\begin{center}
\begin{tabular}{|c|ccccc|}
\hline
\rowcolor{popblueL} Mass (kg) & $45$--$47$ & $48$--$50$ & $51$--$53$ & $54$--$56$ & $57$--$59$ \\
\hline
Frequency $f$ & $4$ & $10$ & $14$ & $8$ & $4$ \\
\hline
\end{tabular}
\end{center}
Estimate the mean mass.
\tcblower
The class mid-points are $46, 49, 52, 55, 58$. We build the working table:
\begin{center}
\begin{tabular}{|c|c|c|c|}
\hline
\rowcolor{popblueL} Class & Mid-point $x$ & $f$ & $fx$ \\
\hline
$45$--$47$ & $46$ & $4$ & $184$ \\
$48$--$50$ & $49$ & $10$ & $490$ \\
$51$--$53$ & $52$ & $14$ & $728$ \\
$54$--$56$ & $55$ & $8$ & $440$ \\
$57$--$59$ & $58$ & $4$ & $232$ \\
\hline
Totals & & $\sum f = 40$ & $\sum fx = 2074$ \\
\hline
\end{tabular}
\end{center}
\[
\bar{x}=\frac{\sum fx}{\sum f}=\frac{2074}{40}=51.85\ \mathrm{kg}.
\]
\end{exoresolu}

\courssub{The Assumed-Mean (Coding) Method}

When the mid-points are large or awkward, the products $fx$ become heavy. The \cle{coding method} (or \cle{assumed-mean method}) lightens the arithmetic by measuring every mid-point from a chosen \cle{assumed mean} $A$ (usually the mid-point of a central class, often the modal class).

\begin{methode}[Assumed-mean (coding) technique]
\begin{enumerate}
\item Choose an assumed mean $A$, usually the mid-point of the class with the largest frequency.
\item For each class compute the deviation $d = x - A$ (or the step deviation $d=\dfrac{x-A}{c}$ when classes have common width $c$).
\item Compute $fd$ for each class and sum to get $\sum fd$.
\item Apply the formula:
\[
\bar{x}=A+\frac{\sum fd}{\sum f}
\qquad\text{or, with step }c:\qquad
\bar{x}=A+c\,\frac{\sum fd}{\sum f}.
\]
\end{enumerate}
\end{methode}

\textbf{Why it works (examinable derivation).}\par
If $d_i = x_i - A$, then $x_i = A + d_i$, so
\[
\bar{x}=\frac{\sum f_i x_i}{\sum f_i}=\frac{\sum f_i (A+d_i)}{\sum f_i}
=\frac{A\sum f_i}{\sum f_i}+\frac{\sum f_i d_i}{\sum f_i}
=A+\frac{\sum f_i d_i}{\sum f_i}.
\]
With $d_i = \dfrac{x_i-A}{c}$ we have $x_i = A + c d_i$ and the same algebra gives $\bar{x}=A+c\dfrac{\sum f_i d_i}{\sum f_i}$. This derivation is short and can be asked in the examination.

\begin{exoresolu}[Coding method with step deviation]
Using the cement data above with $A=52$ and $c=3$, estimate the mean.
\tcblower
\begin{center}
\begin{tabular}{|c|c|c|c|c|c|}
\hline
\rowcolor{popblueL} Class & $x$ & $f$ & $fx$ & $d=\frac{x-52}{3}$ & $fd$ \\
\hline
$45$--$47$ & $46$ & $4$ & $184$ & $-2$ & $-8$ \\
$48$--$50$ & $49$ & $10$ & $490$ & $-1$ & $-10$ \\
$51$--$53$ & $52$ & $14$ & $728$ & $0$ & $0$ \\
$54$--$56$ & $55$ & $8$ & $440$ & $1$ & $8$ \\
$57$--$59$ & $58$ & $4$ & $232$ & $2$ & $8$ \\
\hline
Totals & & $40$ & $2074$ & & $-2$ \\
\hline
\end{tabular}
\end{center}
\[
\bar{x}=A+c\,\frac{\sum fd}{\sum f}=52+3\times\frac{-2}{40}=52-0.15=51.85\ \mathrm{kg},
\]
exactly the same answer as before, with much smaller numbers.
\end{exoresolu}

\begin{attention}[Sign errors in the coding table]
Deviations below $A$ are \emph{negative}; do not drop the signs when summing $\sum fd$. Also, if you use the step deviation $d=\dfrac{x-A}{c}$, you must multiply back by $c$ at the end --- forgetting the factor $c$ is one of the most common errors at the GCE.
\end{attention}

\begin{exoresolu}[Finding an unknown frequency]
The mean of the following distribution is $25$. Find the missing frequency $p$.
\begin{center}
\begin{tabular}{|c|cccc|}
\hline
\rowcolor{popblueL} $x$ & $10$ & $20$ & $30$ & $40$ \\
\hline
$f$ & $5$ & $p$ & $10$ & $5$ \\
\hline
\end{tabular}
\end{center}
\tcblower
$\sum f = 20+p$, $\sum fx = 50 + 20p + 300 + 200 = 550 + 20p$.
\[
\bar{x} = \frac{550+20p}{20+p} = 25
\;\Longrightarrow\;
550+20p = 500 + 25p
\;\Longrightarrow\;
5p = 50
\;\Longrightarrow\;
p = 10.
\]
\end{exoresolu}

\courssub{Effect of Adding or Multiplying by a Constant}

\begin{propriete}[Linear transformations of the mean]
Let $\bar{x}$ be the mean of $x_1,\dots,x_n$, and let $a$ and $b$ be constants.
\begin{enumerate}
\item If $y_i = x_i + b$ (a constant added to every value), then $\bar{y}=\bar{x}+b$.
\item If $y_i = a x_i$ (every value multiplied by a constant), then $\bar{y}=a\bar{x}$.
\item In general, if $y_i = a x_i + b$, then $\bar{y}=a\bar{x}+b$.
\end{enumerate}
\end{propriete}

\textbf{Proof.} For (1): $\bar{y}=\dfrac{\sum (x_i+b)}{n}=\dfrac{\sum x_i}{n}+\dfrac{nb}{n}=\bar{x}+b$. For (2): $\bar{y}=\dfrac{\sum a x_i}{n}=a\dfrac{\sum x_i}{n}=a\bar{x}$. Part (3) follows by combining the two. \hfill$\blacksquare$

\begin{exemplebox}[Using the transformation property]
The mean mark of a class in a test marked out of $20$ is $11.5$. The teacher decides to add $2$ bonus marks to every script, then the school converts all marks to a percentage (multiply by $5$). The new mean is
\[
(11.5+2)\times 5 = 67.5\ \%.
\]
\end{exemplebox}

\begin{exoresolu}[Correcting a mean after detecting an error]
The mean of $20$ observations is $35$. Later it is found that one value, $50$, was wrongly recorded as $5$. Find the corrected mean.
\tcblower
The wrong total is $\sum x_{\text{wrong}} = n\bar{x} = 20 \times 35 = 700$.
The correct total is $\sum x_{\text{correct}} = 700 - 5 + 50 = 745$.
The corrected mean is $\bar{x}_{\text{correct}} = \dfrac{745}{20} = 37.25$.
\end{exoresolu}

\courssub{Combined Mean of Two or More Groups}

Suppose Form Six Science A has $30$ students with mean mark $62$, and Form Six Science B has $20$ students with mean mark $55$. What is the mean of the two classes together? It is \emph{not} $\frac{62+55}{2}=58.5$, because the classes have different sizes. We must recover the \emph{total} marks of each group: since $\bar{x}=\frac{\sum x}{n}$, the total of a group is $n\bar{x}$.

\begin{propriete}[Combined mean of $k$ groups]
If group $i$ has $n_i$ observations and mean $\bar{x}_i$ ($i=1,\dots,k$), the mean of all the observations combined is
\[
\bar{x}_c=\frac{n_1\bar{x}_1+n_2\bar{x}_2+\cdots+n_k\bar{x}_k}{n_1+n_2+\cdots+n_k}.
\]
For two groups: $\bar{x}_c=\dfrac{n_1\bar{x}_1+n_2\bar{x}_2}{n_1+n_2}$. The formula and its derivation from totals are examinable.
\end{propriete}

\textbf{Derivation.} The total of group $i$ is $n_i\bar{x}_i$. Hence the grand total is $\sum n_i\bar{x}_i$ and the grand number of observations is $\sum n_i$, so $\bar{x}_c=\dfrac{\sum n_i\bar{x}_i}{\sum n_i}$. \hfill$\blacksquare$

\begin{methode}[Finding the combined mean]
\begin{enumerate}
\item Identify the number of observations $n_i$ and the mean $\bar{x}_i$ for each group.
\item Compute the total for each group: $T_i = n_i \bar{x}_i$.
\item Find the grand total $T = \sum T_i$ and the grand number $N = \sum n_i$.
\item The combined mean is $\bar{x}_c = T / N$.
\end{enumerate}
\end{methode}

\begin{exoresolu}[Combined mean of two classes]
Science A: $n_1=30$, $\bar{x}_1=62$. Science B: $n_2=20$, $\bar{x}_2=55$. Find the combined mean.
\tcblower
\[
\bar{x}_c=\frac{n_1\bar{x}_1+n_2\bar{x}_2}{n_1+n_2}
=\frac{30(62)+20(55)}{30+20}
=\frac{1860+1100}{50}
=\frac{2960}{50}=59.2.
\]
The combined mean $59.2$ lies between $55$ and $62$, and is closer to $62$ because Science A is the larger group.
\end{exoresolu}

\begin{exoresolu}[Finding an unknown group mean]
The mean height of $25$ boys in a class is \SI{168}{cm}. The mean height of all $40$ students (boys and girls) is \SI{165}{cm}. Find the mean height of the $15$ girls.
\tcblower
Let $\bar{x}_g$ be the girls' mean height. Using the combined-mean formula:
\[
165=\frac{25(168)+15\bar{x}_g}{40}
\;\Longrightarrow\;
6600=4200+15\bar{x}_g
\;\Longrightarrow\;
\bar{x}_g=\frac{2400}{15}=160\ \mathrm{cm}.
\]
\end{exoresolu}

\begin{exoresolu}[Combined mean of three groups]
A company has three branches. Branch A: $50$ staff, mean salary \SI{300000}{FCFA}. Branch B: $30$ staff, mean salary \SI{250000}{FCFA}. Branch C: $20$ staff, mean salary \SI{400000}{FCFA}. Find the overall mean salary.
\tcblower
\[
\bar{x}_c = \frac{50(300000) + 30(250000) + 20(400000)}{50+30+20}
= \frac{15000000 + 7500000 + 8000000}{100}
= \frac{30500000}{100}
= 305000\ \mathrm{FCFA}.
\]
\end{exoresolu}

\begin{attention}[Averaging the averages]
Never compute $\dfrac{\bar{x}_1+\bar{x}_2}{2}$ when the group sizes differ. The combined mean is a \emph{weighted} average: each group mean must be multiplied by its group size. The unweighted average is only correct when $n_1=n_2$.
\end{attention}

\courssub{Weighted Mean}

The combined mean is a special case of a more general idea: the \cle{weighted mean}, used whenever some values ``count more'' than others.

\begin{definition}[Weighted mean]
If values $x_1, x_2, \dots, x_k$ carry weights $w_1, w_2, \dots, w_k$, the weighted mean is
\[
\bar{x}_w=\frac{\sum w_i x_i}{\sum w_i}.
\]
\end{definition}

The link with the combined mean is immediate: the combined mean is the weighted mean of the group means, with the group sizes $n_i$ as weights. Similarly, the mean of a frequency table $\bar{x}=\dfrac{\sum fx}{\sum f}$ is a weighted mean with the frequencies as weights.

\begin{exemplebox}[Weighted mean: a student's average]
In a school, the term mark counts for $30\%$ and the examination mark for $70\%$. A student scores $14/20$ in the term work and $11/20$ in the examination. Her weighted average is
\[
\bar{x}_w=\frac{0.3(14)+0.7(11)}{0.3+0.7}=\frac{4.2+7.7}{1}=11.9/20.
\]
A plain average would give $12.5/20$, overvaluing the term mark.
\end{exemplebox}

\begin{exemplebox}[Weighted mean: price index preparation]
A household buys three commodities. Their prices (FCFA/kg) and quantities bought (kg) are:
\begin{center}
\begin{tabular}{|c|ccc|}
\hline
\rowcolor{popblueL} Commodity & Price $p$ & Quantity $q$ & $pq$ \\
\hline
Rice & $600$ & $50$ & $30000$ \\
Beans & $800$ & $20$ & $16000$ \\
Maize & $400$ & $30$ & $12000$ \\
\hline
Total & & $100$ & $58000$ \\
\hline
\end{tabular}
\end{center}
The weighted mean price per kg (using quantities as weights) is $\frac{58000}{100}=580$ FCFA/kg. This is the basis for Laspeyres and Paasche index numbers studied later.
\end{exemplebox}

\courssub{Mean, Median or Mode? Choosing the Right Average}

The mean uses every value, which is its strength --- and its weakness. A single extreme value can drag the mean far from the bulk of the data.

\begin{exemplebox}[Salaries with an extreme value]
A small company in Douala pays its nine workers monthly salaries (in thousands of FCFA):
\[
80,\ 85,\ 90,\ 90,\ 95,\ 100,\ 100,\ 105,\ 800
\]
where the last figure is the director's salary. The mean is
\[
\bar{x}=\frac{1545}{9}\approx 171.7 \text{ thousand FCFA},
\]
yet \emph{eight of the nine} workers earn less than $110\,000$ FCFA. The median, $95\,000$ FCFA, describes a ``typical'' worker far better. If the company advertises ``average salary $171\,700$ FCFA'', it is technically true but misleading.
\end{exemplebox}

\begin{aretenir}[Which average to use?]
\begin{itemize}
\item \textbf{Mean}: use when the data are roughly symmetric with no extreme values, and when further calculations (variance, combined means, regression) are needed. It is the only average that can be treated algebraically.
\item \textbf{Median}: use when the data contain extreme values or are skewed (salaries, house prices, incomes, reaction times). It is robust and represents the ``typical'' value.
\item \textbf{Mode}: use for categorical data (favourite colour, most popular shoe size) or when the most popular value matters (stocking inventory). For grouped data, the modal class is the class with highest frequency.
\end{itemize}
\end{aretenir}

\begin{savaistu}[The mean in national debates]
When newspapers report the ``average income'' of a country, they usually quote the mean, which is pulled upwards by a small number of very high earners. This is why the \emph{median} income is often much lower, and why economists prefer it when describing the living standard of a typical household. As a statistician, always ask \emph{which} average is being quoted!
\end{savaistu}

\begin{exercice}[Practice]
\begin{enumerate}
\item The mean of the numbers $4, 7, 9, x, 15$ is $9$. Find $x$.
\item A grouped distribution has $\sum f = 50$ and, with assumed mean $A=27$ and class width $c=5$, $\sum fd = 12$. Find the mean.
\item Group X of $40$ candidates has mean mark $58$; group Y of $60$ candidates has mean mark $64$. Find the combined mean.
\item The mean of $20$ observations is $35$. Later it is found that one value, $50$, was wrongly recorded as $5$. Find the corrected mean.
\item The mean of $10$ numbers is $12$. If one number is removed, the mean of the remaining $9$ numbers is $11$. Find the removed number.
\item A student's marks in five subjects are $65, 72, 58, 80, 60$. The weights (coefficients) for these subjects are $3, 4, 2, 5, 1$ respectively. Calculate the weighted mean.
\item The mean of a grouped frequency distribution with $5$ classes is $41$. The class mid-points are $30, 35, 40, 45, 50$ and the corresponding frequencies are $4, 8, p, 10, 6$. Find $p$.
\end{enumerate}
\end{exercice}

\begin{corrige}[Exercise 1]
\begin{enumerate}
\item $\dfrac{4+7+9+x+15}{5}=9 \Rightarrow 35+x=45 \Rightarrow x=10$.
\item $\bar{x}=A+c\dfrac{\sum fd}{\sum f}=27+5\times\dfrac{12}{50}=27+1.2=28.2$.
\item $\bar{x}_c=\dfrac{40(58)+60(64)}{100}=\dfrac{2320+3840}{100}=61.6$.
\item The wrong total is $20\times 35=700$. Corrected total: $700-5+50=745$. Corrected mean: $\dfrac{745}{20}=37.25$.
\item Total of $10$ numbers $= 10 \times 12 = 120$. Total of remaining $9$ numbers $= 9 \times 11 = 99$. Removed number $= 120 - 99 = 21$.
\item $\bar{x}_w = \frac{3(65)+4(72)+2(58)+5(80)+1(60)}{3+4+2+5+1} = \frac{195+288+116+400+60}{15} = \frac{1059}{15} = 70.6$.
\item $\sum f = 28+p$, $\sum fx = 120+280+40p+450+300 = 1150+40p$.
$\frac{1150+40p}{28+p} = 41 \Rightarrow 1150+40p = 1148+41p \Rightarrow p=2$.
\end{enumerate}
\end{corrige}

\begin{aretenir}[Key points of this section]
\begin{itemize}
\item Ungrouped data: $\bar{x}=\dfrac{\sum x}{n}$; frequency table: $\bar{x}=\dfrac{\sum fx}{\sum f}$.
\item Grouped data: use class mid-points $x = \frac{L+U}{2}$; the mean is an \emph{estimate}.
\item Coding method: $\bar{x}=A+\frac{\sum fd}{\sum f}$ or $\bar{x}=A+c\frac{\sum fd}{\sum f}$ with $d=\frac{x-A}{c}$. Derivation is examinable.
\item If $y=ax+b$, then $\bar{y}=a\bar{x}+b$.
\item Combined mean: $\bar{x}_c=\dfrac{n_1\bar{x}_1+n_2\bar{x}_2}{n_1+n_2}$ --- always weight by group sizes.
\item Weighted mean: $\bar{x}_w=\dfrac{\sum w_i x_i}{\sum w_i}$.
\item The mean is sensitive to extreme values; prefer the median for skewed data such as salaries.
\end{itemize}
\end{aretenir}

\coursec{Measures of Dispersion: Variance and Standard Deviation}

In the previous section we studied the \cle{mean}, and we even saw how to combine the means of several groups into a single \emph{combined mean}. The mean tells us where the \emph{centre} of a data set lies. But two data sets can share exactly the same mean and yet behave very differently. Imagine two bakeries in Douala, each producing an average of $500$ loaves per day. Bakery A produces $498, 500, 502, 499, 501$ over five days; Bakery B produces $200, 750, 480, 620, 450$. Both have mean $500$, but Bakery A is far more \emph{consistent}. To distinguish such situations we need \cle{measures of dispersion} — numbers that describe how \emph{spread out} the data are around the centre.

\courssub{Why the Range Is Not Enough}

The simplest measure of spread is the \cle{range}: $\text{range} = x_{\max} - x_{\min}$. It is quick, but it uses only \emph{two} values out of the whole data set and ignores everything in between.

\begin{exemplebox}[The weakness of the range]
Consider the two sets
\[ A:\ 2,\ 50,\ 50,\ 50,\ 98 \qquad B:\ 2,\ 10,\ 40,\ 70,\ 98. \]
Both have range $98 - 2 = 96$ and mean $50$. Yet in $A$ most values sit right at $50$, while in $B$ the values are scattered across the whole interval. The range cannot see this difference.
\end{exemplebox}

\begin{attention}[The range is fragile]
Because the range depends only on the two extreme values, a single \emph{outlier} (one unusually large or small value) can inflate it enormously, even when all the other data are tightly clustered. Never use the range alone to compare the spread of two data sets.
\end{attention}

A natural first improvement is the \cle{mean deviation}: take the distance of each value from the mean, $|x - \bar{x}|$, and average these distances:
\[ \text{Mean deviation} = \frac{1}{n}\sum |x - \bar{x}|. \]
This uses \emph{every} data value. However, the absolute value is awkward to handle algebraically (it is not differentiable at $0$, and it does not combine nicely with other formulas). Statisticians therefore prefer to \emph{square} the deviations instead of taking absolute values. Squaring also removes the sign, and it penalises large deviations more heavily — which is usually desirable. This leads us to the central notions of this section.

\courssub{Variance and Standard Deviation of Ungrouped Data}

\begin{definition}[Variance and standard deviation]
Let $x_1, x_2, \dots, x_n$ be $n$ observations with mean $\bar{x}$. The \cle{variance}, denoted $\sigma^2$ (or $\mathrm{Var}(x)$), is the mean of the squared deviations from the mean:
\[ \sigma^2 = \frac{1}{n}\sum_{i=1}^{n}(x_i - \bar{x})^2. \]
The \cle{standard deviation} $\sigma$ is the positive square root of the variance:
\[ \sigma = \sqrt{\frac{1}{n}\sum_{i=1}^{n}(x_i - \bar{x})^2}. \]
The standard deviation is expressed in the \emph{same units} as the data (the variance is in \emph{squared} units), which makes $\sigma$ the preferred measure for interpretation.
\end{definition}

\begin{attention}[Why divide by $n$ and not $n-1$?]
At GCE Advanced Level we always treat the data as the \emph{whole population} under study, so we divide by $n$. (In more advanced work, when the data are only a \emph{sample} used to estimate a population variance, one divides by $n-1$ instead; this is \emph{not} required here.)
\end{attention}

\begin{propriete}[Computational form of the variance]
The variance can be computed without subtracting the mean from each value:
\[ \sigma^2 = \frac{\sum x^2}{n} - \bar{x}^2 \qquad\text{("mean of squares minus square of the mean").} \]
\textbf{Proof (examinable).} Expanding the square and using $\sum x = n\bar{x}$:
\begin{align*}
\sum (x - \bar{x})^2 &= \sum x^2 - 2\bar{x}\sum x + \sum \bar{x}^2 \\
&= \sum x^2 - 2\bar{x}(n\bar{x}) + n\bar{x}^2 \\
&= \sum x^2 - n\bar{x}^2.
\end{align*}
Dividing by $n$ gives $\sigma^2 = \dfrac{\sum x^2}{n} - \bar{x}^2$. $\blacksquare$
\end{propriete}

\begin{exoresolu}[Variance of ungrouped data]
The marks of $6$ candidates in a quiz are $4,\ 6,\ 7,\ 8,\ 9,\ 8$. Find the variance and the standard deviation.
\tcblower
\textbf{Solution.} First the mean:
\[ \bar{x} = \frac{4+6+7+8+9+8}{6} = \frac{42}{6} = 7. \]
Using the computational form, $\sum x^2 = 16 + 36 + 49 + 64 + 81 + 64 = 310$, so
\[ \sigma^2 = \frac{310}{6} - 7^2 = 51.666\ldots - 49 = 2.666\ldots \approx 2.67, \]
\[ \sigma = \sqrt{2.666\ldots} \approx 1.63 \text{ marks}. \]
\emph{Check by the definition:} the deviations are $-3, -1, 0, 1, 2, 1$; their squares are $9, 1, 0, 1, 4, 1$, summing to $16$; and $\frac{16}{6} = 2.666\ldots$. Same answer. $\blacksquare$
\end{exoresolu}

\courssub{Variance and Standard Deviation of Grouped Data}

When data are presented in a frequency table, each value $x$ (or class mid-point) occurs $f$ times, so each squared deviation must be weighted by its frequency.

\begin{propriete}[Variance of grouped data]
For a frequency distribution with values (or mid-points) $x_i$ and frequencies $f_i$,
\[ \sigma^2 = \frac{\sum f(x - \bar{x})^2}{\sum f} = \frac{\sum f x^2}{\sum f} - \bar{x}^2, \qquad \bar{x} = \frac{\sum fx}{\sum f}. \]
The standard deviation is $\sigma = \sqrt{\sigma^2}$.
\end{propriete}

\begin{methode}[Computing $\sigma^2$ from a grouped table]
\begin{enumerate}
\item Find the mid-point $x$ of each class.
\item Compute the mean $\bar{x} = \dfrac{\sum fx}{\sum f}$ (add a column $fx$).
\item Add a column $fx^2$ (compute $x^2$ then multiply by $f$, or compute $(fx)\times x$).
\item Apply $\sigma^2 = \dfrac{\sum fx^2}{\sum f} - \bar{x}^2$, then take the square root.
\end{enumerate}
\end{methode}

\begin{exoresolu}[Variance of grouped data]
The table shows the masses (kg) of $50$ bags of cement sampled at a depot in Yaoundé. Find the variance and standard deviation.
\begin{center}
\begin{tabular}{|c|c|c|c|c|}
\hline
\rowcolor{popblueL} Class & Mid-point $x$ & $f$ & $fx$ & $fx^2$ \\ \hline
$45$--$47$ & $46$ & $6$ & $276$ & $12\,696$ \\ \hline
$48$--$50$ & $49$ & $18$ & $882$ & $43\,218$ \\ \hline
$51$--$53$ & $52$ & $16$ & $832$ & $43\,264$ \\ \hline
$54$--$56$ & $55$ & $8$ & $440$ & $24\,200$ \\ \hline
$57$--$59$ & $58$ & $2$ & $116$ & $6\,728$ \\ \hline
\rowcolor{popblueL} Totals & & $\sum f = 50$ & $\sum fx = 2546$ & $\sum fx^2 = 130\,106$ \\ \hline
\end{tabular}
\end{center}
\tcblower
\textbf{Solution.}
\[ \bar{x} = \frac{2546}{50} = 50.92 \text{ kg}, \qquad
\sigma^2 = \frac{130\,106}{50} - (50.92)^2 = 2602.12 - 2592.8464 = 9.2736. \]
\[ \sigma = \sqrt{9.2736} \approx 3.05 \text{ kg}. \]
On average, a bag's mass deviates from the mean by roughly $3$ kg.
\end{exoresolu}

\begin{popfigure}
\begin{center}
\begin{tikzpicture}
\begin{axis}[
  width=13cm, height=6.5cm,
  ybar, bar width=7pt,
  ymin=0, ymax=30,
  xlabel={Daily output (units)}, ylabel={Frequency},
  xtick={10,20,30,40,50,60,70,80,90},
  legend style={at={(0.5,1.02)}, anchor=south, legend columns=2, draw=none},
  axis lines=left, every axis/.append style={font=\small}
]
\addplot+[ybar, fill=popblue, draw=popdark] coordinates {(20,2) (30,8) (40,18) (50,26) (60,18) (70,8) (80,2)};
\addlegendentry{Factory A (small spread)}
\addplot+[ybar, fill=poporange, draw=popdark] coordinates {(10,4) (20,10) (30,12) (40,14) (50,15) (60,14) (70,12) (80,10) (90,4)};
\addlegendentry{Factory B (large spread)}
\end{axis}
\end{tikzpicture}
\end{center}
\legende{Two factories with the same mean output ($50$ units) but very different spreads: Factory A has a small standard deviation (consistent), Factory B a large one (erratic).}
\end{popfigure}

\courssub{The Coding Method for Variance}

When the values $x$ are large or the mean is not a whole number, direct computation is heavy. We simplify the arithmetic by a change of variable (\cle{coding}):
\[ d = \frac{x - A}{c} \quad\Longleftrightarrow\quad x = A + cd, \]
where $A$ is a chosen \emph{assumed mean} (usually a central mid-point) and $c$ is the class width (or any convenient common divisor).

\begin{propriete}[Coding formula for the variance]
With $d = \dfrac{x-A}{c}$,
\[ \sigma^2 = c^2\left[\frac{\sum f d^2}{\sum f} - \left(\frac{\sum f d}{\sum f}\right)^2\right]. \]
\textbf{Justification.} Since $x = A + cd$, we have $\bar{x} = A + c\bar{d}$ and $x - \bar{x} = c(d - \bar{d})$. Hence
\[ \sum f(x-\bar{x})^2 = c^2 \sum f(d - \bar{d})^2, \]
and applying the computational form to the $d$'s gives the result. $\blacksquare$
\end{propriete}

\begin{methode}[Coding technique for variance]
\begin{enumerate}
\item Choose an assumed mean $A$ near the centre and the step $c$ (class width).
\item Form the column $d = (x-A)/c$; the $d$'s should be small integers.
\item Compute columns $fd$ and $fd^2$ and their totals.
\item Compute $\sigma^2 = c^2\left[\dfrac{\sum fd^2}{\sum f} - \left(\dfrac{\sum fd}{\sum f}\right)^2\right]$.
\item Take the square root for $\sigma$.
\end{enumerate}
\end{methode}

\begin{exoresolu}[Variance by coding]
Using the cement-bag data above with $A = 52$ and $c = 3$:
\begin{center}
\begin{tabular}{|c|c|c|c|c|}
\hline
\rowcolor{popblueL} $x$ & $f$ & $d=\frac{x-52}{3}$ & $fd$ & $fd^2$ \\ \hline
$46$ & $6$ & $-2$ & $-12$ & $24$ \\ \hline
$49$ & $18$ & $-1$ & $-18$ & $18$ \\ \hline
$52$ & $16$ & $0$ & $0$ & $0$ \\ \hline
$55$ & $8$ & $1$ & $8$ & $8$ \\ \hline
$58$ & $2$ & $2$ & $4$ & $8$ \\ \hline
\rowcolor{popblueL} Totals & $50$ & & $-18$ & $58$ \\ \hline
\end{tabular}
\end{center}
Verify the value of $\sigma^2$ found earlier.
\tcblower
\textbf{Solution.}
\[ \sigma^2 = 3^2\left[\frac{58}{50} - \left(\frac{-18}{50}\right)^2\right] = 9\left[1.16 - 0.1296\right] = 9(1.0304) = 9.2736. \]
Hence $\sigma \approx 3.05$ kg, exactly as before — but with far simpler arithmetic. Note also that $\bar{x} = A + c\bar{d} = 52 + 3\times\frac{-18}{50} = 52 - 1.08 = 50.92$ kg, confirming the mean.
\end{exoresolu}

\begin{attention}[Do not forget the factor $c^2$!]
The most common examination error with the coding method is to compute the variance \emph{of the coded values} $d$ and to forget to multiply by $c^2$ at the end. Remember: the $d$'s are in "step units"; each step is worth $c$ real units, and since variance involves \emph{squares}, the correction factor is $c^2$ (not $c$). For the standard deviation, multiply by $c$ only: $\sigma_x = c\,\sigma_d$.
\end{attention}

\courssub{Effect of Linear Transformations}

The coding formula is a consequence of a general and very useful result.

\begin{propriete}[Variance under a linear transformation]
Let $y = ax + b$ where $a$ and $b$ are constants. Then
\[ \mathrm{Var}(y) = a^2\,\mathrm{Var}(x), \qquad \sigma_y = |a|\,\sigma_x. \]
In particular:
\begin{itemize}
\item \textbf{Adding a constant does not change the variance:} $\mathrm{Var}(x + b) = \mathrm{Var}(x)$. Shifting all the data moves the mean by the same amount, so the deviations $x - \bar{x}$ are unchanged.
\item \textbf{Multiplying by $a$ multiplies the variance by $a^2$:} $\mathrm{Var}(ax) = a^2\mathrm{Var}(x)$.
\end{itemize}
\textbf{Proof.} If $y = ax+b$ then $\bar{y} = a\bar{x} + b$, so $y - \bar{y} = a(x - \bar{x})$ and
\[ \mathrm{Var}(y) = \frac{1}{n}\sum a^2(x-\bar{x})^2 = a^2\,\mathrm{Var}(x). \quad\blacksquare \]
\end{propriete}

\begin{exemplebox}[Using the transformation property]
The marks of a class have mean $52$ and variance $16$. The teacher decides to scale every mark by $y = 1.5x - 10$. Then
\[ \bar{y} = 1.5(52) - 10 = 68, \qquad \mathrm{Var}(y) = (1.5)^2(16) = 36, \qquad \sigma_y = 6. \]
Note that the shift of $-10$ has \emph{no effect} on the variance.
\end{exemplebox}

\courssub{Interpretation: Comparing Consistency}

The standard deviation measures the \emph{typical distance} of the data from the mean. A small $\sigma$ means the values cluster tightly around $\bar{x}$ (consistent, reliable, homogeneous); a large $\sigma$ means wide scatter (erratic, heterogeneous).

\begin{exemplebox}[Two factories]
Two factories each produce an average of $500$ bottles per hour. Factory A has $\sigma = 5$; Factory B has $\sigma = 40$. Factory A's hourly output almost always lies within a few bottles of $500$, whereas Factory B's output swings wildly. A manager would judge Factory A far more \emph{consistent}, even though both have the same mean. The mean alone would have hidden this completely.
\end{exemplebox}

\begin{savaistu}[Why statisticians love the standard deviation]
For many roughly bell-shaped distributions, about $68\%$ of the data lie within one standard deviation of the mean and about $95\%$ within two. This "empirical rule" makes $\sigma$ a powerful yardstick: a value more than $2\sigma$ from the mean is already unusual.
\end{savaistu}

\courssub{Coefficient of Variation}

Suppose we want to compare the dispersion of the heights of students (measured in cm, mean around $165$) with the dispersion of their masses (in kg, mean around $60$). The standard deviations are in \emph{different units} and refer to \emph{different-sized means}, so comparing them directly is meaningless. We therefore use a \emph{relative}, unit-free measure.

\begin{definition}[Coefficient of variation]
The \cle{coefficient of variation} is the standard deviation expressed as a percentage of the mean:
\[ \mathrm{CV} = \frac{\sigma}{\bar{x}} \times 100\%. \]
The smaller the CV, the more consistent the data set \emph{relative to its mean}.
\end{definition}

\begin{exoresolu}[Comparing data with different units]
The heights of a group of students have $\bar{x} = 165$ cm, $\sigma = 8.25$ cm; their masses have $\bar{x} = 60$ kg, $\sigma = 7.2$ kg. Which characteristic is more variable?
\tcblower
\textbf{Solution.}
\[ \mathrm{CV}_{\text{height}} = \frac{8.25}{165}\times 100\% = 5\%, \qquad
\mathrm{CV}_{\text{mass}} = \frac{7.2}{60}\times 100\% = 12\%. \]
Although the masses have a smaller standard deviation in absolute terms, they are \emph{relatively} much more variable ($12\%$ against $5\%$). The CV, being unit-free, allows this fair comparison.
\end{exoresolu}

\begin{attention}[Trap: CV requires a meaningful zero]
The coefficient of variation is only sensible for quantities measured on a scale with a genuine zero (masses, prices, counts). It is meaningless for temperatures in ${}^\circ$C, since $\bar{x}$ could be near $0$ or even negative, making the CV explode or change sign.
\end{attention}

\begin{aretenir}[Key points of this section]
\begin{itemize}
\item The \textbf{range} uses only two values and is distorted by outliers; dispersion must involve \emph{all} the data.
\item Variance: $\sigma^2 = \dfrac{\sum (x-\bar{x})^2}{n} = \dfrac{\sum x^2}{n} - \bar{x}^2$ (ungrouped); $\sigma^2 = \dfrac{\sum fx^2}{\sum f} - \bar{x}^2$ (grouped).
\item Standard deviation $\sigma = \sqrt{\sigma^2}$ is in the same units as the data.
\item Coding with $d = \dfrac{x-A}{c}$: $\sigma^2 = c^2\left[\dfrac{\sum fd^2}{\sum f} - \left(\dfrac{\sum fd}{\sum f}\right)^2\right]$ — \textbf{never forget the $c^2$}.
\item $\mathrm{Var}(ax+b) = a^2\mathrm{Var}(x)$: adding a constant changes nothing; multiplying by $a$ multiplies the variance by $a^2$.
\item Small $\sigma$ $\Rightarrow$ consistent data; large $\sigma$ $\Rightarrow$ scattered data.
\item Coefficient of variation $\mathrm{CV} = \dfrac{\sigma}{\bar{x}}\times 100\%$ compares dispersion across different units or different-sized means.
\end{itemize}
\end{aretenir}

\begin{exercice}[Practice]
\begin{enumerate}
\item Find the variance and standard deviation of $3,\ 5,\ 7,\ 9,\ 11$.
\item A frequency distribution has $\sum f = 40$, $\sum fx = 200$, $\sum fx^2 = 1240$. Find the mean, variance and standard deviation.
\item Using $A = 25$ and $c = 5$, the coded values of a distribution give $\sum f = 60$, $\sum fd = 12$, $\sum fd^2 = 96$. Find $\bar{x}$ and $\sigma$.
\item The marks in Mathematics have mean $55$ and standard deviation $10$; the marks in Physics have mean $40$ and standard deviation $8$. Using the coefficient of variation, determine in which subject the class is more consistent.
\end{enumerate}
\end{exercice}

\begin{corrige}[Exercise 1]
\begin{enumerate}
\item $\bar{x} = 7$; $\sum x^2 = 9+25+49+81+121 = 285$; $\sigma^2 = \frac{285}{5} - 49 = 57 - 49 = 8$; $\sigma = \sqrt{8} \approx 2.83$.
\item $\bar{x} = \frac{200}{40} = 5$; $\sigma^2 = \frac{1240}{40} - 25 = 31 - 25 = 6$; $\sigma = \sqrt{6} \approx 2.45$.
\item $\bar{d} = \frac{12}{60} = 0.2$, so $\bar{x} = A + c\bar{d} = 25 + 5(0.2) = 26$. $\sigma^2 = 25\left[\frac{96}{60} - 0.04\right] = 25(1.56) = 39$, hence $\sigma = \sqrt{39} \approx 6.24$.
\item $\mathrm{CV}_{\text{Maths}} = \frac{10}{55}\times 100\% \approx 18.2\%$; $\mathrm{CV}_{\text{Physics}} = \frac{8}{40}\times 100\% = 20\%$. The class is slightly more consistent in Mathematics.
\end{enumerate}
\end{corrige}

\coursec{Quartiles, Deciles, Percentiles and Interquartile Range}

In the previous section we measured dispersion with the \cle{variance} and the \cle{standard deviation}. Those measures are powerful, but they have two weaknesses: they depend on \emph{every} value in the data set, so a single extreme value can inflate them enormously, and they require heavy computation. There is another family of tools, based purely on \emph{position} in ordered data, which is simpler, more robust, and extremely popular at the GCE: the \cle{quartiles}, \cle{deciles} and \cle{percentiles}. Just as the median splits ordered data into two equal halves, quartiles split it into four equal parts, deciles into ten, and percentiles into a hundred. From these positions we build new measures of spread --- the \cle{interquartile range} and the \cle{percentile range} --- and a beautiful summary picture, the \cle{box-and-whisker plot}.

\courssub{Quartiles of Ungrouped Data}

\textbf{Intuition.}\par
Imagine the 11 best students of an Upper Sixth class in Bamenda ranked by their mock examination mark. The median is the mark of the 6th student: five students are below, five above. Now split each half again: the mark of the 3rd student leaves a quarter of the group below it, and the mark of the 9th student leaves three quarters below it. These two marks are the \emph{lower} and \emph{upper} quartiles.

\begin{definition}[Quartiles, deciles and percentiles]
Let a data set of $N$ values be arranged in \textbf{ascending order}.
\begin{itemize}
\item The \cle{quartiles} $Q_1$, $Q_2$, $Q_3$ are the values that divide the data into \textbf{four equal parts}. $Q_1$ (lower quartile) has about $25\%$ of the data below it, $Q_2$ is the median ($50\%$), and $Q_3$ (upper quartile) has about $75\%$ below it.
\item The \cle{deciles} $D_1, D_2, \dots, D_9$ divide the data into \textbf{ten equal parts}: $D_k$ has about $\tfrac{k}{10}$ of the data below it.
\item The \cle{percentiles} $P_1, P_2, \dots, P_{99}$ divide the data into \textbf{one hundred equal parts}: $P_k$ has about $k\%$ of the data below it.
\end{itemize}
Note that $Q_2 = D_5 = P_{50} = \text{median}$, and $Q_1 = P_{25}$, $Q_3 = P_{75}$.
\end{definition}

\begin{methode}[Quartiles of ungrouped data (Tukey's hinges method)]
\begin{enumerate}
\item Arrange the $N$ values in ascending order.
\item Locate the median $Q_2$:
      \begin{itemize}
      \item If $N$ is odd, $Q_2$ is the middle value (position $\frac{N+1}{2}$).
      \item If $N$ is even, $Q_2$ is the mean of the two middle values (positions $\frac{N}{2}$ and $\frac{N}{2}+1$).
      \end{itemize}
\item Split the ordered data into a \textbf{lower half} and an \textbf{upper half}:
      \begin{itemize}
      \item If $N$ is odd, \emph{exclude} the median itself; the lower half consists of the values strictly below $Q_2$, the upper half of those strictly above.
      \item If $N$ is even, the lower half is the first $N/2$ values, the upper half the last $N/2$ values.
      \end{itemize}
\item $Q_1$ is the median of the lower half; $Q_3$ is the median of the upper half (apply the same rule as in step 2 to each half).
\end{enumerate}
\end{methode}

\begin{exemplebox}[Quartiles of 11 marks (odd $N$)]
The marks of 11 students are, in order:
\[ 4,\ 7,\ 8,\ 10,\ 12,\ 15,\ 18,\ 20,\ 22,\ 25,\ 30. \]
$N=11$ (odd). Median position $= \frac{11+1}{2} = 6$; $Q_2 = 15$ (the 6th value).
Lower half (5 values): $4, 7, 8, 10, 12$; its median is the 3rd value $\Rightarrow Q_1 = 8$.
Upper half (5 values): $18, 20, 22, 25, 30$; its median is the 3rd value $\Rightarrow Q_3 = 22$.
\end{exemplebox}

\begin{exemplebox}[Quartiles of 10 marks (even $N$)]
Marks: $5,\ 8,\ 12,\ 14,\ 16,\ 18,\ 20,\ 22,\ 25,\ 30$.
$N=10$ (even). Median $Q_2 = \frac{16+18}{2} = 17$ (average of 5th and 6th).
Lower half (first 5): $5, 8, 12, 14, 16$; median $Q_1 = 12$ (3rd value).
Upper half (last 5): $18, 20, 22, 25, 30$; median $Q_3 = 22$ (3rd value).
\end{exemplebox}

\courssub{Quartiles of Grouped Data by Interpolation}

For grouped data we cannot read individual values, so we proceed exactly as for the median: locate the \emph{position} of the quartile, find the class containing it using cumulative frequencies, then estimate the value by \cle{linear interpolation}.

\begin{propriete}[Interpolation formula for quartiles of grouped data]
For grouped data with total frequency $N$:
\[ Q_i = L + \left(\frac{\dfrac{iN}{4} - F}{f}\right)\times c, \qquad i = 1, 2, 3, \]
where
\begin{itemize}
\item $L$ = \textbf{lower class boundary} of the quartile class (the first class whose cumulative frequency reaches or exceeds $\tfrac{iN}{4}$),
\item $F$ = cumulative frequency \textbf{before} that class,
\item $f$ = frequency of the quartile class,
\item $c$ = class width.
\end{itemize}
The same pattern gives deciles and percentiles:
\[ D_k = L + \left(\frac{\dfrac{kN}{10} - F}{f}\right)c, \qquad
   P_k = L + \left(\frac{\dfrac{kN}{100} - F}{f}\right)c. \]
\end{propriete}

\begin{attention}[Position is not value!]
The most frequent error at the GCE is to confuse the \textbf{position} $\tfrac{N}{4}$ with the \textbf{value} of $Q_1$. If $N = 80$, then $\tfrac{N}{4} = 20$ is \emph{not} the lower quartile: it only tells you that $Q_1$ is the value of the ``20th observation'', which you must then find inside its class by interpolation. Always finish the interpolation before writing your answer.
\end{attention}

\begin{exoresolu}[Quartiles, a decile and a percentile of grouped data]
The marks of 80 candidates in a statistics test are summarised below.
\begin{center}
\begin{tabular}{|c|c|c|}
\hline
\rowcolor{popblueL} Marks & Frequency $f$ & Cumulative frequency \\
\hline
$0$--$10$ & $8$ & $8$ \\
\hline
$10$--$20$ & $16$ & $24$ \\
\hline
$20$--$30$ & $30$ & $54$ \\
\hline
$30$--$40$ & $18$ & $72$ \\
\hline
$40$--$50$ & $8$ & $80$ \\
\hline
\end{tabular}
\end{center}
Find $Q_1$, $Q_2$, $Q_3$, $D_7$ and $P_{85}$.
\tcblower
Here $N = 80$ and every class has width $c = 10$. Class boundaries are $0, 10, 20, 30, 40, 50$.

\textbf{Lower quartile $Q_1$.} Position $\tfrac{N}{4} = 20$. The first cumulative frequency reaching 20 is 24, so the quartile class is $10$--$20$, with $L = 10$, $F = 8$, $f = 16$:
\[ Q_1 = 10 + \frac{20 - 8}{16}\times 10 = 10 + 7.5 = 17.5. \]

\textbf{Median $Q_2$.} Position $\tfrac{N}{2} = 40$; class $20$--$30$, $L = 20$, $F = 24$, $f = 30$:
\[ Q_2 = 20 + \frac{40 - 24}{30}\times 10 = 20 + 5.333\ldots \approx 25.3. \]

\textbf{Upper quartile $Q_3$.} Position $\tfrac{3N}{4} = 60$; class $30$--$40$, $L = 30$, $F = 54$, $f = 18$:
\[ Q_3 = 30 + \frac{60 - 54}{18}\times 10 = 30 + 3.333\ldots \approx 33.3. \]

\textbf{Seventh decile $D_7$.} Position $\tfrac{7N}{10} = 56$; class $30$--$40$:
\[ D_7 = 30 + \frac{56 - 54}{18}\times 10 = 30 + 1.111\ldots \approx 31.1. \]

\textbf{85th percentile $P_{85}$.} Position $\tfrac{85N}{100} = 68$; class $30$--$40$:
\[ P_{85} = 30 + \frac{68 - 54}{18}\times 10 = 30 + 7.777\ldots \approx 37.8. \]
\end{exoresolu}

\courssub{Graphical Determination from the Ogive}

Because the ogive (cumulative frequency curve) is precisely a graph of ``position against value'', quartiles and percentiles can be read directly from it: locate $\tfrac{N}{4}$, $\tfrac{N}{2}$, $\tfrac{3N}{4}$ on the cumulative frequency axis, project horizontally to the curve, then vertically down to the value axis.

\begin{popfigure}
\begin{center}
\begin{tikzpicture}
\begin{axis}[
width=12cm, height=8cm,
xlabel={Marks}, ylabel={Cumulative frequency},
xmin=0, xmax=52, ymin=0, ymax=85,
xtick={0,10,20,30,40,50}, ytick={0,20,40,60,80},
grid=both, grid style={popline, dashed, very thin},
axis line style={popdark},
]
\addplot[very thick, popblue, mark=*, mark size=1.5pt] coordinates {(0,0) (10,8) (20,24) (30,54) (40,72) (50,80)};
\addplot[dashed, poporange, thick] coordinates {(0,20) (17.5,20) (17.5,0)};
\addplot[dashed, poppurple, thick] coordinates {(0,40) (25.3,40) (25.3,0)};
\addplot[dashed, popgreen, thick] coordinates {(0,60) (33.3,60) (33.3,0)};
\addplot[popdark, thick, <->] coordinates {(17.5,-6) (33.3,-6)};
\node[poporange, below] at (axis cs:17.5,0) {$Q_1$};
\node[poppurple, below] at (axis cs:25.3,0) {$Q_2$};
\node[popgreen, below] at (axis cs:33.3,0) {$Q_3$};
\node[popdark, below] at (axis cs:25.4,-9) {IQR};
\node[popmuted, left, align=center] at (axis cs:2,20) {$\frac{N}{4}$};
\node[popmuted, left, align=center] at (axis cs:2,40) {$\frac{N}{2}$};
\node[popmuted, left, align=center] at (axis cs:2,60) {$\frac{3N}{4}$};
\end{axis}
\end{tikzpicture}
\end{center}
\legende{Reading $Q_1$, $Q_2$, $Q_3$ and the interquartile range from the ogive ($N = 80$).}
\end{popfigure}

\courssub{Interquartile Range and Semi-Interquartile Range}

The quartiles immediately give a measure of dispersion focused on the \emph{middle half} of the data, completely unaffected by extreme values.

\begin{definition}[Interquartile and semi-interquartile range]
The \cle{interquartile range} (IQR) is the difference between the upper and lower quartiles:
\[ \mathrm{IQR} = Q_3 - Q_1. \]
The \cle{semi-interquartile range} (also called the \emph{quartile deviation}) is half of it:
\[ \mathrm{SIQR} = \frac{Q_3 - Q_1}{2}. \]
The IQR contains the middle $50\%$ of the observations. A small IQR means the central half of the data is tightly packed.
\end{definition}

\begin{exemplebox}[Computing the IQR]
For our 80 candidates: $\mathrm{IQR} = 33.3 - 17.5 = 15.8$ marks, and $\mathrm{SIQR} = \tfrac{15.8}{2} = 7.9$ marks. For the 11 ungrouped marks: $\mathrm{IQR} = 22 - 8 = 14$.
\end{exemplebox}

In the same spirit, the \cle{percentile range} $P_{90} - P_{10}$ measures the spread of the middle $80\%$ of the data, discarding the top and bottom tenths. In our example, $P_{10}$ has position $8$, which falls exactly on the boundary $10$ (interpolation gives $10$), and $P_{90}$ has position $72$, falling on the boundary $40$; hence $P_{90} - P_{10} = 40 - 10 = 30$ marks. Examiners like it because it resists the influence of one or two abnormal scores even better than the IQR.

\courssub{Box-and-Whisker Plots (beyond the strict syllabus)}

A \cle{box-and-whisker plot} summarises a distribution with five numbers: minimum, $Q_1$, median $Q_2$, $Q_3$, maximum. The \emph{box} spans the IQR; the \emph{whiskers} extend to the extremes. It is not explicitly required by the Cameroon GCE syllabus, but it is an excellent revision tool and frequently appears in international examinations.

\begin{popfigure}
\begin{center}
\begin{tikzpicture}[scale=0.32]
\draw[popline] (0,0) -- (32,0);
\foreach \x in {0,5,10,15,20,25,30}
  \draw[popdark] (\x,0.15) -- (\x,-0.15) node[below] {\small \x};
\draw[very thick, popblue] (4,2) -- (8,2);
\draw[very thick, popblue, fill=popblueL] (8,1) rectangle (22,3);
\draw[very thick, poporange] (15,1) -- (15,3);
\draw[very thick, popblue] (22,2) -- (30,2);
\draw[very thick, popblue] (4,1.4) -- (4,2.6);
\draw[very thick, popblue] (30,1.4) -- (30,2.6);
\node[above, popdark] at (4,2.7) {\small min $=4$};
\node[above, popdark] at (8,3.1) {\small $Q_1=8$};
\node[above, poporange] at (15,3.1) {\small $Q_2=15$};
\node[above, popdark] at (22,3.1) {\small $Q_3=22$};
\node[above, popdark] at (30,2.7) {\small max $=30$};
\end{tikzpicture}
\end{center}
\legende{Box-and-whisker plot of the 11 marks, aligned with a number line. The box length is the IQR.}
\end{popfigure}

\courssub{Interpretation: Position within a Group}

Percentiles are the natural language of ranking. If a student's mark equals $P_{90}$ of the class, she is in the \cle{top decile}: at least $90\%$ of the candidates scored at or below her mark, and only about $10\%$ scored above. Similarly, scoring below $Q_1$ places a candidate in the bottom quarter --- a signal that remediation is needed --- while scoring above $Q_3$ places him in the top quarter. Notice the subtlety: a percentile describes a \emph{position relative to the group}, not an absolute level of mastery. Being at $P_{95}$ in a weak school may correspond to a lower actual mark than being at $P_{60}$ in a strong one.

\begin{aretenir}[Key points]
\begin{itemize}
\item $Q_1, Q_2, Q_3$ split ordered data into quarters; $Q_2$ is the median; $Q_1 = P_{25}$, $Q_3 = P_{75}$, $D_5 = P_{50}$.
\item For grouped data: position first ($\tfrac{iN}{4}$, $\tfrac{kN}{10}$, $\tfrac{kN}{100}$), then the class, then interpolation $Q_i = L + \frac{iN/4 - F}{f}\times c$.
\item $\mathrm{IQR} = Q_3 - Q_1$; $\mathrm{SIQR} = \tfrac{Q_3 - Q_1}{2}$; percentile range $= P_{90} - P_{10}$.
\item All these measures are \textbf{robust}: extreme values do not affect them.
\item Quartiles and percentiles can be read graphically from the ogive.
\end{itemize}
\end{aretenir}

\begin{exercice}[Practice]
Using the table of the 80 candidates above, find: (a) $D_3$; (b) $P_{40}$; (c) the interquartile range if the class widths were doubled to 20 with frequencies $24, 48, 8$ in classes $0$--$20$, $20$--$40$, $40$--$50$. Comment on the effect of grouping.
\end{exercice}

\begin{corrige}[Exercise 1]
(a) Position $\tfrac{3N}{10} = 24$: this falls exactly on the boundary of the second class (cumulative frequency 24), so $D_3 = 20$. (Interpolation: class $10$--$20$, $L=10$, $F=8$, $f=16$, $c=10$ gives $10 + \frac{24-8}{16}\times10 = 20$.)
(b) Position $\tfrac{40N}{100} = 32$; class $20$--$30$: $P_{40} = 20 + \frac{32-24}{30}\times 10 = 20 + 2.667 \approx 22.7$.
(c) New grouping: $N=80$, classes $0$--$20$ ($f=24$, $c=20$), $20$--$40$ ($f=48$, $c=20$), $40$--$50$ ($f=8$, $c=10$).
$Q_1$: position 20, class $0$--$20$, $L=0$, $F=0$, $f=24$: $Q_1 = 0 + \frac{20}{24}\times 20 \approx 16.7$.
$Q_3$: position 60, class $20$--$40$, $L=20$, $F=24$, $f=48$: $Q_3 = 20 + \frac{60-24}{48}\times 20 = 35$.
Hence $\mathrm{IQR} \approx 35 - 16.7 = 18.3$ instead of $15.8$: wider classes give a coarser (and here larger) estimate of the spread.
\end{corrige}

\coursec{Index Numbers}

In the previous sections we learnt how to summarise a \emph{single} set of data at a \emph{single} moment: quartiles, deciles and percentiles told us where values stand within one distribution. But in real life, the most urgent questions are questions of \emph{change}. Is garri more expensive in Douala this year than last year? By how much has the cost of living risen in Yaoundé since 2020? Has cocoa production in the Centre Region increased or decreased? A single statistic cannot answer such questions, because prices and quantities of many different goods change at the same time and in different proportions. We therefore need a new tool: the \cle{index number}.

\courssub{The Idea of an Index Number}

\begin{experience}[A market observation]
A student notes the price of a basin of garri at Mokolo market: FCFA $8\,000$ in $2022$ and FCFA $10\,000$ in $2024$. The price has risen by FCFA $2\,000$, that is by $\frac{2000}{8000}\times 100 = 25\%$. We say the \emph{price relative} is $1.25$, or, multiplied by $100$, the \emph{price index} is $125$. The number $125$ tells us at a glance: ``the price is $25\%$ higher than in the reference year''. Index numbers simply generalise this idea to several goods at once.
\end{experience}

\begin{definition}[Index number and base period]
An \cle{index number} is a statistical measure that expresses the value of a variable (price, quantity, production, value) at a given period as a percentage of its value at a reference period called the \cle{base period} (or base year). The index at the base period is conventionally $100$.
\begin{itemize}
\item An index \textbf{above} $100$ indicates an \textbf{increase} relative to the base period.
\item An index \textbf{below} $100$ indicates a \textbf{decrease} relative to the base period.
\item An index of $100$ indicates \textbf{no change}.
\end{itemize}
\end{definition}

\begin{aretenir}[Reading an index]
An index of $118$ means an increase of $18\%$ since the base period; an index of $92$ means a decrease of $8\%$. Never say ``the price is $118$'' — the index has \emph{no unit}; it is a pure percentage relative to the base.
\end{aretenir}

\courssub{Simple Index Numbers}

\begin{definition}[Simple price, quantity and value indices]
Let $p_0$, $q_0$ be the price and quantity of an item in the base year, and $p_n$, $q_n$ the corresponding values in the current year.
\begin{itemize}
\item \cle{Simple price index} (price relative): $\displaystyle I_p = \frac{p_n}{p_0}\times 100$.
\item \cle{Simple quantity index}: $\displaystyle I_q = \frac{q_n}{q_0}\times 100$.
\item \cle{Simple value index}: $\displaystyle I_v = \frac{p_n q_n}{p_0 q_0}\times 100 = I_p \times I_q / 100$.
\end{itemize}
\end{definition}

\begin{exemplebox}[Simple price index]
A litre of palm oil cost FCFA $1\,200$ in $2021$ (base year) and FCFA $1\,500$ in $2024$. Then
\[ I_p = \frac{1500}{1200}\times 100 = 125. \]
The price of palm oil has risen by $25\%$ since $2021$.
\end{exemplebox}

When several commodities are involved, two elementary constructions exist.

\begin{definition}[Simple aggregate index and average of relatives]
For $k$ commodities:
\begin{itemize}
\item \cle{Simple aggregate price index}: $\displaystyle I = \frac{\sum_{i=1}^k p_{n,i}}{\sum_{i=1}^k p_{0,i}}\times 100$ (ratio of total prices).
\item \cle{Simple average of price relatives}: $\displaystyle I = \frac{1}{k}\sum_{i=1}^k \frac{p_{n,i}}{p_{0,i}}\times 100$ (arithmetic mean of the individual indices).
\end{itemize}
\end{definition}

\begin{attention}[Limitations of simple indices]
Simple indices treat every item as equally important, which is unrealistic: a household spends far more on garri than on salt. They are also sensitive to the units of measurement (e.g., price per kg vs price per tonne). This motivates \emph{weighted} indices.
\end{attention}

\courssub{Weighted Index Numbers}

\begin{definition}[Weighted index number]
A \cle{weighted index number} assigns to each commodity a \cle{weight} reflecting its relative importance (usually the quantity consumed, the value spent, or the expenditure share). A weighted price index has the general form
\[ I = \frac{\sum p_n w}{\sum p_0 w}\times 100, \]
where the $w$ are the weights. The whole question is: \emph{which} quantities should serve as weights — those of the base year or those of the current year?
\end{definition}

\begin{methode}[Choosing weights in a weighted index]
\begin{enumerate}
\item Identify the purpose of the index (cost of living, production, sales value).
\item Choose weights that reflect the \textbf{importance} of each item: quantities consumed ($q$), expenditure shares ($p_0q_0$), or output volumes.
\item Decide the \textbf{period} of the weights: base-period quantities $q_0$ give the \emph{Laspeyres} index; current-period quantities $q_n$ give the \emph{Paasche} index; geometric mean of both gives \emph{Fisher's ideal index}.
\item Keep the \textbf{same} weights in the numerator and in the denominator: only the \emph{prices} must change between top and bottom.
\item State clearly the base year and the weights used, so the index can be interpreted and reproduced.
\end{enumerate}
\end{methode}

\courssub{Laspeyres and Paasche Price Indices}

\begin{propriete}[Laspeyres price index]
The \cle{Laspeyres price index} uses \textbf{base-period quantities} $q_0$ as weights:
\[ I_L = \frac{\sum p_n q_0}{\sum p_0 q_0}\times 100. \]
It answers the question: ``How much would the \emph{old} basket of goods cost at \emph{new} prices, compared with its cost at old prices?'' It measures the change in cost of purchasing a \emph{fixed} base-period basket.
\end{propriete}

\begin{propriete}[Paasche price index]
The \cle{Paasche price index} uses \textbf{current-period quantities} $q_n$ as weights:
\[ I_P = \frac{\sum p_n q_n}{\sum p_0 q_n}\times 100. \]
It answers the question: ``How much does the \emph{present} basket cost at new prices, compared with what the \emph{same present} basket would have cost at old prices?'' It measures the change in cost of purchasing the \emph{current} basket.
\end{propriete}

\begin{propriete}[Fisher's ideal price index]
The \cle{Fisher's ideal price index} is the geometric mean of the Laspeyres and Paasche indices:
\[ I_F = \sqrt{I_L \times I_P}. \]
It satisfies the \emph{time reversal} and \emph{factor reversal} tests, making it theoretically superior, though less used in practice due to computational complexity.
\end{propriete}

\begin{attention}[Do not mix the two periods!]
A very common examination error is to write something like $\dfrac{\sum p_n q_0}{\sum p_0 q_n}\times 100$, mixing base-year quantities in the numerator with current-year quantities in the denominator. This is meaningless: within \emph{one} formula, the quantities must all come from the \emph{same} period. Check: Laspeyres $\to$ $q_0$ everywhere; Paasche $\to$ $q_n$ everywhere.
\end{attention}

\begin{exoresolu}[Laspeyres, Paasche and Fisher indices for a market basket]
The table below gives the prices (FCFA per unit) and quantities consumed of four goods in $2022$ (base year, index $0$) and $2024$ (current year, index $n$).
\begin{center}
\begin{tabular}{|l|c|c|c|c|}
\hline
\rowcolor{popblueL} Good & $p_0$ & $q_0$ & $p_n$ & $q_n$ \\
\hline
Garri (basin) & $8\,000$ & $12$ & $10\,000$ & $10$ \\
\hline
Plantain (bunch) & $2\,000$ & $20$ & $2\,500$ & $24$ \\
\hline
Palm oil (litre) & $1\,200$ & $30$ & $1\,500$ & $28$ \\
\hline
Fish (kg) & $2\,500$ & $8$ & $3\,500$ & $6$ \\
\hline
\end{tabular}
\end{center}
Compute the Laspeyres, Paasche and Fisher price indices for $2024$ (base $2022 = 100$) and interpret the results.
\tcblower
\textbf{Solution.} We build the four products needed.
\begin{align*}
\sum p_0 q_0 &= 8000(12)+2000(20)+1200(30)+2500(8)\\
&= 96\,000+40\,000+36\,000+20\,000 = 192\,000,\\
\sum p_n q_0 &= 10\,000(12)+2500(20)+1500(30)+3500(8)\\
&= 120\,000+50\,000+45\,000+28\,000 = 243\,000,\\
\sum p_0 q_n &= 8000(10)+2000(24)+1200(28)+2500(6)\\
&= 80\,000+48\,000+33\,600+15\,000 = 176\,600,\\
\sum p_n q_n &= 10\,000(10)+2500(24)+1500(28)+3500(6)\\
&= 100\,000+60\,000+42\,000+21\,000 = 223\,000.
\end{align*}
\textbf{Laspeyres index:}
\[ I_L = \frac{\sum p_n q_0}{\sum p_0 q_0}\times 100 = \frac{243\,000}{192\,000}\times 100 = 126.5625 \approx 126.6. \]
\textbf{Paasche index:}
\[ I_P = \frac{\sum p_n q_n}{\sum p_0 q_n}\times 100 = \frac{223\,000}{176\,600}\times 100 = 126.2741 \approx 126.3. \]
\textbf{Fisher's ideal index:}
\[ I_F = \sqrt{I_L \times I_P} = \sqrt{126.5625 \times 126.2741} = \sqrt{15980.5} \approx 126.4. \]
\textbf{Interpretation.} All three indices exceed $100$: prices have risen by about $26\%$ between $2022$ and $2024$. Laspeyres says the \emph{old} basket ($q_0$) would cost $26.6\%$ more today; Paasche says the \emph{current} basket ($q_n$) costs $26.3\%$ more than it would have at $2022$ prices; Fisher gives a compromise at $26.4\%$. The two answers are close here, but they need not be in general.
\end{exoresolu}

\begin{popfigure}
\begin{tikzpicture}
\begin{axis}[
    width=0.85\linewidth, height=6.2cm,
    xlabel={Year}, ylabel={Consumer price index (base $2019=100$)},
    xmin=2019, xmax=2024, ymin=95, ymax=135,
    xtick={2019,2020,2021,2022,2023,2024},
    ytick={100,110,120,130},
    grid=both, grid style={popline, dashed},
    legend style={at={(0.02,0.98)}, anchor=north west, font=\small},
]
\addplot[color=popblue, very thick, mark=*, mark size=2pt] coordinates {
    (2019,100) (2020,104) (2021,109) (2022,116) (2023,124) (2024,131)
};
\addlegendentry{Price index}
\addplot[color=popmuted, dashed, domain=2019:2024] {100};
\end{axis}
\end{tikzpicture}
\legende{A consumer price index rising from $100$ (base year $2019$) to $131$ in $2024$: an overall increase of $31\%$ in five years.}
\end{popfigure}

\courssub{Quantity and Value Indices}

The same weighting principles apply to quantities and values.

\begin{propriete}[Laspeyres and Paasche quantity indices]
\begin{itemize}
\item \cle{Laspeyres quantity index}: $\displaystyle I_{L,q} = \frac{\sum q_n p_0}{\sum q_0 p_0}\times 100$ (base-period prices as weights).
\item \cle{Paasche quantity index}: $\displaystyle I_{P,q} = \frac{\sum q_n p_n}{\sum q_0 p_n}\times 100$ (current-period prices as weights).
\end{itemize}
\end{propriete}

\begin{propriete}[Value index]
The \cle{value index} compares total expenditure (or revenue) between two periods:
\[ I_v = \frac{\sum p_n q_n}{\sum p_0 q_0}\times 100. \]
Note that $I_v = \frac{I_p \times I_q}{100}$ only when the \emph{same} weighting system is used for both price and quantity indices (e.g., both Laspeyres or both Paasche). In general, $I_v \neq I_L \times I_{L,q} / 100$.
\end{propriete}

\courssub{Chain-Base vs Fixed-Base Indices}

\begin{definition}[Fixed-base and chain-base indices]
\begin{itemize}
\item A \cle{fixed-base index} compares every period to the \emph{same} base period (e.g., $2019=100$ for all years).
\item A \cle{chain-base index} compares each period to the \emph{immediately preceding} period (link relatives), then chains them together: $I_{0\to n} = \frac{I_{0\to1}}{100} \times \frac{I_{1\to2}}{100} \times \cdots \times \frac{I_{n-1\to n}}{100} \times 100$.
\end{itemize}
\end{definition}

\begin{aretenir}[When to use which]
Fixed-base indices are standard for long-term comparisons and official publications (CPI). Chain-base indices are useful when the basket of goods changes frequently, as they allow updating weights each period without breaking the series.
\end{aretenir}

\courssub{Comparing Laspeyres and Paasche}

\begin{aretenir}[Laspeyres versus Paasche]
\begin{itemize}
\item \textbf{Laspeyres} ($q_0$ weights): easy to compute repeatedly, since the basket is fixed once and for all; but it tends to \emph{overstate} inflation because it ignores the fact that consumers substitute away from goods whose prices rise fastest (substitution bias).
\item \textbf{Paasche} ($q_n$ weights): reflects current consumption habits, but requires new quantity data every period and tends to \emph{understate} inflation because it uses the basket that has already adjusted to price changes.
\item \textbf{Fisher} (geometric mean): theoretically ideal, satisfies time and factor reversal tests, but requires both sets of quantities.
\item Both Laspeyres and Paasche are valid; they answer \emph{different} questions. In practice, most national consumer price indices are of Laspeyres type, with the basket updated periodically (e.g., every 5 years).
\end{itemize}
\end{aretenir}

\begin{savaistu}[Index numbers in Cameroon]
In Cameroon, the National Institute of Statistics (INS) publishes a \cle{consumer price index} (CPI) that tracks the cost of a fixed basket of goods and services — food, housing, transport, clothing — in cities such as Yaoundé and Douala. Government and employers use it for \cle{cost-of-living adjustments}: when the CPI rises sharply, salaries and pensions may be revised so that purchasing power is preserved. The CPI is essentially a Laspeyres-type index: a fixed basket, repriced every month. The INS also produces \cle{producer price indices} (PPI) for agriculture and industry, and \cle{quantity indices} for major export crops like cocoa, coffee, and cotton.
\end{savaistu}

\begin{exercice}[Computing and interpreting indices]
The table below shows the prices (FCFA) and quantities of three items bought by a family in $2021$ (base) and $2023$.
\begin{center}
\begin{tabular}{|l|c|c|c|c|}
\hline
\rowcolor{popblueL} Item & $p_0$ & $q_0$ & $p_n$ & $q_n$ \\
\hline
Rice (kg) & $500$ & $40$ & $650$ & $36$ \\
\hline
Beans (kg) & $800$ & $15$ & $1\,000$ & $18$ \\
\hline
Oil (litre) & $1\,100$ & $10$ & $1\,400$ & $12$ \\
\hline
\end{tabular}
\end{center}
\begin{enumerate}
\item Compute the simple price index of each item for $2023$ (base $2021 = 100$).
\item Compute the Laspeyres price index.
\item Compute the Paasche price index.
\item Compute Fisher's ideal price index.
\item Interpret the two indices in one sentence each.
\end{enumerate}
\end{exercice}

\begin{corrige}[Exercise 1]
\textbf{1.} Simple price indices:
\begin{itemize}
\item Rice: $\frac{650}{500}\times100 = 130.0$
\item Beans: $\frac{1000}{800}\times100 = 125.0$
\item Oil: $\frac{1400}{1100}\times100 = 127.3$
\end{itemize}

\textbf{2.} Laspeyres:
\begin{align*}
\sum p_0 q_0 &= 500(40)+800(15)+1100(10) = 20\,000+12\,000+11\,000 = 43\,000,\\
\sum p_n q_0 &= 650(40)+1000(15)+1400(10) = 26\,000+15\,000+14\,000 = 55\,000,\\
I_L &= \frac{55\,000}{43\,000}\times 100 = 127.9069 \approx 127.9.
\end{align*}

\textbf{3.} Paasche:
\begin{align*}
\sum p_0 q_n &= 500(36)+800(18)+1100(12) = 18\,000+14\,400+13\,200 = 45\,600,\\
\sum p_n q_n &= 650(36)+1000(18)+1400(12) = 23\,400+18\,000+16\,800 = 58\,200,\\
I_P &= \frac{58\,200}{45\,600}\times 100 = 127.6316 \approx 127.6.
\end{align*}

\textbf{4.} Fisher's ideal index:
\[ I_F = \sqrt{127.9069 \times 127.6316} = \sqrt{16325.5} \approx 127.8. \]

\textbf{5.} Laspeyres: the $2021$ basket would cost about $27.9\%$ more in $2023$. Paasche: the $2023$ basket costs about $27.6\%$ more than it would have cost in $2021$. Fisher: a compromise estimate of $27.8\%$ increase. All three confirm a strong rise in the cost of living, close to $28\%$.
\end{corrige}

\begin{attention}[Final warnings for the examination]
\begin{itemize}
\item Always \textbf{state the base year} and give the index to at least one decimal place.
\item An index of $126.6$ means a $26.6\%$ \emph{increase}, not a $126.6\%$ increase.
\item Do not add units (FCFA) to an index number.
\item Keep quantities consistent within one formula: $q_0$ for Laspeyres, $q_n$ for Paasche, in \emph{both} numerator and denominator.
\item For Fisher's index, use the \emph{exact} values of $I_L$ and $I_P$ (not rounded) before taking the square root.
\item Distinguish clearly between \emph{price}, \emph{quantity} and \emph{value} indices in your notation and interpretation.
\end{itemize}
\end{attention}

Index numbers thus condense the evolution of a whole economy into a single, readable figure. In the next section, we return to the \emph{shape} of distributions and learn how to measure their asymmetry through the notion of \emph{skewness}.

\coursec{Skewness and Shape of Distributions}

In the previous section we used index numbers to summarise how prices and quantities evolve over time. We now return to the data themselves and ask a final, very visual question: \emph{what does the shape of a distribution tell us?} Two classes may have the same mean mark of $12/20$, yet in one class the marks are balanced around $12$ while in the other most students scored below $12$ with a few brilliant candidates pulling the average up. The mean alone cannot distinguish these situations. The missing ingredient is \cle{skewness} — a measure of the \emph{asymmetry} of a distribution.

\courssub{Symmetric and Skewed Distributions}

\begin{experience}[Looking at shapes]
Consider the marks of two groups of Upper Sixth students in a mock examination.

Group A: $8, 9, 10, 10, 11, 11, 11, 12, 12, 12, 13, 13, 14, 14, 15$.

Group B: $4, 5, 6, 7, 8, 8, 9, 9, 10, 10, 11, 12, 14, 16, 19$.

For Group A the frequencies rise gently to a peak in the middle and fall off in the same way on both sides: the histogram is a \emph{mirror image} about its centre. For Group B the bulk of the marks sits at the lower end, but a long \emph{tail} stretches towards the high marks ($14, 16, 19$). The two distributions have completely different shapes, and this difference has practical consequences: in Group B, the "typical" student is weaker than the mean suggests. Indeed, both groups have mean $11.5$, but the median of Group A is $12$ while that of Group B is only $9$: the few high marks in Group B drag the mean upwards without representing most students.
\end{experience}

\begin{definition}[Symmetric distribution]
A distribution is \cle{symmetric} if its graph (histogram, frequency polygon or frequency curve) can be divided into two halves that are mirror images of each other about a vertical line through the centre. The two tails are of equal length and shape.
\end{definition}

\begin{definition}[Skewness]
\cle{Skewness} is the degree of asymmetry of a distribution. A distribution that is not symmetric is said to be \cle{skewed}.
\begin{itemize}
\item A distribution is \cle{positively skewed} (or \emph{skewed to the right}) if it has a long tail stretching towards the \textbf{higher} values, with most observations concentrated at the lower end.
\item A distribution is \cle{negatively skewed} (or \emph{skewed to the left}) if it has a long tail stretching towards the \textbf{lower} values, with most observations concentrated at the upper end.
\end{itemize}
\end{definition}

\begin{attention}[The tail decides — not the peak!]
The most common mistake at the GCE is to name the skewness after the side where the \emph{bulk} of the data lies. This is wrong. The direction of skewness is the direction of the \textbf{long tail}:
\begin{itemize}
\item tail to the \textbf{right} (high values) $\Rightarrow$ \textbf{positive} skewness;
\item tail to the \textbf{left} (low values) $\Rightarrow$ \textbf{negative} skewness.
\end{itemize}
A useful memory aid: the skewness "pulls" in the direction of the tail, exactly as the few extreme values pull the mean towards them.
\end{attention}

\begin{popfigure}
\begin{center}
\begin{tikzpicture}
\begin{axis}[
  width=14cm, height=6.5cm,
  axis lines=left,
  xlabel={value}, ylabel={frequency density},
  xtick=\empty, ytick=\empty,
  xmin=0, xmax=10, ymin=0, ymax=0.62,
  legend style={at={(0.98,0.97)}, anchor=north east, font=\small},
  every axis plot/.append style={very thick}
]
% symmetric (normal-like)
\addplot[popblue, smooth, domain=0:10, samples=100] {0.55*exp(-((x-5)^2)/2)};
\addlegendentry{symmetric}
% positively skewed (right tail)
\addplot[poporange, smooth, domain=0:10, samples=120] {0.9*x*exp(-x)*1.35};
\addlegendentry{positively skewed}
% negatively skewed (left tail)
\addplot[popgreen, smooth, domain=0:10, samples=120] {0.9*(10-x)*exp(-(10-x))*1.35};
\addlegendentry{negatively skewed}
\end{axis}
\end{tikzpicture}
\end{center}
\legende{Three typical shapes: a symmetric bell-shaped curve, a positively skewed curve (long right tail) and a negatively skewed curve (long left tail).}
\end{popfigure}

\courssub{Relative Positions of Mean, Median and Mode}

The three measures of central tendency react differently to extreme values. The \cle{mode} sits at the peak of the distribution and ignores the tails completely. The \cle{median} depends only on the order of the data, so it is only slightly affected by a few extreme values. The \cle{mean}, however, uses every value, so it is \emph{dragged} towards the long tail. This simple observation explains the following fundamental property.

\begin{propriete}[Mean, median and mode according to shape]
\begin{itemize}
\item For a \textbf{symmetric} (unimodal) distribution:
\[ \text{mean} = \text{median} = \text{mode}. \]
\item For a \textbf{positively skewed} distribution (tail to the right):
\[ \text{mode} < \text{median} < \text{mean}. \]
\item For a \textbf{negatively skewed} distribution (tail to the left):
\[ \text{mean} < \text{median} < \text{mode}. \]
\end{itemize}
Moreover, for moderately skewed distributions, the median usually lies about one third of the way from the mode to the mean, which gives the empirical relation
\[ \bar{x} - M_o \approx 3(\bar{x} - M_e), \]
and motivates the second Pearson coefficient below. (This property is quoted, not proved, at the GCE.)
\end{propriete}

\begin{popfigure}
\begin{center}
\begin{tikzpicture}[scale=1.0]
% ---- Symmetric ----
\begin{scope}
\draw[->] (0,0) -- (4.6,0);
\draw[very thick, popblue, domain=0.4:4.2, samples=60] plot (\x, {1.5*exp(-((\x-2.3)^2)/0.8)});
\draw[dashed, popdark] (2.3,0) -- (2.3,1.7);
\node[below, font=\small, align=center] at (2.3,0) {mean = median\\ = mode};
\node[font=\small\bfseries, popblue] at (2.3,-1.4) {Symmetric};
\end{scope}
% ---- Positive skew ----
\begin{scope}[xshift=6.2cm]
\draw[->] (0,0) -- (4.6,0);
\draw[very thick, poporange, domain=0.1:4.5, samples=80] plot (\x, {2.6*\x*exp(-\x*1.4)});
\draw[dashed, popdark] (0.72,0) -- (0.72,1.05);
\draw[dashed, popdark] (1.15,0) -- (1.15,0.86);
\draw[dashed, popdark] (1.75,0) -- (1.75,0.55);
\node[below, font=\small] at (0.72,0) {mode};
\node[below, font=\small] at (1.15,-0.45) {median};
\node[below, font=\small] at (1.75,-0.9) {mean};
\node[font=\small\bfseries, poporange] at (2.3,-1.4) {Positively skewed};
\end{scope}
% ---- Negative skew ----
\begin{scope}[xshift=12.4cm]
\draw[->] (0,0) -- (4.6,0);
\draw[very thick, popgreen, domain=0.1:4.5, samples=80] plot (\x, {2.6*(4.6-\x)*exp(-(4.6-\x)*1.4)});
\draw[dashed, popdark] (2.85,0) -- (2.85,0.55);
\draw[dashed, popdark] (3.45,0) -- (3.45,0.86);
\draw[dashed, popdark] (3.88,0) -- (3.88,1.05);
\node[below, font=\small] at (2.85,-0.9) {mean};
\node[below, font=\small] at (3.45,-0.45) {median};
\node[below, font=\small] at (3.88,0) {mode};
\node[font=\small\bfseries, popgreen] at (2.3,-1.4) {Negatively skewed};
\end{scope}
\end{tikzpicture}
\end{center}
\legende{Ordering of mean, median and mode for each shape. The mean is always pulled towards the long tail.}
\end{popfigure}

\begin{aretenir}[Why the mean moves most]
The mean is the \emph{centre of gravity} of the distribution: a single very large value (a millionaire in a village, a mark of $19$ in a weak class) pulls it strongly. The median only counts positions, and the mode only looks at the peak. Hence: \textbf{mean is most affected, mode least affected, median in between.}
\end{aretenir}

\courssub{Pearson's Coefficients of Skewness}

Visual inspection is useful, but the examiner expects a \emph{numerical} measure. Karl Pearson proposed coefficients based on the distance between the mean and the mode (or median), standardised by the standard deviation $\sigma$ so that the result is a pure number with no units. Dividing by $\sigma$ is essential: a gap of $2$ marks between mean and mode is large when $\sigma = 1$ but negligible when $\sigma = 10$.

\begin{propriete}[Pearson's coefficients of skewness]
For a distribution with mean $\bar{x}$, mode $M_o$, median $M_e$ and standard deviation $\sigma \neq 0$:
\[ S_k = \frac{\bar{x} - M_o}{\sigma} \qquad\text{(first coefficient)} \]
\[ S_k = \frac{3(\bar{x} - M_e)}{\sigma} \qquad\text{(second coefficient)} \]
The second form is preferred when the mode is ill-defined or difficult to estimate (common for grouped data); the factor $3$ comes from the empirical relation $\bar{x} - M_o \approx 3(\bar{x} - M_e)$ for moderately skewed distributions. (Formulae quoted; not examinable as a proof.)
\end{propriete}

\textbf{Interpretation.}
\begin{itemize}
\item $S_k > 0$: positively skewed (mean above mode/median, right tail).
\item $S_k = 0$: symmetric.
\item $S_k < 0$: negatively skewed (left tail).
\end{itemize}
As a rule of thumb, $|S_k| < 0.5$ indicates a fairly symmetric distribution, $0.5 \le |S_k| \le 1$ a moderately skewed one, and $|S_k| > 1$ a highly skewed one.

\begin{exoresolu}[Pearson's coefficient for examination marks]
In a mock GCE examination, the marks of a class give: mean $\bar{x} = 11.4$, median $M_e = 12.0$, standard deviation $\sigma = 3.0$. Compute Pearson's second coefficient of skewness and comment.
\tcblower
\textbf{Solution.}
\[ S_k = \frac{3(\bar{x} - M_e)}{\sigma} = \frac{3(11.4 - 12.0)}{3.0} = \frac{3 \times (-0.6)}{3.0} = -0.6. \]
Since $S_k < 0$, the distribution is \textbf{negatively skewed}: most students obtained high marks, with a tail of weaker marks stretching to the left. Since $|S_k| = 0.6$ lies between $0.5$ and $1$, the skewness is \emph{moderate}. Note the consistency check: mean $<$ median, exactly as expected for a left-skewed distribution.
\end{exoresolu}

\begin{exoresolu}[Pearson's first coefficient from raw data]
The daily sales (in thousands of FCFA) of a small shop in Bamenda over one week were:
\[ 20,\ 22,\ 22,\ 23,\ 24,\ 25,\ 40. \]
Compute Pearson's first coefficient of skewness and interpret it.
\tcblower
\textbf{Solution.} The mode is clearly $M_o = 22$. The mean is
\[ \bar{x} = \frac{20 + 22 + 22 + 23 + 24 + 25 + 40}{7} = \frac{176}{7} \approx 25.14. \]
For the standard deviation we compute the squared deviations from $\bar{x} \approx 25.14$:
\begin{align*}
\sum (x_i - \bar{x})^2 &\approx (20-25.14)^2 + 2(22-25.14)^2 + (23-25.14)^2 \\
&\quad + (24-25.14)^2 + (25-25.14)^2 + (40-25.14)^2 \\
&\approx 26.42 + 19.72 + 4.58 + 1.30 + 0.02 + 220.90 = 272.94,
\end{align*}
so $\sigma = \sqrt{272.94/7} \approx \sqrt{38.99} \approx 6.24$. Hence
\[ S_k = \frac{\bar{x} - M_o}{\sigma} \approx \frac{25.14 - 22}{6.24} \approx 0.50. \]
Since $S_k > 0$, the sales are \textbf{positively skewed}: the single exceptional day of $40\,000$ FCFA pulls the mean above the mode. The shopkeeper should quote the mode or median, not the mean, to describe a "typical" day.
\end{exoresolu}

\courssub{Quartile Coefficient of Skewness}

Pearson's coefficients depend on the mean and standard deviation, which are themselves sensitive to extreme values. A measure based only on quartiles is more \emph{robust}. Recall that for a symmetric distribution the median $Q_2$ lies exactly halfway between $Q_1$ and $Q_3$, i.e. $Q_3 - Q_2 = Q_2 - Q_1$. Any imbalance between these two gaps signals skewness: if the upper gap is wider, the right tail is longer and the skew is positive.

\begin{propriete}[Quartile coefficient of skewness (Bowley's coefficient)]
For a distribution with quartiles $Q_1$, $Q_2$ (median) and $Q_3$:
\[ S_{kQ} = \frac{(Q_3 - Q_2) - (Q_2 - Q_1)}{Q_3 - Q_1} = \frac{Q_3 + Q_1 - 2Q_2}{Q_3 - Q_1}. \]
It always satisfies $-1 \le S_{kQ} \le 1$. The sign is interpreted exactly as for Pearson's coefficient: positive for right skew, negative for left skew, zero for symmetry. (Formula quoted; not examinable as a proof.)
\end{propriete}

\begin{exoresolu}[Quartile coefficient for waiting times]
The waiting times (in minutes) of patients at a district hospital give $Q_1 = 8$, $Q_2 = 15$, $Q_3 = 34$. Compute the quartile coefficient of skewness and interpret it.
\tcblower
\textbf{Solution.}
\[ S_{kQ} = \frac{Q_3 + Q_1 - 2Q_2}{Q_3 - Q_1} = \frac{34 + 8 - 2(15)}{34 - 8} = \frac{42 - 30}{26} = \frac{12}{26} \approx 0.46. \]
Since $S_{kQ} > 0$, the distribution of waiting times is \textbf{positively skewed}: half the patients wait at most $15$ minutes, but a minority waits much longer (the gap $Q_3 - Q_2 = 19$ exceeds $Q_2 - Q_1 = 7$). The skewness is moderate since $|S_{kQ}| < 0.5$.
\end{exoresolu}

\begin{attention}[Common mistakes with skewness coefficients]
\begin{itemize}
\item \textbf{Wrong sign convention.} $S_k > 0$ means the tail is on the \emph{right} (positive skew). Do not describe such a distribution as "skewed to the left".
\item \textbf{Forgetting the factor 3} in Pearson's second coefficient: it is $3(\bar{x} - M_e)/\sigma$, not $(\bar{x} - M_e)/\sigma$.
\item \textbf{Mixing the two quartile gaps.} In Bowley's formula the numerator is $(Q_3 - Q_2) - (Q_2 - Q_1)$; reversing the order changes the sign and reverses the conclusion.
\item \textbf{Comparing coefficients from different formulae.} Pearson's and Bowley's coefficients measure skewness on different scales; do not expect them to be equal for the same data.
\item \textbf{Claiming skewness from the mean alone.} You must compare at least two measures (mean vs median, or the two quartile gaps) before concluding.
\item \textbf{Using the coefficient without units check.} All three coefficients are pure numbers; if your answer carries a unit (marks, minutes, FCFA), you have forgotten to divide by $\sigma$ or by $Q_3 - Q_1$.
\end{itemize}
\end{attention}

\courssub{Skewness in Real Life}

\begin{savaistu}[Where do skewed distributions appear?]
\begin{itemize}
\item \textbf{Incomes.} In Cameroon, as in most countries, the distribution of household incomes is strongly \emph{positively skewed}: most households earn modest incomes, while a small number earn very large ones. This is why the \emph{median} income, not the mean, is usually quoted in economic reports — the mean is inflated by the extreme tail and would overstate the "typical" income.
\item \textbf{Examination marks.} An easy paper produces a \emph{negatively skewed} mark distribution (most marks high, a tail of low marks); a very difficult paper produces a \emph{positively skewed} one. Examiners use this to judge whether a paper was well calibrated.
\item \textbf{Waiting times and lifetimes.} Waiting times at a hospital or a bank, and the lifetimes of electronic components, are typically positively skewed: many short waits, a few very long ones.
\end{itemize}
\end{savaistu}

\begin{exercice}[Identifying and measuring skewness]
The ages of workers in a small company in Douala are summarised by: mean $= 36.2$ years, median $= 33.0$ years, mode $= 29.5$ years, $\sigma = 9.6$ years, $Q_1 = 27$, $Q_2 = 33$, $Q_3 = 44$.
\begin{enumerate}
\item Without any calculation, state the expected direction of skewness from the positions of the mean, median and mode.
\item Compute Pearson's first coefficient of skewness.
\item Compute Pearson's second coefficient of skewness.
\item Compute Bowley's quartile coefficient of skewness.
\item Comment on the consistency of your answers.
\end{enumerate}
\end{exercice}

\begin{corrige}[Exercise 1]
\begin{enumerate}
\item Since mode $(29.5) <$ median $(33.0) <$ mean $(36.2)$, the distribution is \textbf{positively skewed} (tail towards the higher ages).
\item $S_k = \dfrac{\bar{x} - M_o}{\sigma} = \dfrac{36.2 - 29.5}{9.6} = \dfrac{6.7}{9.6} \approx 0.70$.
\item $S_k = \dfrac{3(\bar{x} - M_e)}{\sigma} = \dfrac{3(36.2 - 33.0)}{9.6} = \dfrac{9.6}{9.6} = 1.00$.
\item $S_{kQ} = \dfrac{Q_3 + Q_1 - 2Q_2}{Q_3 - Q_1} = \dfrac{44 + 27 - 66}{44 - 27} = \dfrac{5}{17} \approx 0.29$.
\item All three coefficients are \textbf{positive}, confirming a positively skewed distribution: most workers are relatively young, with a tail of older workers. The magnitudes differ because the coefficients use different scales, but the \emph{sign} — and hence the conclusion — is consistent.
\end{enumerate}
\end{corrige}

\begin{exercice}[From summary statistics to a conclusion]
Two schools enter candidates for the same examination. School X has $\bar{x} = 10.8$, $M_e = 10.5$, $\sigma = 2.4$. School Y has $Q_1 = 9$, $Q_2 = 12$, $Q_3 = 14$.
\begin{enumerate}
\item Compute Pearson's second coefficient for School X and interpret it.
\item Compute Bowley's coefficient for School Y and interpret it.
\item Which school's marks are more nearly symmetric? Justify.
\end{enumerate}
\end{exercice}

\begin{corrige}[Exercise 2]
\begin{enumerate}
\item $S_k = \dfrac{3(10.8 - 10.5)}{2.4} = \dfrac{0.9}{2.4} = 0.375$. Slight \textbf{positive} skewness: a small tail of high marks.
\item $S_{kQ} = \dfrac{14 + 9 - 2(12)}{14 - 9} = \dfrac{-1}{5} = -0.2$. Slight \textbf{negative} skewness: a small tail of low marks.
\item School Y: $|S_{kQ}| = 0.2 < 0.375 = |S_k|$, so its marks are more nearly symmetric (both distributions are only mildly skewed, since both coefficients are below $0.5$ in absolute value).
\end{enumerate}
\end{corrige}

\begin{aretenir}[Key points of the section]
\begin{itemize}
\item A distribution is \cle{symmetric} when its two halves are mirror images; then mean $=$ median $=$ mode.
\item \cle{Positive skewness}: long tail to the right, mode $<$ median $<$ mean. \cle{Negative skewness}: long tail to the left, mean $<$ median $<$ mode.
\item Pearson: $S_k = \dfrac{\bar{x} - M_o}{\sigma}$ or $S_k = \dfrac{3(\bar{x} - M_e)}{\sigma}$.
\item Bowley (quartiles): $S_{kQ} = \dfrac{Q_3 + Q_1 - 2Q_2}{Q_3 - Q_1}$, always between $-1$ and $1$.
\item Sign $>$ 0: right skew; sign $<$ 0: left skew; value near 0: near symmetry. The \textbf{tail}, never the peak, gives the direction.
\end{itemize}
\end{aretenir}

\newpage

\coursec{Integration activity}

\courssub{Aims of the activity}

This integration activity mobilises, in one authentic study, all the competences developed in this chapter. By the end of the activity, you should be able to:

\begin{itemize}
\item collect, organise and present real data in a \cle{grouped frequency table} with correct \cle{class boundaries};
\item draw and interpret a \cle{histogram} and a \cle{cumulative frequency curve} (ogive);
\item calculate and interpret measures of \cle{central tendency} (mean, median, mode) and of \cle{dispersion} (standard deviation, quartiles, semi-interquartile range) for grouped data;
\item construct and interpret \cle{Laspeyres} and \cle{Paasche} price indices;
\item measure the \cle{skewness} of a distribution using Pearson's coefficient and the quartile coefficient, and describe the shape of the distribution;
\item communicate statistical findings in a short advisory report and defend them orally.
\end{itemize}

\courssub{Situation}

\textbf{Statistical Audit of the GBHS Bafoussam Canteen.}\par\medskip

The canteen of Government Bilingual High School Bafoussam serves about 400 students every school day. The manager, Madam Nkwain, has noticed that her revenue is unstable and that some students complain that ``food has become too dear this year''. Before the next Parent--Teacher Association meeting, she asks the Upper Sixth Pure Mathematics with Statistics class to carry out a \emph{statistical audit} of the canteen.

Your class accepts the mandate and collects the following data during one school week.

\textbf{Data set A — Daily spending.} Thirty volunteer students record the total amount (in FCFA) each of them spends at the canteen per day. The grouped results are:

\begin{center}
\begin{tabular}{|c|c|}
\hline
\rowcolor{popblueL} Spending (FCFA) & Number of students \\
\hline
100 -- 199 & 3 \\
\hline
200 -- 299 & 6 \\
\hline
300 -- 399 & 9 \\
\hline
400 -- 499 & 7 \\
\hline
500 -- 599 & 4 \\
\hline
600 -- 699 & 1 \\
\hline
\end{tabular}
\end{center}

\textbf{Data set B — Prices and quantities.} For five representative items, the class records the average price and the quantity sold per day, for last year (base year, index 0) and this year (current year, index 1):

\begin{center}
\begin{tabular}{|l|c|c|c|c|}
\hline
\rowcolor{popblueL} Item & $p_0$ (FCFA) & $q_0$ & $p_1$ (FCFA) & $q_1$ \\
\hline
Plate of rice & 250 & 80 & 300 & 70 \\
\hline
Plate of beans & 150 & 60 & 200 & 55 \\
\hline
Puff-puff & 50 & 200 & 75 & 180 \\
\hline
Bottle of water & 100 & 150 & 150 & 140 \\
\hline
Bread & 100 & 90 & 125 & 85 \\
\hline
\end{tabular}
\end{center}

The class must now analyse these data and write a one-page report advising Madam Nkwain.

\courssub{Tasks}

\begin{enumerate}
\item Rewrite Data set A as a grouped frequency table showing \textbf{class boundaries}, \textbf{mid-values}, \textbf{cumulative frequencies} and the columns $fx$ and $fx^2$.
\item Draw the \textbf{histogram} of the spending distribution, using the class boundaries on the horizontal axis.
\item Construct the \textbf{cumulative frequency table} and draw the \textbf{ogive} (``less than'' curve). Use the ogive to estimate the median and the quartiles graphically.
\item Calculate the \textbf{mean}, \textbf{median} and \textbf{modal} daily spending, and state the \textbf{range} of the distribution.
\item Calculate the \textbf{variance} and the \textbf{standard deviation} of the daily spending.
\item Calculate $Q_1$, $Q_3$, the \textbf{interquartile range} and the \textbf{semi-interquartile range}. Compare your values of $Q_1$, $Q_2$, $Q_3$ with the graphical estimates of question 3.
\item Compute \textbf{Pearson's coefficient of skewness} and the \textbf{quartile coefficient of skewness}. Describe the shape of the spending distribution.
\item Using Data set B, compute the \textbf{Laspeyres price index} and the \textbf{Paasche price index} for this year (base: last year $= 100$). Interpret each index in one sentence for Madam Nkwain.
\item Write a one-page report (half a page of figures, half a page of advice) answering: How much does a typical student spend? Is spending spread out? Have prices really increased, and by how much? What do you advise the manager?
\item Each group presents its results in 5 minutes; the class compares the answers and discusses any discrepancies.
\end{enumerate}

\begin{methode}[Strategy for a grouped-data audit]
\begin{enumerate}
\item Start from the raw grouped table and build the \emph{working table}: boundaries, mid-values $x$, cumulative frequencies $F$, then $fx$ and $fx^2$. Every later calculation reads from this table.
\item Draw the histogram (equal class widths here, so height $=$ frequency) and the ogive (cumulative frequency against \emph{upper} class boundaries).
\item Compute the averages first (mean, median, mode), then the dispersion (variance, standard deviation, quartiles), then the skewness coefficients which combine the two.
\item Treat index numbers separately: build the four totals $\sum p_0q_0$, $\sum p_1q_0$, $\sum p_0q_1$, $\sum p_1q_1$ before applying the Laspeyres and Paasche formulae.
\item Finish by interpreting every number in the language of the situation (FCFA, students, prices) — a statistician who only computes has done half the job.
\end{enumerate}
\end{methode}

\courssub{Model answer}

\textbf{1. Extended frequency table.} The classes are given with \emph{class limits} $100$--$199$, $200$--$299$, \dots\ Since spending is recorded to the nearest franc, each class of width $100$ has \emph{class boundaries} obtained by shifting the limits by $0.5$: the class $100$--$199$ becomes $99.5$--$199.5$, and so on. The mid-value of each class is the mean of its boundaries, e.g.\ $x = \dfrac{99.5 + 199.5}{2} = 149.5$. Cumulative frequencies are running totals of $f$; the columns $fx$ and $fx^2$ prepare the mean and variance.

\begin{center}
\begin{tabular}{|c|c|c|c|c|c|c|}
\hline
\rowcolor{popblueL} Class & Boundaries & $x$ & $f$ & $F$ & $fx$ & $fx^2$ \\
\hline
100--199 & 99.5--199.5 & 149.5 & 3 & 3 & 448.5 & 67050.75 \\
\hline
200--299 & 199.5--299.5 & 249.5 & 6 & 9 & 1497.0 & 373501.50 \\
\hline
300--399 & 299.5--399.5 & 349.5 & 9 & 18 & 3145.5 & 1099352.25 \\
\hline
400--499 & 399.5--499.5 & 449.5 & 7 & 25 & 3146.5 & 1414351.75 \\
\hline
500--599 & 499.5--599.5 & 549.5 & 4 & 29 & 2198.0 & 1207801.00 \\
\hline
600--699 & 599.5--699.5 & 649.5 & 1 & 30 & 649.5 & 421850.25 \\
\hline
\rowcolor{popblueL} Total & & & 30 & & 11085 & 4583907.5 \\
\hline
\end{tabular}
\end{center}

Check: $\sum f = 30$ (all students counted), $\sum fx = 11085$, $\sum fx^2 = 4583907.5$.

\textbf{2. Histogram.} All class widths are equal to $100$, so the height of each rectangle equals the frequency. Rectangles are drawn over the boundary intervals $[99.5;199.5]$, \dots, $[599.5;699.5]$ with heights $3, 6, 9, 7, 4, 1$. The tallest bar corresponds to the modal class $299.5$--$399.5$.

\begin{popfigure}
\begin{tikzpicture}
\begin{axis}[ybar, bar width=0.55cm, width=11cm, height=6cm,
xlabel={Spending (FCFA)}, ylabel={Frequency},
xtick={149.5,249.5,349.5,449.5,549.5,649.5},
xticklabels={100--199,200--299,300--399,400--499,500--599,600--699},
x tick label style={rotate=30, anchor=east, font=\small},
ymin=0, ymax=10, axis lines=left, ymajorgrids]
\addplot[fill=popblue, draw=popdark] coordinates {(149.5,3) (249.5,6) (349.5,9) (449.5,7) (549.5,4) (649.5,1)};
\end{axis}
\end{tikzpicture}
\legende{Histogram of daily spending at the canteen}
\end{popfigure}

\textbf{3. Cumulative frequency table and ogive.} For a ``less than'' ogive, each cumulative frequency is plotted against the \emph{upper} class boundary:

\begin{center}
\begin{tabular}{|c|c|}
\hline
\rowcolor{popblueL} Upper boundary (FCFA) & Cumulative frequency $F$ \\
\hline
99.5 & 0 \\
\hline
199.5 & 3 \\
\hline
299.5 & 9 \\
\hline
399.5 & 18 \\
\hline
499.5 & 25 \\
\hline
599.5 & 29 \\
\hline
699.5 & 30 \\
\hline
\end{tabular}
\end{center}

(The point $(99.5, 0)$ anchors the curve: nobody spends less than $99.5$ FCFA.) Join the points with a smooth curve or straight segments. To read the quartiles, locate $N/4 = 7.5$, $N/2 = 15$ and $3N/4 = 22.5$ on the vertical axis, move across to the curve, then down to the spending axis.

\begin{popfigure}
\begin{tikzpicture}
\begin{axis}[width=11cm, height=7cm,
xlabel={Spending (FCFA)}, ylabel={Cumulative frequency},
xmin=50, xmax=750, ymin=0, ymax=32,
xtick={99.5,199.5,299.5,399.5,499.5,599.5,699.5},
x tick label style={rotate=30, anchor=east, font=\small},
axis lines=left, grid=major]
\addplot[color=popblue, thick, mark=*, mark size=1.5pt] coordinates {(99.5,0) (199.5,3) (299.5,9) (399.5,18) (499.5,25) (599.5,29) (699.5,30)};
\addplot[color=poporange, dashed] coordinates {(50,15) (366.2,15) (366.2,0)};
\addplot[color=popgreen, dashed] coordinates {(50,7.5) (274.5,7.5) (274.5,0)};
\addplot[color=poppurple, dashed] coordinates {(50,22.5) (463.8,22.5) (463.8,0)};
\end{axis}
\end{tikzpicture}
\legende{Ogive of daily spending, with graphical readings of $Q_1$, $Q_2$, $Q_3$}
\end{popfigure}

Reading across from $F = 7.5$, $15$ and $22.5$ gives approximately $Q_1 \approx 275$, $Q_2 \approx 366$ and $Q_3 \approx 464$ FCFA. Graphical readings are estimates: anything within about $\pm 5$ FCFA is acceptable.

\textbf{4. Mean, median, mode, range.} The mean of grouped data is
\[ \bar{x} = \frac{\sum fx}{\sum f} = \frac{11085}{30} = 369.5 \text{ FCFA}. \]
For the median, $N/2 = 15$, which falls in the class $299.5$--$399.5$ (the first class whose cumulative frequency, $18$, reaches $15$). With lower boundary $L = 299.5$, cumulative frequency before the class $F_b = 9$, class frequency $f_m = 9$ and class width $h = 100$:
\[ Q_2 = L + \frac{\tfrac{N}{2} - F_b}{f_m}\times h = 299.5 + \frac{15 - 9}{9}\times 100 = 299.5 + 66.67 \approx 366.2 \text{ FCFA}. \]
The modal class is $299.5$--$399.5$ (highest frequency, $9$). With $\Delta_1 = 9 - 6 = 3$ (excess over the previous class) and $\Delta_2 = 9 - 7 = 2$ (excess over the next class):
\[ M_o = L + \frac{\Delta_1}{\Delta_1 + \Delta_2}\times h = 299.5 + \frac{3}{3+2}\times 100 = 359.5 \text{ FCFA}. \]
Range $\approx 699.5 - 99.5 = 600$ FCFA (using the extreme class boundaries, since the exact raw values are unknown).

\textbf{5. Variance and standard deviation.} Using the computational formula for grouped data,
\[ \sigma^2 = \frac{\sum fx^2}{\sum f} - \bar{x}^2 = \frac{4583907.5}{30} - 369.5^2 = 152796.92 - 136530.25 \approx 16266.67 \text{ FCFA}^2, \]
\[ \sigma = \sqrt{16266.67} \approx 127.5 \text{ FCFA}. \]
So daily spending typically deviates from the mean by about $127$ FCFA.

\textbf{6. Quartiles.} $N/4 = 7.5$ falls in the class $199.5$--$299.5$ (cumulative frequency $9$ is the first to reach $7.5$), with $L = 199.5$, $F_b = 3$, $f = 6$:
\[ Q_1 = 199.5 + \frac{7.5 - 3}{6}\times 100 = 199.5 + 75 = 274.5 \text{ FCFA}. \]
$3N/4 = 22.5$ falls in the class $399.5$--$499.5$ (cumulative frequency $25$), with $L = 399.5$, $F_b = 18$, $f = 7$:
\[ Q_3 = 399.5 + \frac{22.5 - 18}{7}\times 100 = 399.5 + 64.29 \approx 463.8 \text{ FCFA}. \]
Interquartile range: $Q_3 - Q_1 = 463.8 - 274.5 = 189.3$ FCFA; semi-interquartile range: $\dfrac{189.3}{2} \approx 94.6$ FCFA. These calculated values agree very well with the ogive readings of question 3 ($275$, $366$, $464$), which confirms both the graph and the formulae.

\textbf{7. Skewness.} Pearson's second coefficient (used here because the mode is only an estimate for grouped data):
\[ S_k = \frac{3(\bar{x} - Q_2)}{\sigma} = \frac{3(369.5 - 366.2)}{127.5} = \frac{9.9}{127.5} \approx 0.08. \]
Quartile (Bowley's) coefficient:
\[ S_q = \frac{Q_3 + Q_1 - 2Q_2}{Q_3 - Q_1} = \frac{463.8 + 274.5 - 2(366.2)}{189.3} = \frac{5.9}{189.3} \approx 0.03. \]
Both coefficients are small and positive (each lies between $-1$ and $+1$, and values near $0$ indicate near-symmetry). The spending distribution is therefore \emph{almost symmetrical, with a very slight positive skew}: a few students spend much more than the others, which is consistent with the ordering $\bar{x} > Q_2 > M_o$ ($369.5 > 366.2 > 359.5$).

\textbf{8. Index numbers.} First compute the four totals needed by the formulae:

\begin{center}
\begin{tabular}{|l|c|c|c|c|}
\hline
\rowcolor{popblueL} Item & $p_0q_0$ & $p_1q_0$ & $p_0q_1$ & $p_1q_1$ \\
\hline
Rice & 20000 & 24000 & 17500 & 21000 \\
\hline
Beans & 9000 & 12000 & 8250 & 11000 \\
\hline
Puff-puff & 10000 & 15000 & 9000 & 13500 \\
\hline
Water & 15000 & 22500 & 14000 & 21000 \\
\hline
Bread & 9000 & 11250 & 8500 & 10625 \\
\hline
\rowcolor{popblueL} Total & 63000 & 84750 & 57250 & 77125 \\
\hline
\end{tabular}
\end{center}

Laspeyres index (base-year quantities as weights):
\[ L = \frac{\sum p_1 q_0}{\sum p_0 q_0}\times 100 = \frac{84750}{63000}\times 100 \approx 134.5. \]
Paasche index (current-year quantities as weights):
\[ P = \frac{\sum p_1 q_1}{\sum p_0 q_1}\times 100 = \frac{77125}{57250}\times 100 \approx 134.7. \]

\emph{Interpretation.} At last year's quantities, the same basket of items costs about $34.5\%$ more this year (Laspeyres); at this year's quantities, the increase is about $34.7\%$ (Paasche). The two indices are very close, so the conclusion is robust: \textbf{canteen prices have risen by roughly one third}, whatever quantities are used as weights.

\textbf{9. Model report (summary).} \emph{A typical student spends about 370 FCFA per day (mean $369.5$, median $366.2$ FCFA); half of the students spend between $274.5$ and $463.8$ FCFA. Spending is moderately spread out ($\sigma \approx 127.5$ FCFA, semi-interquartile range $\approx 94.6$ FCFA) and the distribution is nearly symmetric with a slight positive skew: a small group of big spenders pulls the mean above the median. Prices have increased by about $34.5\%$ on average (Laspeyres $134.5$, Paasche $134.7$), while quantities sold have fallen slightly for every item. Advice to the manager: the students' complaints are justified by the data; she should renegotiate with suppliers, keep at least one complete meal below $400$ FCFA so that the median student can still afford lunch, and monitor the prices of water and rice, which contribute most to the index.}

\textbf{10. Class discussion.} Groups compare their ogive readings, rounding choices and reports; small numerical differences (e.g.\ median $366$ vs $366.2$ FCFA) are acceptable provided the method — correct boundaries, correct interpolation formula, correct reading of the ogive — is sound.

\begin{attention}[Common errors to avoid]
Using class \emph{limits} instead of \cle{class boundaries} in the median and quartile formulae (this shifts every answer by $0.5$); forgetting to subtract $\bar{x}^2$ in the variance formula, or forgetting the square root at the end; swapping $q_0$ and $q_1$ between the Laspeyres and Paasche formulae; plotting the ogive against mid-values instead of upper class boundaries; concluding ``strong skew'' from a coefficient close to zero — near-zero means near-\emph{symmetric}.
\end{attention}

\begin{exercice}[Going further]
The following year, the same audit is repeated and the new prices (in FCFA) are: rice $330$, beans $210$, puff-puff $80$, water $150$, bread $130$, with quantities sold $75$, $60$, $190$, $145$, $88$ respectively. Taking last year (Data set B, index 0) as base, compute the new Laspeyres and Paasche price indices and comment on the trend.
\end{exercice}

\begin{corrige}[Exercise 1]
The new totals are $\sum p_2 q_0 = 330(80) + 210(60) + 80(200) + 150(150) + 130(90) = 26400 + 12600 + 16000 + 22500 + 11700 = 89200$ and $\sum p_2 q_1 = 330(75) + 210(60) + 80(190) + 150(145) + 130(88) = 24750 + 12600 + 15200 + 21750 + 11440 = 85740$, while $\sum p_0 q_1 = 250(75) + 150(60) + 50(190) + 100(145) + 100(88) = 18750 + 9000 + 9500 + 14500 + 8800 = 60550$. Hence
\[ L = \frac{89200}{63000}\times 100 \approx 141.6, \qquad P = \frac{85740}{60550}\times 100 \approx 141.6. \]
Prices have continued to rise, from an index of about $134.5$ to about $141.6$: the increase is slowing down but has not stopped.
\end{corrige}

\newpage

\coursec{Methods and worked examples}

\courssub{Methods and Worked Examples}

\begin{methode}[Constructing a Grouped Frequency Distribution]
\begin{enumerate}
    \item \textbf{Identify the range:} $R = x_{\max} - x_{\min}$.
    \item \textbf{Choose the number of classes $k$:} Typically $5 \le k \le 20$. Sturges' rule $k \approx 1 + 3.322\log_{10}n$ may be used as a guide.
    \item \textbf{Determine the class width $c$:} $c = \dfrac{R}{k}$. Round \textbf{up} to a convenient number (e.g., 5, 10, 20).
    \item \textbf{Set class limits:} Start at a value $\le x_{\min}$ (preferably a multiple of $c$). Define \cle{class boundaries} (real limits) to avoid gaps: for discrete data with unit gap, lower boundary $= \text{lower limit} - 0.5$, upper boundary $= \text{upper limit} + 0.5$. For continuous data, boundaries coincide with limits.
    \item \textbf{Tally and count frequencies $f$:} Use a tally column for accuracy.
    \item \textbf{Compute mid-points $x$:} $x = \dfrac{\text{Lower Limit} + \text{Upper Limit}}{2}$.
    \item \textbf{Check:} $\sum f = n$.
\end{enumerate}
\end{methode}

\begin{exoresolu}[Grouped Frequency Table -- Mathematics Mock Scores]
The scores (out of 100) of 50 Upper Sixth students in a Mathematics mock examination are:
\begin{center}
\begin{tabular}{cccccccccc}
42 & 67 & 55 & 78 & 34 & 89 & 51 & 63 & 72 & 48 \\
59 & 61 & 75 & 82 & 45 & 68 & 54 & 70 & 38 & 85 \\
57 & 64 & 79 & 49 & 66 & 52 & 73 & 80 & 41 & 69 \\
58 & 62 & 77 & 46 & 71 & 53 & 76 & 81 & 44 & 65 \\
56 & 60 & 74 & 47 & 70 & 50 & 69 & 83 & 43 & 66
\end{tabular}
\end{center}
Construct a grouped frequency distribution table with 7 classes. Include columns for Class Boundaries, Mid-points ($x$), Frequency ($f$), and Cumulative Frequency ($CF$).
\tcblower
\textbf{Solution:}
\begin{enumerate}
    \item \textbf{Range:} $x_{\max} = 89$, $x_{\min} = 34$, so $R = 89 - 34 = 55$.
    \item \textbf{Number of classes:} $k = 7$ (given).
    \item \textbf{Class width:} $c = \dfrac{55}{7} \approx 7.86$. Round \textbf{up} to $c = 8$.
    \item \textbf{Class limits:} Start at 34:
    34--41,\ 42--49,\ 50--57,\ 58--65,\ 66--73,\ 74--81,\ 82--89.
    \item \textbf{Class boundaries} (discrete data, gap $=1$): subtract 0.5 from each lower limit, add 0.5 to each upper limit.
    \item \textbf{Mid-points $x$:} average of the class limits, e.g. $(34+41)/2 = 37.5$.
    \item \textbf{Tally and frequency:} count the data points falling in each class.
\end{enumerate}

\begin{center}
\begin{tabular}{|c|c|c|c|c|c|}\hline
\rowcolor{popblueL}
\textbf{Class Limits} & \textbf{Class Boundaries} & \textbf{Mid-point ($x$)} & \textbf{Tally} & \textbf{Freq ($f$)} & \textbf{Cum. Freq ($CF$)} \\ \hline
34 -- 41 & 33.5 -- 41.5 & 37.5 & \texttt{||||} & 4 & 4 \\ \hline
42 -- 49 & 41.5 -- 49.5 & 45.5 & \texttt{|||| ||||} & 8 & 12 \\ \hline
50 -- 57 & 49.5 -- 57.5 & 53.5 & \texttt{|||| ||||} & 8 & 20 \\ \hline
58 -- 65 & 57.5 -- 65.5 & 61.5 & \texttt{|||| |||| |} & 11 & 31 \\ \hline
66 -- 73 & 65.5 -- 73.5 & 69.5 & \texttt{|||| ||||} & 9 & 40 \\ \hline
74 -- 81 & 73.5 -- 81.5 & 77.5 & \texttt{|||| ||} & 7 & 47 \\ \hline
82 -- 89 & 81.5 -- 89.5 & 85.5 & \texttt{|||} & 3 & 50 \\ \hline
\textbf{Total} & & & & \textbf{50} & \\ \hline
\end{tabular}
\end{center}
\textbf{Check:} $\sum f = 4+8+8+11+9+7+3 = 50 = n$. \checkmark
\begin{attention}[Common Mistake]
Using class limits (34--41) instead of boundaries (33.5--41.5) when drawing the Histogram or Ogive, or when applying interpolation formulas. \textbf{Always use boundaries for graphs and for interpolation.} The mid-point is calculated from the limits: $(34+41)/2 = 37.5$.
\end{attention}
\end{exoresolu}

\begin{methode}[Drawing Histogram, Frequency Polygon and Ogive]
\begin{enumerate}
    \item \textbf{Histogram:} Horizontal axis = \textbf{class boundaries}. Vertical axis = \textbf{frequency density} ($f/c$) if class widths are unequal; \textbf{frequency ($f$)} if all widths are equal. Bars touch (no gaps).
    \item \textbf{Frequency polygon:} Plot the points $(\text{mid-point } x,\ \text{frequency } f)$ and join them with straight lines. Close the polygon by adding an extra class with $f=0$ at each end (mid-points $x_1 - c$ and $x_k + c$) joined to the horizontal axis.
    \item \textbf{Ogive (cumulative frequency curve):} Plot the points $(\text{upper class boundary},\ CF)$. Join with a smooth curve (or straight segments). Start at $(\text{lower boundary of first class},\ 0)$.
    \item \textbf{Reading quantiles:} For the Median ($Q_2$), $Q_1$, $Q_3$, locate $n/2$, $n/4$, $3n/4$ on the vertical axis, trace horizontally to the curve, then vertically down to the horizontal axis.
\end{enumerate}
\end{methode}

\begin{exoresolu}[Graphical Representation and Estimation]
Using the grouped frequency table from the previous example (Mock Scores):
\begin{enumerate}
    \item Draw the Histogram and Frequency Polygon on the same axes.
    \item Draw the Ogive (Cumulative Frequency Curve).
    \item Use the Ogive to estimate the \textbf{Median}, \textbf{Lower Quartile ($Q_1$)}, \textbf{Upper Quartile ($Q_3$)}, and the \textbf{Interquartile Range (IQR)}.
    \item Estimate the number of students who scored \textbf{above 75}.
\end{enumerate}
\tcblower
\textbf{Solution:}

\textbf{1. Histogram and Frequency Polygon.}
Class widths are equal ($c=8$), so bar heights equal the frequencies.
Mid-points: 37.5, 45.5, 53.5, 61.5, 69.5, 77.5, 85.5.
Extra mid-points to close the polygon: $37.5-8=29.5$ and $85.5+8=93.5$ (both with $f=0$).

\begin{popfigure}
\begin{tikzpicture}[scale=0.65]
\begin{axis}[
    width=14cm, height=8cm,
    xlabel={Scores (marks)}, ylabel={Frequency},
    xmin=25, xmax=95, ymin=0, ymax=13,
    xtick={33.5,41.5,49.5,57.5,65.5,73.5,81.5,89.5},
    ytick={0,2,4,6,8,10,12},
    axis lines=middle,
    enlargelimits=false,
    tick label style={font=\footnotesize},
    label style={font=\small}
]
\draw[popblue, fill=popblueL, opacity=0.6] (33.5,0) rectangle (41.5,4);
\draw[popblue, fill=popblueL, opacity=0.6] (41.5,0) rectangle (49.5,8);
\draw[popblue, fill=popblueL, opacity=0.6] (49.5,0) rectangle (57.5,8);
\draw[popblue, fill=popblueL, opacity=0.6] (57.5,0) rectangle (65.5,11);
\draw[popblue, fill=popblueL, opacity=0.6] (65.5,0) rectangle (73.5,9);
\draw[popblue, fill=popblueL, opacity=0.6] (73.5,0) rectangle (81.5,7);
\draw[popblue, fill=popblueL, opacity=0.6] (81.5,0) rectangle (89.5,3);
\draw[poporange, thick, mark=*, mark size=2pt] plot coordinates {
    (29.5,0) (37.5,4) (45.5,8) (53.5,8) (61.5,11) (69.5,9) (77.5,7) (85.5,3) (93.5,0)
};
\end{axis}
\end{tikzpicture}
\legende{Histogram (blue bars) and frequency polygon (orange line) for the mock scores}
\end{popfigure}

\textbf{2. Ogive.}
Plot $(\text{upper boundary}, CF)$: $(41.5,4)$, $(49.5,12)$, $(57.5,20)$, $(65.5,31)$, $(73.5,40)$, $(81.5,47)$, $(89.5,50)$, starting from $(33.5, 0)$.

\begin{popfigure}
\begin{tikzpicture}[scale=0.65]
\begin{axis}[
    width=14cm, height=8cm,
    xlabel={Scores (upper class boundaries)}, ylabel={Cumulative frequency},
    xmin=25, xmax=95, ymin=0, ymax=55,
    xtick={33.5,41.5,49.5,57.5,65.5,73.5,81.5,89.5},
    ytick={0,10,20,30,40,50},
    axis lines=middle,
    enlargelimits=false,
    tick label style={font=\footnotesize},
    label style={font=\small},
    clip=false,
    no markers
]
\addplot[popgreen, thick, smooth, tension=0.5] coordinates {
    (33.5,0) (41.5,4) (49.5,12) (57.5,20) (65.5,31) (73.5,40) (81.5,47) (89.5,50)
};
\addplot[popgreen, only marks, mark=*, mark size=2pt] coordinates {
    (33.5,0) (41.5,4) (49.5,12) (57.5,20) (65.5,31) (73.5,40) (81.5,47) (89.5,50)
};
\draw[dashed, popmuted] (25,25) -- (62,25) -- (62,0) node[below, font=\tiny] {Median $\approx 62$};
\draw[dashed, popmuted] (25,12.5) -- (50,12.5) -- (50,0) node[below, font=\tiny] {$Q_1 \approx 50$};
\draw[dashed, popmuted] (25,37.5) -- (71,37.5) -- (71,0) node[below, font=\tiny] {$Q_3 \approx 71$};
\draw[dashed, poppink] (75.5,0) -- (75.5,42) -- (25,42) node[left, font=\tiny] {CF $\approx 42$};
\end{axis}
\end{tikzpicture}
\legende{Ogive with graphical estimation of the median, quartiles, and $CF$ at 75.5}
\end{popfigure}

\textbf{3. Estimates from the Ogive (graphical reading), $n = 50$:}
\begin{itemize}
    \item \textbf{Median ($Q_2$):} position $n/2 = 25$. From the graph: $\approx \mathbf{62}$.
    \item \textbf{Lower quartile ($Q_1$):} position $n/4 = 12.5$. From the graph: $\approx \mathbf{50}$.
    \item \textbf{Upper quartile ($Q_3$):} position $3n/4 = 37.5$. From the graph: $\approx \mathbf{71}$.
    \item \textbf{IQR} $= Q_3 - Q_1 \approx 71 - 50 = \mathbf{21}$ marks.
\end{itemize}

\textbf{4. Students scoring above 75:}
The score 75 lies in the class 74--81 (boundaries 73.5--81.5). We need $CF$ at the boundary 75.5.
Linear interpolation on the ogive segment from $(73.5, 40)$ to $(81.5, 47)$:
gradient $= \dfrac{47-40}{81.5-73.5} = \dfrac{7}{8} = 0.875$.
\[ CF(75.5) = 40 + 0.875 \times (75.5 - 73.5) = 40 + 1.75 = 41.75 \approx 42. \]
Students scoring at most 75: $\approx 42$. Hence students scoring \textbf{above 75} $= 50 - 42 = \mathbf{8}$.
\end{exoresolu}

\begin{methode}[Measures of Central Tendency: Mean, Median, Mode, Combined Mean]
\begin{enumerate}
    \item \textbf{Mean ($\bar{x}$):}
    \begin{itemize}
        \item Ungrouped: $\bar{x} = \dfrac{\sum x}{n}$.
        \item Grouped: $\bar{x} = \dfrac{\sum fx}{\sum f}$. For large numbers use \textbf{coding (assumed mean)}: $\bar{x} = A + \dfrac{\sum fd}{\sum f} \times c$, where $d = \dfrac{x-A}{c}$, $A$ = assumed mean (mid-point of a central class), $c$ = class width.
    \end{itemize}
    \item \textbf{Median (grouped):} identify the median class, the first class where $CF \ge n/2$.
    \[ \text{Median} = L_m + \left( \frac{\frac{n}{2} - CF_{m-1}}{f_m} \right) \times c \]
    $L_m$ = lower boundary of median class, $f_m$ = frequency of median class, $CF_{m-1}$ = cumulative frequency before the median class.
    \item \textbf{Mode (grouped):} identify the modal class (highest $f$).
    \[ \text{Mode} = L_{mo} + \left( \frac{\Delta_1}{\Delta_1 + \Delta_2} \right) \times c, \qquad \Delta_1 = f_{mo} - f_{mo-1},\quad \Delta_2 = f_{mo} - f_{mo+1}. \]
    $L_{mo}$ = lower boundary of the modal class.
    \item \textbf{Combined mean:} for $k$ groups with sizes $n_i$ and means $\bar{x}_i$:
    \[ \bar{x}_{\text{combined}} = \frac{\sum n_i \bar{x}_i}{\sum n_i} \]
\end{enumerate}
\end{methode}

\begin{exoresolu}[Central Tendency Calculations and Combined Mean]
Using the Mock Scores grouped data (table of the first example):
\begin{enumerate}
    \item Calculate the \textbf{Mean} using the coding method (assumed mean $A = 61.5$).
    \item Calculate the \textbf{Median} and \textbf{Mode} using the interpolation formulas.
    \item A second stream of 30 students sat the same examination and obtained a mean score of 58.2. Calculate the \textbf{Combined Mean} for all 80 students.
\end{enumerate}
\tcblower
\textbf{Data recall:} $n=50$, $c=8$. Median class: $n/2 = 25$; $CF$ first reaches $\ge 25$ at the class 58--65 (boundaries 57.5--65.5), so the median class is the 4th class. Modal class: highest $f = 11$, also the class 58--65.

\textbf{1. Mean (coding method).}
Let $A = 61.5$ and $d = \dfrac{x - 61.5}{8}$.

\begin{center}
\begin{tabular}{|c|c|c|c|c|}\hline
\rowcolor{popblueL}
Class & Mid-point $x$ & Freq $f$ & $d = \frac{x-61.5}{8}$ & $fd$ \\ \hline
34--41 & 37.5 & 4 & $-3$ & $-12$ \\ \hline
42--49 & 45.5 & 8 & $-2$ & $-16$ \\ \hline
50--57 & 53.5 & 8 & $-1$ & $-8$ \\ \hline
58--65 & 61.5 & 11 & 0 & 0 \\ \hline
66--73 & 69.5 & 9 & 1 & 9 \\ \hline
74--81 & 77.5 & 7 & 2 & 14 \\ \hline
82--89 & 85.5 & 3 & 3 & 9 \\ \hline
\textbf{Total} & & \textbf{50} & & \textbf{$-4$} \\ \hline
\end{tabular}
\end{center}

\[ \bar{x} = A + \frac{\sum fd}{\sum f} \times c = 61.5 + \frac{-4}{50} \times 8 = 61.5 - 0.64 = \mathbf{60.86} \]

\textbf{2. Median.}
$L_m = 57.5$, $f_m = 11$, $CF_{m-1} = 20$, $c = 8$, $n/2 = 25$.
\[ \text{Median} = 57.5 + \left( \frac{25 - 20}{11} \right) \times 8 = 57.5 + \frac{40}{11} \approx 57.5 + 3.64 = \mathbf{61.14} \]

\textbf{3. Mode.}
$L_{mo} = 57.5$, $f_{mo} = 11$, $f_{mo-1} = 8$, $f_{mo+1} = 9$, $c = 8$.
$\Delta_1 = 11 - 8 = 3$, \quad $\Delta_2 = 11 - 9 = 2$.
\[ \text{Mode} = 57.5 + \left( \frac{3}{3+2} \right) \times 8 = 57.5 + 4.8 = \mathbf{62.3} \]

\textbf{4. Combined mean.}
Group 1: $n_1 = 50$, $\bar{x}_1 = 60.86$. Group 2: $n_2 = 30$, $\bar{x}_2 = 58.2$.
\[ \bar{x}_{\text{comb}} = \frac{n_1\bar{x}_1 + n_2\bar{x}_2}{n_1+n_2} = \frac{50(60.86) + 30(58.2)}{80} = \frac{3043 + 1746}{80} = \frac{4789}{80} \approx \mathbf{59.86} \]
\end{exoresolu}

\begin{methode}[Measures of Dispersion: Variance and Standard Deviation]
\begin{enumerate}
    \item \textbf{Ungrouped data:}
    \[ \sigma^2 = \frac{\sum (x - \bar{x})^2}{n} = \frac{\sum x^2}{n} - \bar{x}^2 \quad ; \quad \sigma = \sqrt{\sigma^2} \]
    (For a sample, some conventions use $s^2 = \frac{\sum (x-\bar{x})^2}{n-1}$; at GCE Advanced Level the divisor $n$ is used unless stated otherwise.)
    \item \textbf{Grouped data:}
    \[ \sigma^2 = \frac{\sum f(x - \bar{x})^2}{\sum f} = \frac{\sum fx^2}{\sum f} - \bar{x}^2 \]
    \textbf{Coding method (recommended in examinations):} with $d = \dfrac{x-A}{c}$,
    \[ \sigma^2 = c^2 \left[ \frac{\sum fd^2}{\sum f} - \left( \frac{\sum fd}{\sum f} \right)^2 \right] \quad ; \quad \sigma = c \sqrt{ \frac{\sum fd^2}{\sum f} - \left( \frac{\sum fd}{\sum f} \right)^2 } \]
    \item \textbf{Coefficient of variation:} $CV = \dfrac{\sigma}{\bar{x}} \times 100\%$ (compares relative variability between data sets).
\end{enumerate}
\end{methode}

\begin{exoresolu}[Variance and Standard Deviation -- Grouped Data]
For the Mock Scores data (50 students), calculate:
\begin{enumerate}
    \item the variance ($\sigma^2$) and standard deviation ($\sigma$) using the coding method ($A=61.5$, $c=8$);
    \item the coefficient of variation ($CV$);
    \item the percentage of students within \textbf{one standard deviation} of the mean (i.e., between $\bar{x}-\sigma$ and $\bar{x}+\sigma$), assuming uniform distribution within classes.
\end{enumerate}
\tcblower
\textbf{1. Variance and standard deviation.}
Extend the coding table with $d^2$ and $fd^2$.

\begin{center}
\begin{tabular}{|c|c|c|c|c|c|}\hline
\rowcolor{popblueL}
Class & $x$ & $f$ & $d$ & $d^2$ & $fd^2$ \\ \hline
34--41 & 37.5 & 4 & $-3$ & 9 & 36 \\ \hline
42--49 & 45.5 & 8 & $-2$ & 4 & 32 \\ \hline
50--57 & 53.5 & 8 & $-1$ & 1 & 8 \\ \hline
58--65 & 61.5 & 11 & 0 & 0 & 0 \\ \hline
66--73 & 69.5 & 9 & 1 & 1 & 9 \\ \hline
74--81 & 77.5 & 7 & 2 & 4 & 28 \\ \hline
82--89 & 85.5 & 3 & 3 & 9 & 27 \\ \hline
\textbf{Total} & & \textbf{50} & \textbf{$-4$} & & \textbf{140} \\ \hline
\end{tabular}
\end{center}

\[ \frac{\sum fd^2}{\sum f} = \frac{140}{50} = 2.8, \qquad \left( \frac{\sum fd}{\sum f} \right)^2 = \left( \frac{-4}{50} \right)^2 = 0.0064 \]
\[ \sigma^2 = c^2 \left[ 2.8 - 0.0064 \right] = 64 \times 2.7936 = \mathbf{178.79} \]
\[ \sigma = \sqrt{178.79} \approx \mathbf{13.37} \text{ marks} \]

\textbf{2. Coefficient of variation.}
With $\bar{x} = 60.86$:
\[ CV = \frac{13.37}{60.86} \times 100\% \approx \mathbf{21.97\%} \]
\emph{Interpretation:} the spread of scores is about 22\% of the mean score — moderate variability.

\textbf{3. Percentage within $\bar{x} \pm \sigma$:}
$\bar{x} - \sigma = 60.86 - 13.37 = 47.49$ and $\bar{x} + \sigma = 60.86 + 13.37 = 74.23$.
We need the frequency between 47.49 and 74.23, assuming uniformity within each class:
\begin{itemize}
    \item \textbf{Class 42--49} (boundaries 41.5--49.5, $f=8$): included width $= 49.5 - 47.49 = 2.01$ out of 8. Included frequency $\approx 8 \times \frac{2.01}{8} = 2.01$.
    \item \textbf{Class 50--57:} fully inside. Frequency $= 8$.
    \item \textbf{Class 58--65:} fully inside. Frequency $= 11$.
    \item \textbf{Class 66--73:} fully inside. Frequency $= 9$.
    \item \textbf{Class 74--81} (boundaries 73.5--81.5, $f=7$): included width $= 74.23 - 73.5 = 0.73$ out of 8. Included frequency $\approx 7 \times \frac{0.73}{8} = 0.64$.
\end{itemize}
Total included frequency $\approx 2.01 + 8 + 11 + 9 + 0.64 = 30.65 \approx 31$ students.
Percentage $= \dfrac{31}{50} \times 100\% = \mathbf{62\%}$.
(For a Normal distribution one would expect about $68\%$; this data set is slightly skewed.)
\end{exoresolu}

\begin{methode}[Quantiles for Grouped Data: Quartiles, Deciles, Percentiles]
\begin{enumerate}
    \item \textbf{General formula for a quantile $Q$ at cumulative position $p$:} locate the first class where $CF \ge p$, then
    \[ Q = L_Q + \left( \frac{p - CF_{Q-1}}{f_Q} \right) \times c \]
    where $L_Q$ = lower class boundary of the quantile class, $CF_{Q-1}$ = cumulative frequency before that class, $f_Q$ = frequency of that class, $c$ = class width.
    \item \textbf{Positions:}
    \begin{itemize}
        \item Quartiles: $Q_1$ at $n/4$, $Q_2 \equiv$ Median at $n/2$, $Q_3$ at $3n/4$.
        \item Deciles: $D_k$ at $kn/10$, for $k = 1, \dots, 9$.
        \item Percentiles: $P_k$ at $kn/100$, for $k = 1, \dots, 99$.
    \end{itemize}
    \item \textbf{Interquartile range:} $IQR = Q_3 - Q_1$.
    \item \textbf{Semi-interquartile range (quartile deviation):} $SIQR = \dfrac{Q_3 - Q_1}{2}$.
\end{enumerate}
\end{methode}

\begin{exoresolu}[Quantiles, IQR and SIQR Calculation]
Using the Mock Scores data ($n=50$), calculate by formula:
\begin{enumerate}
    \item $Q_1$, $Q_3$, $D_5$, $P_{90}$;
    \item the Interquartile Range (IQR) and Semi-Interquartile Range (SIQR);
    \item the 70th percentile ($P_{70}$).
\end{enumerate}
\tcblower
\textbf{Reference table (boundaries, $f$, $CF$):}
\begin{center}
\begin{tabular}{|c|c|c|}\hline
\rowcolor{popblueL}
Class Boundaries & $f$ & $CF$ \\ \hline
33.5 -- 41.5 & 4 & 4 \\ \hline
41.5 -- 49.5 & 8 & 12 \\ \hline
49.5 -- 57.5 & 8 & 20 \\ \hline
57.5 -- 65.5 & 11 & 31 \\ \hline
65.5 -- 73.5 & 9 & 40 \\ \hline
73.5 -- 81.5 & 7 & 47 \\ \hline
81.5 -- 89.5 & 3 & 50 \\ \hline
\end{tabular}
\end{center}

\textbf{1. Quartiles and other quantiles ($c=8$):}
\begin{itemize}
    \item \textbf{$Q_1$:} position $= 50/4 = 12.5$. First $CF \ge 12.5$: class 41.5--49.5, with $L=41.5$, $CF_{prev}=4$, $f=8$.
    \[ Q_1 = 41.5 + \frac{12.5 - 4}{8} \times 8 = 41.5 + 8.5 = \mathbf{50.0} \]
    \item \textbf{$Q_3$:} position $= 3(50)/4 = 37.5$. First $CF \ge 37.5$: class 65.5--73.5, with $L=65.5$, $CF_{prev}=31$, $f=9$.
    \[ Q_3 = 65.5 + \frac{37.5 - 31}{9} \times 8 = 65.5 + \frac{52}{9} \approx 65.5 + 5.78 = \mathbf{71.28} \]
    \item \textbf{$D_5$:} position $= 5(50)/10 = 25$. Class 57.5--65.5, $L=57.5$, $CF_{prev}=20$, $f=11$.
    \[ D_5 = 57.5 + \frac{25-20}{11} \times 8 \approx 57.5 + 3.64 = \mathbf{61.14} \quad \text{(equals the Median, as expected)} \]
    \item \textbf{$P_{90}$:} position $= 90(50)/100 = 45$. First $CF \ge 45$: class 73.5--81.5, $L=73.5$, $CF_{prev}=40$, $f=7$.
    \[ P_{90} = 73.5 + \frac{45-40}{7} \times 8 = 73.5 + \frac{40}{7} \approx 73.5 + 5.71 = \mathbf{79.21} \]
\end{itemize}

\textbf{2. IQR and SIQR:}
\[ IQR = Q_3 - Q_1 = 71.28 - 50.00 = \mathbf{21.28}, \qquad SIQR = \frac{IQR}{2} = \frac{21.28}{2} = \mathbf{10.64} \]

\textbf{3. $P_{70}$:}
Position $= 70(50)/100 = 35$. First $CF \ge 35$: class 57.5--65.5, $L=57.5$, $CF_{prev}=20$, $f=11$.
\[ P_{70} = 57.5 + \frac{35-20}{11} \times 8 = 57.5 + \frac{120}{11} \approx 57.5 + 10.91 = \mathbf{68.41} \]
\end{exoresolu}

\begin{methode}[Index Numbers: Laspeyres and Paasche Price Indices]
\begin{enumerate}
    \item \textbf{Notation:} base year (0), current year ($n$); price $p$, quantity $q$. So $p_0, q_0$ are base-year values and $p_n, q_n$ current-year values.
    \item \textbf{Price relative} of an item: $I = \dfrac{p_n}{p_0} \times 100$.
    \item \textbf{Laspeyres price index (base-weighted):} uses \textbf{base-year quantities $q_0$} as weights.
    \[ L = \frac{\sum p_n q_0}{\sum p_0 q_0} \times 100 \]
    \emph{Advantage:} fixed weights, easy to compare over time. \emph{Disadvantage:} ignores changes in consumption patterns (substitution bias).
    \item \textbf{Paasche price index (current-weighted):} uses \textbf{current-year quantities $q_n$} as weights.
    \[ P = \frac{\sum p_n q_n}{\sum p_0 q_n} \times 100 \]
    \emph{Advantage:} reflects the current consumption pattern. \emph{Disadvantage:} weights change every period, so pure price change is harder to isolate.
    \item \textbf{Typical relationship:} when consumers substitute away from goods whose prices rise fastest, $L > P$.
\end{enumerate}
\end{methode}

\begin{exoresolu}[Laspeyres and Paasche Indices -- Household Commodities]
A household consumes three commodities: rice, beans and oil. Data for the base year (2020) and the current year (2024):

\begin{center}
\begin{tabular}{|l|c|c|c|c|}\hline
\rowcolor{popblueL}
\textbf{Commodity} & \textbf{Price 2020 ($p_0$)} & \textbf{Qty 2020 ($q_0$)} & \textbf{Price 2024 ($p_n$)} & \textbf{Qty 2024 ($q_n$)} \\ \hline
Rice (kg) & 400 FCFA & 50 & 600 FCFA & 40 \\ \hline
Beans (kg) & 800 FCFA & 20 & 1000 FCFA & 25 \\ \hline
Oil (litre) & 1200 FCFA & 10 & 1500 FCFA & 8 \\ \hline
\end{tabular}
\end{center}

Calculate:
\begin{enumerate}
    \item the \textbf{Laspeyres price index} for 2024 (base 2020 $= 100$);
    \item the \textbf{Paasche price index} for 2024 (base 2020 $= 100$);
    \item interpret the results and explain why they differ.
\end{enumerate}
\tcblower
\textbf{Working table} (values in FCFA):

\begin{center}
\begin{tabular}{|l|c|c|c|c|}\hline
\rowcolor{popblueL}
Item & $p_n q_0$ & $p_0 q_0$ & $p_n q_n$ & $p_0 q_n$ \\ \hline
Rice & $600\times50=30\,000$ & $400\times50=20\,000$ & $600\times40=24\,000$ & $400\times40=16\,000$ \\ \hline
Beans & $1000\times20=20\,000$ & $800\times20=16\,000$ & $1000\times25=25\,000$ & $800\times25=20\,000$ \\ \hline
Oil & $1500\times10=15\,000$ & $1200\times10=12\,000$ & $1500\times8=12\,000$ & $1200\times8=9\,600$ \\ \hline
\textbf{Sum} & \textbf{65\,000} & \textbf{48\,000} & \textbf{61\,000} & \textbf{45\,600} \\ \hline
\end{tabular}
\end{center}

\textbf{1. Laspeyres index ($L$):}
\[ L = \frac{\sum p_n q_0}{\sum p_0 q_0} \times 100 = \frac{65\,000}{48\,000} \times 100 \approx \mathbf{135.4} \]

\textbf{2. Paasche index ($P$):}
\[ P = \frac{\sum p_n q_n}{\sum p_0 q_n} \times 100 = \frac{61\,000}{45\,600} \times 100 \approx \mathbf{133.8} \]

\textbf{3. Interpretation:}
\begin{itemize}
    \item Both indices exceed 100, indicating a general price increase (inflation) of about 34--35\% between 2020 and 2024.
    \item Here Laspeyres ($135.4$) $>$ Paasche ($133.8$).
    \item \textbf{Reason:} the household reduced its consumption of rice ($50 \to 40$) and oil ($10 \to 8$), whose prices rose sharply (by 50\% and 25\%), and increased its consumption of beans ($20 \to 25$), whose price rose less (25\%).
    \item Laspeyres keeps the \emph{old} quantities (heavy weight on the now-expensive rice and oil), so it overstates the increase in the cost of living.
    \item Paasche uses the \emph{new} quantities (less weight on the expensive items), reflecting the substitution effect, and therefore gives a lower index.
\end{itemize}
\end{exoresolu}

\begin{methode}[Skewness: Pearson's and Quartile Coefficients]
\begin{enumerate}
    \item \textbf{Symmetry:} Mean $=$ Median $=$ Mode $\implies$ symmetrical distribution (zero skewness).
    \item \textbf{Positive skew (right skew):} long tail on the right; typically Mean $>$ Median $>$ Mode.
    \item \textbf{Negative skew (left skew):} long tail on the left; typically Mean $<$ Median $<$ Mode.
    \item \textbf{Pearson's coefficient of skewness:}
    \begin{itemize}
        \item $Sk_P = \dfrac{\text{Mean} - \text{Mode}}{\sigma}$ (when the mode is well defined);
        \item $Sk_P = \dfrac{3(\text{Mean} - \text{Median})}{\sigma}$ (more common; uses the empirical relation Mean $-$ Mode $\approx 3(\text{Mean} - \text{Median})$).
    \end{itemize}
    Values usually lie between $-3$ and $+3$. $Sk_P > 0$: positive skew; $Sk_P < 0$: negative skew.
    \item \textbf{Quartile (Bowley's) coefficient of skewness:}
    \[ Sk_Q = \frac{Q_3 + Q_1 - 2Q_2}{Q_3 - Q_1} = \frac{Q_3 + Q_1 - 2\,\text{Median}}{IQR} \]
    Always between $-1$ and $+1$; robust, since quartiles are not affected by extreme values.
    \begin{itemize}
        \item $Sk_Q = 0$: symmetrical ($Q_3 - Q_2 = Q_2 - Q_1$);
        \item $Sk_Q > 0$: positive skew ($Q_3 - Q_2 > Q_2 - Q_1$);
        \item $Sk_Q < 0$: negative skew.
    \end{itemize}
\end{enumerate}
\end{methode}

\begin{exoresolu}[Skewness Analysis of the Mock Scores]
Using the statistics calculated for the Mock Scores ($n=50$):
$\bar{x} = 60.86$, Median $= 61.14$, Mode $= 62.30$, $\sigma = 13.37$, $Q_1 = 50.00$, $Q_2 = 61.14$, $Q_3 = 71.28$, $IQR = 21.28$.
\begin{enumerate}
    \item Determine the direction of skewness by comparing the Mean, Median and Mode.
    \item Calculate \textbf{Pearson's coefficient of skewness} (median-based formula).
    \item Calculate the \textbf{Quartile coefficient of skewness}.
    \item Comment on the consistency of the two measures.
\end{enumerate}
\tcblower
\textbf{1. Direction of skewness:}
Mean ($60.86$) $<$ Median ($61.14$) $<$ Mode ($62.30$).
\cle{Conclusion:} the distribution is \textbf{negatively skewed (left-skewed)}: the tail stretches towards the lower scores, while the bulk of the distribution lies at the higher scores.

\textbf{2. Pearson's coefficient ($Sk_P$):}
\[ Sk_P = \frac{3(\bar{x} - \text{Median})}{\sigma} = \frac{3(60.86 - 61.14)}{13.37} = \frac{-0.84}{13.37} \approx \mathbf{-0.063} \]
(The mode-based version gives $\frac{60.86 - 62.30}{13.37} \approx -0.108$, also negative; the median-based formula is preferred as it is more stable.)

\textbf{3. Quartile coefficient ($Sk_Q$):}
\[ Sk_Q = \frac{Q_3 + Q_1 - 2Q_2}{Q_3 - Q_1} = \frac{71.28 + 50.00 - 2(61.14)}{21.28} = \frac{-1.00}{21.28} \approx \mathbf{-0.047} \]

\textbf{4. Comment:}
\begin{itemize}
    \item Both coefficients are \textbf{negative}, confirming \textbf{negative skewness}.
    \item Both are \textbf{close to zero} ($|Sk_P| \approx 0.06$, $|Sk_Q| \approx 0.05$), so the distribution is \textbf{nearly symmetrical}.
    \item The slight negative skew reflects a small number of weaker students pulling the mean just below the median.
    \item Pearson's coefficient uses the mean and standard deviation (sensitive to extreme values); the quartile coefficient is robust. Their agreement strengthens the conclusion.
\end{itemize}
\begin{aretenir}[Exam Tip]
When asked to ``comment on the skewness'' of a distribution, always:
\begin{enumerate}
    \item state the order of Mean, Median and Mode;
    \item state the direction (positive / negative / approximately symmetrical);
    \item compute a coefficient (Pearson's or the quartile coefficient);
    \item interpret its magnitude (``highly skewed'', ``moderately skewed'', ``approximately symmetric'').
\end{enumerate}
\end{aretenir}
\end{exoresolu}

\newpage

\coursec{Exercises}

\courssub{I check my knowledge}

\begin{exercice}[MCQ: Data Types and Collection]
Which of the following variables is \emph{continuous}?
\begin{enumerate}
\item The number of students in a classroom.
\item The height of a student measured in centimetres.
\item The number of cars passing a checkpoint in an hour.
\item The number of defective items in a batch of 50.
\end{enumerate}
\end{exercice}

\begin{exercice}[True/False with Justification]
State whether each statement is true or false. Justify your answer briefly.
\begin{enumerate}
\item The mode of a data set is always unique.
\item The standard deviation of a set of values is always less than or equal to the range.
\item In a grouped frequency distribution, the median class is the class with the highest frequency.
\item The Laspeyres price index uses current year quantities as weights.
\end{enumerate}
\end{exercice}

\begin{exercice}[Definitions and Terminology]
Define the following terms precisely:
\begin{enumerate}
\item \cle{Class boundaries} in a grouped frequency distribution.
\item \cle{Frequency density} and why it is used when drawing a histogram with unequal class widths.
\item \cle{Semi-interquartile range} (quartile deviation).
\item \cle{Positively skewed distribution} in terms of the relative positions of mean, median and mode.
\end{enumerate}
\end{exercice}

\begin{exercice}[Direct Calculations]
The marks (out of 20) of 10 students in a test are:
\[
12,\ 15,\ 14,\ 12,\ 18,\ 15,\ 13,\ 16,\ 15,\ 14.
\]
Calculate:
\begin{enumerate}
\item The mean.
\item The median.
\item The mode.
\item The variance and standard deviation (use the formula for ungrouped data).
\end{enumerate}
\end{exercice}

\courssub{I practise}

\begin{exercice}[Grouped Data: Mean, Median, Mode by Interpolation]
The table below shows the distribution of masses (in kg) of 80 packets of rice sold in a market in Douala.

\begin{center}
\begin{tabular}{|c|c|c|c|c|c|c|}\hline
Mass (kg) & 1.0--1.4 & 1.5--1.9 & 2.0--2.4 & 2.5--2.9 & 3.0--3.4 & 3.5--3.9 \\ \hline
Frequency & 8 & 15 & 22 & 18 & 12 & 5 \\ \hline
\end{tabular}
\end{center}

\begin{enumerate}
\item Construct the cumulative frequency table (using true class boundaries).
\item Estimate the mean mass using the coding method with an assumed mean of 2.2 kg.
\item Estimate the median mass by linear interpolation.
\item Estimate the modal mass by linear interpolation (use the histogram method formula).
\end{enumerate}
\end{exercice}

\begin{exercice}[Histogram with Unequal Class Widths and Frequency Polygon]
The following data gives the ages of 100 patients attending a health centre in Yaoundé.

\begin{center}
\begin{tabular}{|c|c|c|c|c|c|}\hline
Age (years) & 0--9 & 10--19 & 20--29 & 30--49 & 50--69 \\ \hline
Number of patients & 15 & 20 & 25 & 30 & 10 \\ \hline
\end{tabular}
\end{center}

\begin{enumerate}
\item Calculate the frequency density for each class.
\item Draw a histogram to represent the data (use graph paper; show frequency density on the vertical axis).
\item Superimpose the frequency polygon on the same axes.
\item Use the histogram to estimate the mode graphically.
\end{enumerate}
\end{exercice}

\begin{exercice}[Cumulative Frequency Curve (Ogive) and Percentiles]
The table shows the distribution of monthly incomes (in thousands of CFA francs) of 120 workers.

\begin{center}
\begin{tabular}{|c|c|c|c|c|c|c|}\hline
Income (1000 FCFA) & 50--59 & 60--69 & 70--79 & 80--89 & 90--99 & 100--109 \\ \hline
Frequency & 10 & 18 & 32 & 35 & 15 & 10 \\ \hline
\end{tabular}
\end{center}

\begin{enumerate}
\item Construct the cumulative frequency table (using upper class boundaries).
\item Draw the ogive (cumulative frequency curve).
\item From the ogive, estimate:
\begin{enumerate}
\item The median income.
\item The first quartile ($Q_1$) and third quartile ($Q_3$).
\item The 90th percentile ($P_{90}$).
\item The semi-interquartile range.
\end{enumerate}
\end{enumerate}
\end{exercice}

\begin{exercice}[Combined Mean and Comparison of Consistency]
Two classes, A and B, sat the same Mathematics test.
Class A has 35 students with a mean mark of 12.4 and a standard deviation of 2.1.
Class B has 45 students with a mean mark of 10.8 and a standard deviation of 3.0.

\begin{enumerate}
\item Calculate the combined mean mark for the two classes together.
\item Which class shows more consistent performance? Justify your answer using the coefficient of variation.
\item If the pass mark is 10, and assuming the marks are roughly normally distributed, which class would you expect to have a higher pass rate? Explain briefly.
\end{enumerate}
\end{exercice}

\courssub{I prepare for the GCE}

\begin{exercice}[GCE Style: Full Analysis of a Grouped Distribution] \textbf{[20 marks]}
The table below shows the distribution of marks obtained by 80 candidates in a GCE Advanced Level Pure Mathematics with Statistics paper.

\begin{center}
\begin{tabular}{|c|c|c|c|c|c|c|c|}\hline
Marks & 10--19 & 20--29 & 30--39 & 40--49 & 50--59 & 60--69 & 70--79 \\ \hline
Frequency & 5 & 10 & 15 & 20 & 18 & 8 & 4 \\ \hline
\end{tabular}
\end{center}

\begin{enumerate}
\item Construct the cumulative frequency table using true class boundaries. \textbf{[2 marks]}
\item Calculate the mean mark using the coding method (choose a suitable assumed mean). \textbf{[4 marks]}
\item Calculate the standard deviation of the marks. \textbf{[4 marks]}
\item Estimate the median mark by linear interpolation. \textbf{[3 marks]}
\item Estimate the modal mark by linear interpolation. \textbf{[3 marks]}
\item Draw a histogram to represent the data (use frequency density if necessary) and superimpose the frequency polygon. \textbf{[2 marks]}
\item Comment on the performance of the candidates based on the values obtained. \textbf{[2 marks]}
\end{enumerate}
\end{exercice}

\begin{exercice}[GCE Style: Index Numbers and Skewness] \textbf{[18 marks]}
A household buys four food items: Rice, Beans, Fish, and Oil. The prices (in FCFA per kg) and quantities consumed (in kg per month) for the base year 2023 and the current year 2024 are given below.

\begin{center}
\begin{tabular}{|c|c|c|c|c|}\hline
Item & Price 2023 ($p_0$) & Quantity 2023 ($q_0$) & Price 2024 ($p_1$) & Quantity 2024 ($q_1$) \\ \hline
Rice & 600 & 50 & 720 & 55 \\
Beans & 800 & 20 & 900 & 18 \\
Fish & 2500 & 5 & 3000 & 4 \\
Oil & 1200 & 10 & 1300 & 12 \\ \hline
\end{tabular}
\end{center}

\begin{enumerate}
\item Compute the simple aggregate price index for 2024 with 2023 as base year. \textbf{[2 marks]}
\item Compute the Laspeyres price index for 2024. \textbf{[3 marks]}
\item Compute the Paasche price index for 2024. \textbf{[3 marks]}
\item Interpret the difference between the Laspeyres and Paasche indices in this context. \textbf{[2 marks]}
\end{enumerate}

A distribution of weekly wages has the following summary statistics:
Mean = 52,000 FCFA, Median = 50,000 FCFA, Mode = 46,000 FCFA, Standard deviation = 8,000 FCFA.

\begin{enumerate}
\setcounter{enumi}{4}
\item Calculate Pearson's first coefficient of skewness (using mode). \textbf{[2 marks]}
\item Calculate Pearson's second coefficient of skewness (using median). \textbf{[2 marks]}
\item State the direction of skewness and justify your answer. \textbf{[2 marks]}
\item Given that $Q_1 = 44,000$ FCFA and $Q_3 = 58,000$ FCFA, compute the quartile coefficient of skewness and the semi-interquartile range. \textbf{[2 marks]}
\end{enumerate}
\end{exercice}

\begin{exercice}[GCE Style: Mixed Problem – Reconstruction and Questionnaire Critique] \textbf{[17 marks]}

\textbf{Part A: Reconstruction of a Frequency Distribution} \textbf{[10 marks]}
A frequency distribution of the heights (in cm) of 60 plants was partially destroyed. The remaining data is:

\begin{center}
\begin{tabular}{|c|c|c|c|c|c|}\hline
Height (cm) & 10--14 & 15--19 & 20--24 & 25--29 & 30--34 \\ \hline
Frequency & 5 & $f_1$ & 18 & $f_2$ & 7 \\ \hline
\end{tabular}
\end{center}

It is known that the total frequency is 60 and the mean height is 22.5 cm.

\begin{enumerate}
\item Find the missing frequencies $f_1$ and $f_2$. \textbf{[4 marks]}
\item Using the completed table, calculate the variance and standard deviation of the heights. \textbf{[4 marks]}
\item Estimate the median height by linear interpolation. \textbf{[2 marks]}
\end{enumerate}

\textbf{Part B: Questionnaire Critique} \textbf{[7 marks]}
A student designs the following questionnaire to investigate the favourite subjects of Upper Sixth students in a school. Identify \textbf{three} sources of bias or poor design in the questionnaire and rewrite the faulty questions to eliminate the bias.

\begin{enumerate}
\item \emph{``Don't you agree that Mathematics is the most important subject?''} \hfill (Yes / No)
\item \emph{``Which subject do you like: Mathematics, Physics, or Chemistry?''} \hfill (Tick one)
\item \emph{``How many hours do you spend studying each day?''} \hfill (0--1, 1--2, 2--3, 3--4, 4+)
\item \emph{``What is your favourite subject? \_\_\_\_\_\_\_\_\_\_''} \hfill (Open ended)
\end{enumerate}
\end{exercice}



\newpage

\coursec{Solutions to the exercises}

\courssub{I check my knowledge}

\begin{corrige}[Exercise 1 — MCQ: Data Types and Collection]
Answer: \textbf{(b)}. Height can take any value within an interval (e.g.\ $172.4$ cm), so it is \cle{continuous}. The other three variables can only take whole-number values (they result from \emph{counting}), so they are \cle{discrete}.
\end{corrige}

\begin{corrige}[Exercise 2 — True/False with Justification]
\begin{enumerate}
\item \textbf{False.} A data set may have no mode (all values distinct), one mode, or several modes (bimodal, multimodal). Example: $2, 2, 5, 5, 7$ has two modes.
\item \textbf{True.} The standard deviation measures spread \emph{about the mean} while the range measures the total spread from minimum to maximum. In fact it can be proved that $\sigma \leq \dfrac{\text{range}}{2}$ for every data set, so certainly $\sigma \leq$ range.
\item \textbf{False.} The median class is the class containing the $\dfrac{N}{2}$th value (located from the cumulative frequency). The class with the highest frequency is the \emph{modal} class.
\item \textbf{False.} The Laspeyres index uses \emph{base year} quantities $q_0$ as weights: $L = \dfrac{\sum p_1 q_0}{\sum p_0 q_0}\times 100$. It is the \emph{Paasche} index that uses current year quantities $q_1$.
\end{enumerate}
\end{corrige}

\begin{corrige}[Exercise 3 — Definitions and Terminology]
\begin{enumerate}
\item The \cle{class boundaries} (true class limits) are the values that separate two consecutive classes without any gap. For a class $10$--$19$ followed by $20$--$29$ (data recorded to the nearest unit), the class boundaries are $9.5$--$19.5$ and $19.5$--$29.5$.
\item The \cle{frequency density} of a class is $\text{frequency density} = \dfrac{\text{frequency}}{\text{class width}}$. In a histogram, the \emph{area} of each bar (not its height) must be proportional to the frequency; when class widths are unequal, plotting frequency density on the vertical axis guarantees this.
\item The \cle{semi-interquartile range} is $Q = \dfrac{Q_3 - Q_1}{2}$, where $Q_1$ and $Q_3$ are the first and third quartiles. It measures the spread of the middle $50\%$ of the data.
\item A distribution is \cle{positively skewed} (skewed to the right) when it has a long tail towards the high values; then $\text{mode} < \text{median} < \text{mean}$.
\end{enumerate}
\end{corrige}

\begin{corrige}[Exercise 4 — Direct Calculations]
\begin{enumerate}
\item $\sum x = 144$, so $\bar{x} = \dfrac{144}{10} = 14.4$.
\item Ordered data: $12, 12, 13, 14, 14, 15, 15, 15, 16, 18$. Median $= \dfrac{14+15}{2} = 14.5$.
\item Mode $= 15$ (it occurs 3 times).
\item $\sum x^2 = 2104$, so
\[
\sigma^2 = \frac{\sum x^2}{n} - \bar{x}^2 = \frac{2104}{10} - 14.4^2 = 210.4 - 207.36 = 3.04,
\qquad \sigma = \sqrt{3.04} \approx 1.74.
\]
\end{enumerate}
\end{corrige}

\courssub{I practise}

\begin{corrige}[Exercise 5 — Grouped Data: Mean, Median, Mode by Interpolation]
True class boundaries: $0.95$--$1.45$, $1.45$--$1.95$, \dots, $3.45$--$3.95$; midpoints: $1.2, 1.7, 2.2, 2.7, 3.2, 3.7$; class width $c = 0.5$.
\begin{enumerate}
\item Cumulative frequencies (against upper class boundaries): $8,\ 23,\ 45,\ 63,\ 75,\ 80$.
\item With $A = 2.2$ and $u = \dfrac{x - A}{0.5}$: the values of $u$ are $-2, -1, 0, 1, 2, 3$ and
\[
\sum fu = 8(-2) + 15(-1) + 22(0) + 18(1) + 12(2) + 5(3) = 26.
\]
\[
\bar{x} = A + c\,\frac{\sum fu}{N} = 2.2 + 0.5\times\frac{26}{80} = 2.2 + 0.1625 \approx 2.36\ \text{kg}.
\]
\item $\dfrac{N}{2} = 40$ falls in the class $1.95$--$2.45$ (cumulative frequency 23 before it, frequency 22):
\[
\text{Median} = 1.95 + \frac{40 - 23}{22}\times 0.5 = 1.95 + 0.386 \approx 2.34\ \text{kg}.
\]
\item Modal class: $1.95$--$2.45$ (frequency 22). With $\Delta_1 = 22 - 15 = 7$ and $\Delta_2 = 22 - 18 = 4$:
\[
\text{Mode} = 1.95 + \frac{\Delta_1}{\Delta_1 + \Delta_2}\times 0.5 = 1.95 + \frac{7}{11}\times 0.5 \approx 2.27\ \text{kg}.
\]
\end{enumerate}
\end{corrige}

\begin{corrige}[Exercise 6 — Histogram with Unequal Class Widths and Frequency Polygon]
\begin{enumerate}
\item Class boundaries: $-0.5$--$9.5$, $9.5$--$19.5$, $19.5$--$29.5$, $29.5$--$49.5$, $49.5$--$69.5$; widths $10, 10, 10, 20, 20$.
\begin{center}
\begin{tabular}{|c|c|c|c|}\hline
\rowcolor{popblueL} Class & Width & Frequency & Frequency density \\ \hline
0--9 & 10 & 15 & 1.5 \\ \hline
10--19 & 10 & 20 & 2.0 \\ \hline
20--29 & 10 & 25 & 2.5 \\ \hline
30--49 & 20 & 30 & 1.5 \\ \hline
50--69 & 20 & 10 & 0.5 \\ \hline
\end{tabular}
\end{center}
\item The histogram has bars over $[-0.5, 9.5]$, $[9.5, 19.5]$, $[19.5, 29.5]$, $[29.5, 49.5]$, $[49.5, 69.5]$ with heights $1.5, 2.0, 2.5, 1.5, 0.5$ respectively. Note how the wide classes ($30$--$49$ and $50$--$69$) get \emph{shorter} bars than their raw frequencies would suggest: this is exactly why frequency density is needed.
\item The frequency polygon joins the points $(4.5, 1.5)$, $(14.5, 2.0)$, $(24.5, 2.5)$, $(39.5, 1.5)$, $(59.5, 0.5)$ (midpoints of the class boundaries against frequency density), extended to the axis at both ends.
\item Drawing the diagonal lines in the tallest bar (the standard histogram construction) gives a modal age of approximately $24$ years.
\end{enumerate}
\end{corrige}

\begin{corrige}[Exercise 7 — Cumulative Frequency Curve (Ogive) and Percentiles]
\begin{enumerate}
\item Upper class boundaries and cumulative frequencies:
\begin{center}
\begin{tabular}{|c|c|c|c|c|c|c|}\hline
\rowcolor{popblueL} Upper boundary & 59.5 & 69.5 & 79.5 & 89.5 & 99.5 & 109.5 \\ \hline
Cumulative frequency & 10 & 28 & 60 & 95 & 110 & 120 \\ \hline
\end{tabular}
\end{center}
\item Plot the points $(59.5, 10), (69.5, 28), (79.5, 60), (89.5, 95), (99.5, 110), (109.5, 120)$ and join them with a smooth S-shaped curve.
\item Reading from the ogive (values confirmed by linear interpolation):
\begin{enumerate}
\item Median at $\dfrac{N}{2} = 60$: median $\approx 79.5$ thousand FCFA.
\item $Q_1$ at $\dfrac{N}{4} = 30$: $Q_1 = 69.5 + \dfrac{30-28}{32}\times 10 \approx 70.1$; $Q_3$ at $\dfrac{3N}{4} = 90$: $Q_3 = 79.5 + \dfrac{90-60}{35}\times 10 \approx 88.1$ (thousand FCFA).
\item $P_{90}$ at $0.9 \times 120 = 108$: since the cumulative frequency 108 lies between 95 (at 89.5) and 110 (at 99.5),
\[
P_{90} = 89.5 + \frac{108 - 95}{15}\times 10 = 89.5 + 8.67 \approx 98.2\ \text{thousand FCFA}.
\]
\item Semi-interquartile range $= \dfrac{Q_3 - Q_1}{2} \approx \dfrac{88.1 - 70.1}{2} = 9.0$ thousand FCFA.
\end{enumerate}
\end{enumerate}
\end{corrige}

\begin{corrige}[Exercise 8 — Combined Mean and Comparison of Consistency]
\begin{enumerate}
\item Combined mean:
\[
\bar{x} = \frac{n_A \bar{x}_A + n_B \bar{x}_B}{n_A + n_B} = \frac{35(12.4) + 45(10.8)}{80} = \frac{434 + 486}{80} = \frac{920}{80} = 11.5.
\]
\item Coefficient of variation: $CV = \dfrac{\sigma}{\bar{x}}\times 100\%$.
\[
CV_A = \frac{2.1}{12.4}\times 100\% \approx 16.9\%, \qquad CV_B = \frac{3.0}{10.8}\times 100\% \approx 27.8\%.
\]
Since $CV_A < CV_B$, \textbf{Class A} is more consistent (its marks are less dispersed relative to its mean).
\item Class A: mean 12.4 with small spread 2.1, so the pass mark 10 lies about $1.1$ standard deviations below the mean: the great majority pass. Class B: mean 10.8 with spread 3.0, so the pass mark lies only about $0.27$ standard deviations below the mean: a substantial fraction of the candidates fall below 10. We therefore expect a \textbf{higher pass rate in Class A}.
\end{enumerate}
\end{corrige}

\courssub{I prepare for the GCE}

\begin{corrige}[Exercise 9 — GCE Style: Full Analysis of a Grouped Distribution]
Midpoints: $14.5, 24.5, 34.5, 44.5, 54.5, 64.5, 74.5$; class width $c = 10$; boundaries $9.5$--$19.5$, \dots, $69.5$--$79.5$.
\begin{enumerate}
\item Cumulative frequencies: $5,\ 15,\ 30,\ 50,\ 68,\ 76,\ 80$ against upper boundaries $19.5, 29.5, \dots, 79.5$. \hfill \textbf{[2 marks: 1 for boundaries, 1 for cumulative frequencies]}
\item Take $A = 44.5$, $u = \dfrac{x - 44.5}{10}$, so $u = -3, -2, -1, 0, 1, 2, 3$:
\[
\sum fu = 5(-3) + 10(-2) + 15(-1) + 20(0) + 18(1) + 8(2) + 4(3) = -4.
\]
\[
\bar{x} = 44.5 + 10\times\frac{-4}{80} = 44.5 - 0.5 = 44.0.
\]
\textbf{[4 marks: 1 for coding, 1 for $\sum fu$, 1 for formula, 1 for answer]}
\item $\sum fu^2 = 5(9) + 10(4) + 15(1) + 20(0) + 18(1) + 8(4) + 4(9) = 186$.
\[
\sigma^2 = c^2\left[\frac{\sum fu^2}{N} - \left(\frac{\sum fu}{N}\right)^2\right] = 100\left[\frac{186}{80} - \left(\frac{-4}{80}\right)^2\right] = 100(2.325 - 0.0025) = 232.25,
\]
\[
\sigma = \sqrt{232.25} \approx 15.2.
\]
\textbf{[4 marks: 1 for $\sum fu^2$, 1 for formula, 1 for variance, 1 for $\sigma$]}
\item $\dfrac{N}{2} = 40$ lies in the class $39.5$--$49.5$ (cumulative frequency 30 before it, frequency 20):
\[
\text{Median} = 39.5 + \frac{40 - 30}{20}\times 10 = 39.5 + 5 = 44.5.
\]
\textbf{[3 marks: 1 for locating class, 1 for formula, 1 for answer]}
\item Modal class $39.5$--$49.5$; $\Delta_1 = 20 - 15 = 5$, $\Delta_2 = 20 - 18 = 2$:
\[
\text{Mode} = 39.5 + \frac{5}{5+2}\times 10 = 39.5 + 7.14 \approx 46.6.
\]
\textbf{[3 marks: 1 for modal class, 1 for formula, 1 for answer]}
\item All classes have equal width 10, so the histogram bars have heights equal to the frequencies over the intervals $[9.5, 19.5], \dots, [69.5, 79.5]$. The frequency polygon joins $(14.5, 5), (24.5, 10), (34.5, 15), (44.5, 20), (54.5, 18), (64.5, 8), (74.5, 4)$. \hfill \textbf{[2 marks: 1 for histogram, 1 for polygon]}
\item The mean mark (44.0) is below 50, and the large standard deviation (15.2) shows marks are widely spread. Since $\text{mean} < \text{median} < \text{mode}$, the distribution is slightly negatively skewed: a tail of weak candidates pulls the mean down, while the bulk of candidates scored in the 40--59 region. Overall performance is mediocre, with a small group of very weak candidates. \hfill \textbf{[2 marks: 1 for level of performance, 1 for shape/spread comment]}
\end{enumerate}
\emph{Examiner's expectations:} correct use of class boundaries, a clear coding table, interpolation formulas quoted before substitution, and a comment that uses the computed statistics (not a vague remark).
\end{corrige}

\begin{corrige}[Exercise 10 — GCE Style: Index Numbers and Skewness]
\begin{enumerate}
\item Simple aggregate price index:
\[
I = \frac{\sum p_1}{\sum p_0}\times 100 = \frac{720 + 900 + 3000 + 1300}{600 + 800 + 2500 + 1200}\times 100 = \frac{5920}{5100}\times 100 \approx 116.1.
\]
\textbf{[2 marks: 1 for formula and sums, 1 for answer]}
\item Laspeyres (base year quantities):
\[
\sum p_1 q_0 = 720(50) + 900(20) + 3000(5) + 1300(10) = 36000 + 18000 + 15000 + 13000 = 82000,
\]
\[
\sum p_0 q_0 = 600(50) + 800(20) + 2500(5) + 1200(10) = 30000 + 16000 + 12500 + 12000 = 70500,
\]
\[
L = \frac{82000}{70500}\times 100 \approx 116.3.
\]
\textbf{[3 marks: 1 for each aggregate sum, 1 for index]}
\item Paasche (current year quantities):
\[
\sum p_1 q_1 = 720(55) + 900(18) + 3000(4) + 1300(12) = 39600 + 16200 + 12000 + 15600 = 83400,
\]
\[
\sum p_0 q_1 = 600(55) + 800(18) + 2500(4) + 1200(12) = 33000 + 14400 + 10000 + 14400 = 71800,
\]
\[
P = \frac{83400}{71800}\times 100 \approx 116.2.
\]
\textbf{[3 marks: 1 for each aggregate sum, 1 for index]}
\item Both indices show a price rise of about $16\%$. The Laspeyres index (116.3) is slightly higher than the Paasche index (116.2) because the household slightly reduced its consumption of the items whose prices rose most (fish, beans) and increased consumption of the item whose price rose least (oil). Laspeyres keeps the old basket and therefore marginally overstates the rise in the cost of living, while Paasche reflects the new basket. \hfill \textbf{[2 marks]}
\item Pearson's first coefficient:
\[
Sk = \frac{\text{mean} - \text{mode}}{\sigma} = \frac{52000 - 46000}{8000} = \frac{6000}{8000} = 0.75.
\]
\textbf{[2 marks: 1 for formula, 1 for answer]}
\item Pearson's second coefficient:
\[
Sk = \frac{3(\text{mean} - \text{median})}{\sigma} = \frac{3(52000 - 50000)}{8000} = \frac{6000}{8000} = 0.75.
\]
\textbf{[2 marks: 1 for formula, 1 for answer]}
\item Both coefficients are positive, and $\text{mode} < \text{median} < \text{mean}$: the distribution is \textbf{positively skewed} (a long tail of high wages pulls the mean above the median and mode). \hfill \textbf{[2 marks: 1 for direction, 1 for justification]}
\item Quartile coefficient of skewness:
\[
Sk_Q = \frac{Q_3 + Q_1 - 2Q_2}{Q_3 - Q_1} = \frac{58000 + 44000 - 2(50000)}{58000 - 44000} = \frac{2000}{14000} \approx 0.14,
\]
confirming positive skewness. Semi-interquartile range:
\[
Q = \frac{Q_3 - Q_1}{2} = \frac{58000 - 44000}{2} = 7000\ \text{FCFA}.
\]
\textbf{[2 marks: 1 for each quantity]}
\end{enumerate}
\emph{Examiner's expectations:} index formulas stated with correct weights, aggregate sums shown explicitly, and skewness conclusions supported both by the sign of the coefficients and by the ordering of mean, median and mode.
\end{corrige}

\begin{corrige}[Exercise 11 — GCE Style: Mixed Problem – Reconstruction and Questionnaire Critique]
\textbf{Part A.} Midpoints: $12, 17, 22, 27, 32$.
\begin{enumerate}
\item Total frequency: $5 + f_1 + 18 + f_2 + 7 = 60 \Rightarrow f_1 + f_2 = 30$.
Mean condition: $\sum fx = 60(22.5) = 1350$:
\[
5(12) + 17f_1 + 18(22) + 27f_2 + 7(32) = 1350 \Rightarrow 680 + 17f_1 + 27f_2 = 1350 \Rightarrow 17f_1 + 27f_2 = 670.
\]
Substituting $f_2 = 30 - f_1$: $17f_1 + 27(30 - f_1) = 670 \Rightarrow 810 - 10f_1 = 670 \Rightarrow f_1 = 14$, $f_2 = 16$. \hfill \textbf{[4 marks: 1 per equation, 1 per value]}
\item $\sum fx^2 = 5(144) + 14(289) + 18(484) + 16(729) + 7(1024) = 720 + 4046 + 8712 + 11664 + 7168 = 32310$.
\[
\sigma^2 = \frac{\sum fx^2}{N} - \bar{x}^2 = \frac{32310}{60} - 22.5^2 = 538.5 - 506.25 = 32.25, \qquad \sigma = \sqrt{32.25} \approx 5.68\ \text{cm}.
\]
\textbf{[4 marks: 1 for $\sum fx^2$, 1 for formula, 1 for variance, 1 for $\sigma$]}
\item Cumulative frequencies: $5, 19, 37, 53, 60$. $\dfrac{N}{2} = 30$ lies in the class $19.5$--$24.5$:
\[
\text{Median} = 19.5 + \frac{30 - 19}{18}\times 5 = 19.5 + 3.06 \approx 22.6\ \text{cm}.
\]
\textbf{[2 marks: 1 for locating class, 1 for answer]}
\end{enumerate}
\textbf{Part B.} Any three of the following \textbf{[7 marks: up to 2 marks per fault identified and corrected, 1 mark for general presentation]}:
\begin{enumerate}
\item \textbf{Leading question:} it pushes the respondent to agree. Rewrite: \emph{``In your opinion, how important is Mathematics compared with your other subjects?''} (Very important / Important / Not important).
\item \textbf{Restricted choice list:} only three science subjects are offered, so students who prefer other subjects cannot answer honestly. Rewrite: \emph{``Which of the following subjects do you like most?''} with a fuller list (Mathematics, Physics, Chemistry, Biology, Economics, Literature, \dots) plus an \emph{``Other: \_\_\_\_''} option.
\item \textbf{Overlapping response categories:} a student studying exactly 1, 2, 3 or 4 hours fits two boxes. Rewrite with non-overlapping intervals: 0--0.9, 1--1.9, 2--2.9, 3--3.9, 4 or more (or ``less than 1, at least 1 but less than 2, \dots'').
\item The open-ended question (4) is acceptable in principle, though answers may be hard to classify; offering a tick-box list with an ``Other'' option makes processing easier.
\end{enumerate}
\emph{Examiner's expectations:} in Part A, the two equations must be set up from the total frequency and the mean before solving; in Part B, each fault must be \emph{named} (leading question, restricted options, overlapping classes) and a corrected version of the question must actually be written out.
\end{corrige}

\newpage

\coursec{Summary sheet}

\begin{popfigure}
\begin{tikzpicture}[mindmap, grow cyclic, every node/.style={concept, text=white, font=\small, align=center},
root concept/.append style={concept color=popdark, font=\small\bfseries},
level 1/.append style={level distance=3.8cm, sibling angle=51.43},
level 1 concept/.append style={font=\scriptsize}]
\node[root concept]{Descriptive Statistics}
child[concept color=popblue]{ node[align=center]{Data \& collection\\ qualitative, quantitative, primary, secondary} }
child[concept color=popgreen]{ node[align=center]{Tables \& grouping\\ tally, limits, boundaries, class mark} }
child[concept color=poporange]{ node[align=center]{Graphs\\ bar, pie, line, histogram, polygon, ogive} }
child[concept color=poppurple]{ node[align=center]{Central tendency\\ mode, median, mean, combined mean} }
child[concept color=poppink]{ node[align=center]{Dispersion\\ range, variance, s.d., quartiles, IQR} }
child[concept color=popgold]{ node[align=center]{Index numbers\\ Laspeyres, Paasche, weighted} }
child[concept color=popmuted]{ node[align=center]{Skewness\\ Pearson \& quartile coefficients} };
\end{tikzpicture}
\legende{Mind map of the chapter: from raw data to summary measures, index numbers and the shape of a distribution.}
\end{popfigure}

\begin{bilan}
\textbf{Key definitions.}
\begin{itemize}
\item \cle{Statistics}: the science of the collection, organisation, presentation, analysis and interpretation of numerical data. A \cle{population} is the whole group under study; a \cle{sample} is a representative part of it; each member is a \cle{statistical unit}; a measurable characteristic is a \cle{variable}.
\item \cle{Qualitative data} (attributes that cannot be measured numerically: colour, sex, tribe) vs \cle{quantitative data} (numerical: \cle{discrete} -- obtained by counting, e.g.\ number of children; or \cle{continuous} -- obtained by measuring, e.g.\ height, mass).
\item \cle{Primary data} (collected first-hand by interview, observation, questionnaire or experiment) vs \cle{secondary data} (taken from already existing records, e.g.\ school registers, MINEDUB reports).
\item \cle{Raw data}: data as collected, before any organisation. A \cle{frequency} $f$ is the number of times a value occurs; $\sum f = N$ is the total frequency.
\item For a grouped distribution: \cle{class limits} (the stated end values), \cle{class boundaries} (the true end values, e.g.\ $9.5$--$19.5$ for the class $10$--$19$), \cle{class width} $c$ (difference between successive lower boundaries), \cle{class mark} (midpoint) $x_i=\dfrac{\text{lower limit}+\text{upper limit}}{2}$.
\item \cle{Cumulative frequency}: the running total of frequencies; plotted against \emph{upper} class boundaries it gives the \cle{ogive} (cumulative frequency curve).
\item \cle{Index number}: a ratio comparing the value of a variable in a given period to its value in a \cle{base period}, usually expressed as a percentage.
\item \cle{Skewness}: the degree of asymmetry (lack of symmetry) of a distribution.
\end{itemize}

\textbf{Formulas to know.}
\begin{itemize}
\item Mean (ungrouped): $\bar{x}=\dfrac{\sum x}{n}$; grouped: $\bar{x}=\dfrac{\sum f_i x_i}{\sum f_i}$; with assumed mean $A$ and coding $u_i=\dfrac{x_i-A}{c}$: $\bar{x}=A+c\,\dfrac{\sum f_i u_i}{\sum f_i}$.
\item \cle{Combined mean} of two groups: $\bar{x}=\dfrac{n_1\bar{x}_1+n_2\bar{x}_2}{n_1+n_2}$ (extends naturally to $k$ groups: $\bar{x}=\dfrac{\sum n_i\bar{x}_i}{\sum n_i}$).
\item Median (grouped): $M_e=L+\dfrac{\frac{N}{2}-F}{f}\times c$, where $L$ = lower boundary of the median class, $F$ = cumulative frequency before it, $f$ = its frequency, $c$ = its width.
\item Mode (grouped): $M_o=L+\dfrac{\Delta_1}{\Delta_1+\Delta_2}\times c$, with $\Delta_1=f_m-f_{\text{before}}$, $\Delta_2=f_m-f_{\text{after}}$, where $f_m$ is the largest frequency (modal class).
\item \cle{Range} $=$ largest value $-$ smallest value.
\item Variance (ungrouped): $\sigma^2=\dfrac{\sum(x-\bar{x})^2}{n}=\dfrac{\sum x^2}{n}-\bar{x}^2$; standard deviation $\sigma=\sqrt{\sigma^2}$.
\item Variance (grouped): $\sigma^2=\dfrac{\sum f_i x_i^2}{\sum f_i}-\bar{x}^2$, or by coding $\sigma^2=c^2\left(\dfrac{\sum f_i u_i^2}{\sum f_i}-\left(\dfrac{\sum f_i u_i}{\sum f_i}\right)^2\right)$.
\item Combined variance of two groups: $\sigma^2=\dfrac{n_1(\sigma_1^2+d_1^2)+n_2(\sigma_2^2+d_2^2)}{n_1+n_2}$ with $d_i=\bar{x}_i-\bar{x}$.
\item Quartiles by interpolation: $Q_k=L+\dfrac{\frac{kN}{4}-F}{f}\times c$ ($k=1,2,3$; note $Q_2=M_e$); similarly deciles $D_k$ using $\dfrac{kN}{10}$ and percentiles $P_k$ using $\dfrac{kN}{100}$. \cle{Interquartile range} $=Q_3-Q_1$; \cle{semi-interquartile range} $=\dfrac{Q_3-Q_1}{2}$.
\item Simple price index: $\dfrac{p_n}{p_0}\times 100$. \cle{Laspeyres index}: $L=\dfrac{\sum p_n q_0}{\sum p_0 q_0}\times 100$ (base-period quantities as weights). \cle{Paasche index}: $P=\dfrac{\sum p_n q_n}{\sum p_0 q_n}\times 100$ (current-period quantities as weights). Weighted index number: $\dfrac{\sum wI}{\sum w}$, where $I$ are the individual indices and $w$ their weights.
\item \cle{Pearson's coefficient of skewness}: $S_k=\dfrac{\bar{x}-M_o}{\sigma}$; when the mode is awkward to obtain, the empirical approximation $S_k\approx\dfrac{3(\bar{x}-M_e)}{\sigma}$ is used. \cle{Quartile coefficient of skewness}: $S_Q=\dfrac{Q_3+Q_1-2Q_2}{Q_3-Q_1}$. $S>0$: positively skewed (tail to the right); $S<0$: negatively skewed (tail to the left); $S=0$: symmetrical.
\end{itemize}

\textbf{Methods.}
\begin{itemize}
\item Grouping data: find the range, choose the number of classes (usually between 5 and 15), fix a convenient class width $c$, write the classes, \cle{tally} the raw data, then count the frequencies and check that $\sum f=N$.
\item Histogram: rectangles drawn on the \emph{class boundaries} with no gaps, area proportional to frequency (for equal class widths, height $=$ frequency). The \cle{frequency polygon} joins the midpoints of the tops of the bars (or the points $(x_i,f_i)$).
\item Ogive: plot cumulative frequency against \emph{upper} class boundaries, join the points, and read the median at $\dfrac{N}{2}$, $Q_1$ at $\dfrac{N}{4}$ and $Q_3$ at $\dfrac{3N}{4}$ on the cumulative axis.
\item Pie chart angle for each category: $\dfrac{\text{frequency}}{\text{total frequency}}\times 360^{\circ}$. Bar chart: separated bars of equal width, height $=$ frequency. Line graph: values plotted against time (or another ordered variable) and joined.
\end{itemize}

\textbf{Common mistakes.}
\begin{itemize}
\item Using class \emph{limits} instead of class \emph{boundaries} for histograms, ogives and interpolation formulas.
\item Forgetting to multiply by the class width $c$ (for the mean) or by $c^2$ (for the variance) when decoding a coded calculation.
\item Confusing $\sum f x^2$ with $\left(\sum f x\right)^2$ in the variance formula.
\item Mixing up Laspeyres (base-period quantities $q_0$) and Paasche (current-period quantities $q_n$).
\item Reading the median or quartiles from the frequency polygon instead of the ogive.
\item Taking $\dfrac{N}{2}$ as the median itself instead of the \emph{position} of the median.
\item Giving a negative variance or a standard deviation larger than the range without checking the arithmetic.
\end{itemize}

\textbf{GCE tips.}
\begin{itemize}
\item Show the full frequency table with columns $x_i$, $f_ix_i$, $f_ix_i^2$ (or the coding columns $u_i$, $f_iu_i$, $f_iu_i^2$): method marks are awarded for correct table work.
\item Quote the interpolation formula before substituting, and give the final answer with units and sensible rounding (usually 1 or 2 decimal places, or 3 significant figures).
\item For skewness questions, state the sign \emph{and} interpret it in words: for a positively skewed distribution, mean $>$ median $>$ mode and the tail stretches to the right; the reverse holds for negative skewness.
\item Quick checks: a variance can never be negative; an index number near 100 means little change from the base period; an index of 125 means a $25\%$ increase.
\item In pie chart questions, check that your angles add up to $360^{\circ}$ (allowing for rounding).
\end{itemize}
\end{bilan}

\findecours

\end{document}
